%----------------------------------------------------------------------------------------------------%
%----------------------------------------------------------------------------------------------------%
%%%%%%%%%%%                                              HEADER                                                %%%%%%%%%%%
%----------------------------------------------------------------------------------------------------%
%----------------------------------------------------------------------------------------------------%
%
% Universale header.tex – gilt für Arbeitsblätter, Leistungskontrollen, Handreichungen, Nachschlagewerke.
% Neue Pakete und TikZ-Bibliotheken werden HIER ergänzt (nicht im Dokument) und unten protokolliert.
% Quelle: https://github.com/mq90/latex-skills-schule  (CC BY 4.0, Martin Quandt)
%
% Änderungsprotokoll
% 2026-11-09  zusammengeführt aus allen Handreichungs-Headern: TikZ-Bibliotheken fit, shapes.multipart,
%             arrows.meta, shapes.callouts, circuits.logic.US, circuits.logic.IEC; Pakete array, booktabs,
%             longtable; xskak, chessboard, sudoku, tikzpeople, tipa (safe); MusiXTeX als auskommentierte
%             Option (definiert \cup als Achtelnote um – bricht $A \cup B$)
% 2026-11-11  Naturwissenschaften: circuitikz [european] (Schaltpläne), mhchem [version=4] (\ce),
%             chemfig (Strukturformeln), pgf-PeriodicTable (\pgfPT), genealogytree [all] (Stammbäume).
%             NICHT siunitx: kollidiert mit units (\unit[10]{cm}); siunitx nur in eigenständigen Dokumenten
% 2026-11-11  TikZ-Bibliothek decorations.pathmorphing (Federn, Wellen: decoration=coil/snake)
% 2026-11-23  twemojis: farbige Emojis mit pdflatex, z.B. \twemoji[height=1cm]{man standing} (Skill latex-sachskizzen)
%






%----------------------------------------------------------------------------------------------------%
%%%%%%%%%%%%%                            Gestaltung des Dokuments                                %%%%%%%%%%
%----------------------------------------------------------------------------------------------------%



%Definition der Dokumentenklasse
\documentclass[12pt]{article} %book, article, report, standalone



%Pakete für eine saubere Kodierung
\usepackage[utf8]{inputenc} %UTF-8 als Zeichenkodierung  ... ABER nicht immer notwendig, je nach Compiler
\usepackage[T1]{fontenc} %Verwendet westeuropäische Schriftkodierung
\usepackage[ngerman]{babel} %Ermöglicht deutsche Übersetzung und Silbentrennung
\usepackage{lmodern} %Verwendet moderne Schriftarten



%Pakete für die Gestaltung des Dokuments
\usepackage{geometry} %Ermöglicht die Ausrichtung des Blattes mit \geometry{...}
\geometry{a4paper, top=15mm, left=15mm, right=20mm, bottom=15mm, headsep=10mm, footskip=10mm} %gewählte Ausrichtung des Blattes
\pagestyle{empty} %Art der Seitennummerierung, z.B.: empty , plain
\usepackage{fancyhdr} %Ermöglicht die Gestaltung einer Kopfzeile für LK's oder KA's 
\usepackage{lscape} %Ermöglicht eventuelles Querformat des Blattes mit \begin{landscape} ... \end{landscape}
\usepackage{pdfpages} %Ermöglicht einbinden PDF-Dokumente bzw. einzelene Seiten mit \includepdf[pages={1,3,5}]{Dateiname} bzw. pages=-
\usepackage[german]{layout} %Ermöglicht es mit den Befehl \layout sich die Default-Layout beschreiben zu lassen (Bild vom Layout)
%----------------------------------------------------------------------------------------------------%








%----------------------------------------------------------------------------------------------------%
%%%%%%%%%%%%                                        Packages                                                  %%%%%%%%%%%
%----------------------------------------------------------------------------------------------------%



%Pakete für die Strukturierung des Dokuments
\usepackage[shortlabels]{enumitem} %Ermöglicht die Nummerierung selbst festzulegen mit \begin{enumerate}[Art der Nummerierung] ...   \end{enumerate}
\parindent=0px %Setzt die Absatzeinrückung auf 0 Pixel und ernfernt damit unnötige Leerzeichen vor dem ersten Wort eines neuen Abschnitts
\usepackage{setspace} %Ermöglicht verschiedene Funktion zur Abstandsregelung, z.B.: \hspace{2.5cm} ; \onehalfspacing
\onehalfspacing %Setzt den Zeilenabstand auf 1,5
\usepackage[style=german]{csquotes}%Ermöglicht den deutschen Zitierstil darstellen zu lassen mt \enquote{Zitat}
\usepackage[breaklinks=true,%Zeilenumbruch der Verlinkung ist bündig
			colorlinks=true,%Links können farbig sein
			linkcolor=blue,%Linkfarbe wird auf blau gesetzt
			urlcolor=darkgray,%Farb der URL wird auf dunkelgrau gesetzt
			plainpages=false,%
			hyperfootnotes=false,%
			pdfpagelabels=false%
		]{hyperref} %Ermöglicht die Verlinkung innerhalb des Dokuments 
\usepackage{multicol} %Ermöglicht gleichgroße Aufteilung des Dokuments in n Spalten mit \begin{multicols}{n} ... \end{multicols}
\usepackage{xcolor,colortbl} %Ermöglicht Zeilen, Spalten und Zellen einer Tabelle zu färben
\usepackage{caption} %Ermöglicht die Bezeichnung und Strukturierung von Abbildungen oder Quellcode
\usepackage{fancybox} % Ermöglicht die Umrandung einzelner Abschnitte
\usepackage[most]{tcolorbox} %Ermöglicht die Nutzung von Colorboxen für Merksätze, Lösungen, Definitionen, usw mit \newtcolorbox
\tcbuselibrary{breakable} %Aufteilen der Colorboxen, welche über mehrere Seiten gehen
\usepackage{ifthen} % Ermöglicht \ifthenelse Befehle
\allowdisplaybreaks %Erlaubt Seitenumbrüche innerhalb von align-Umgebung



%Nützliche Pakete für Symbole und mathematischen Formelsatz
\usepackage{amsmath, amssymb, mathrsfs, amsfonts, amsthm, dsfont} %sämtliche Pakete für mathematische Zeichen, Symbole und Formelsatz
\usepackage{units} %Ermöglicht die Angabe von Einheiten mit \unit[Wert]{Einheit}
\usepackage{gensymb} %Ermöglicht Nutzung von \degree für das macro des Grad-Symbols °
\usepackage[official]{eurosym} %Ermöglicht die Nutzung von \euro{}
\usepackage{pifont} %Ermöglicht die Nutzung verschiedener Fonts, z.B.: \ding{42} , \ding{52} , \ding{56}
\usepackage{tikzsymbols} %Ermöglicht die Nutzung von u.a. Smileys und anderen Symbolen, z.B.: \Winkey , \Smiley[1.5]



%Pakete zur Darstellung von Grafiken
\usepackage{tikz} %\begin{tikzpicture} .... \end{tikzpicture}
\usepackage{pgf, pgfplots} %Ermöglicht 2D und 3D Plof-Funktion
\pgfplotsset{compat=1.17}
\usetikzlibrary{arrows,%Pfeile
			calc,%Rechnen in Koordinaten, z.B. mit ($(1,2)+5*(0,1)$) auf Punkt (1,7) kommen
			angles, plotmarks, quotes, babel, intersections, shapes.geometric, decorations.pathreplacing, decorations.pathmorphing, automata, positioning, shadows, patterns, patterns.meta,%
			fit,%Rahmen um mehrere Knoten (UML, Automaten)
			shapes.multipart,%geteilte Knoten (UML-Klassen)
			arrows.meta,%moderne Pfeilspitzen, z.B. -{Triangle[open]}
			shapes.callouts,%Sprechblasen
			circuits.logic.US, circuits.logic.IEC%Logikgatter (digitale Schaltungen)
			} %nützliche Bibliotheken für TikZ
\usepackage{animate} %animierte Darstellungen
\usepackage{graphicx} %\includegraphics[width=..px,height=..px]{Dateiname}
\usepackage{multirow} %Zellen über mehrere Zeilen mit \multirow{n}{*}{Text}
\usepackage{array} %erweiterte Spaltentypen, z.B. >{\centering\arraybackslash}p{2cm}
\usepackage{booktabs} %schöne Tabellenlinien \toprule, \midrule, \bottomrule
\usepackage{longtable} %Tabellen über mehrere Seiten



%gelegentlich hilfreiche Zusatzpakete
\usepackage{hologo} %\hologo{LaTeX}
\usepackage{suetterl} %Sütterlin als Schriftart mit \mathfrak{Text}
\usepackage{qrcode} %Einbindung eines QR-Codes 
\usepackage{url} %Einbindung eines URL-Links
\usepackage{ccicons} %Creative Commons Symbole
\usepackage{listings} %\begin{lstlisting} ... \end{lstlisting}
\usepackage{struktex} %Struktogrammdarstellung
\usepackage{ulem} %Unterstreichung mit \uline{} ggf. mit \hfill im Inneren für Unterstrich bis zum Ende der Zeile
\usepackage{xskak} %Schach: Partien und Züge mit \newchessgame, \mainline{1. e4 e5} (lädt skak)
\usepackage{chessboard} %Schachbretter mit \chessboard[setfen=...]
\usepackage{sudoku} %Sudoku-Gitter mit \begin{sudoku} |5|3| | ... |. \end{sudoku}
\usepackage{tikzpeople} %Personen-Figuren als TikZ-Knoten, z.B. \node[alice, minimum size=2cm]{};
\usepackage{twemojis} %farbige Emojis (Twemoji, CC BY 4.0) als Vektorgrafik: \twemoji[height=1cm]{man standing}, im Text \texttwemoji{light bulb}
\usepackage[safe]{tipa} %Lautschrift mit \textipa{...}; [safe], weil tipa sonst \; \: \! \* \s umdefiniert (bricht z.B. \(\;\)) – Langformen siehe latex-lautschrift
\usepackage[european]{circuitikz} %Schaltpläne mit deutschen Symbolen: \begin{circuitikz} \draw (0,0) to[lamp] (2,0); \end{circuitikz}
\usepackage[version=4]{mhchem} %Reaktionsgleichungen mit \ce{2 H2 + O2 -> 2 H2O}
\usepackage{chemfig} %Strukturformeln mit \chemfig{H-C(-[2]H)(-[6]H)-H}, freie Elektronenpaare mit \charge
\usepackage{pgf-PeriodicTable} %Periodensystem mit \pgfPT[...], deutsche Namen mit \pgfPTsetLanguage{de}
\usepackage[all]{genealogytree} %Stammbäume mit \begin{genealogypicture} ... \end{genealogypicture}
%\usepackage{musixtex} %Notensatz – NUR für Musik-Dokumente einkommentieren: definiert u.a. \cup als Note um (bricht $A \cup B$)
%\usepackage[list=true, font=large, labelfont=bf, labelformat=brace, position=top]{subcaption}



%abgeänderte Befehle
\renewcommand{\arraystretch}{1.25} %
%\renewcommand{\familydefault}{\sfdefault} %serifenlose Schriftart
\renewcommand{\lstlistingname}{Quellcode} %Umbenennung
\renewcommand{\lstlistlistingname}{Quellcodeverzeichnis} %Umbenennung
%----------------------------------------------------------------------------------------------------%
%----------------------------------------------------------------------------------------------------%
%----------------------------------------------------------------------------------------------------%
%%%% %%%%%%                                     NEWCOMMANDS                                         %%%%%%%%%%%%
%----------------------------------------------------------------------------------------------------%
%----------------------------------------------------------------------------------------------------%
%
% Universale newcommands.tex – alle eigenen Befehle, Farben, Boxen, Kopfzeilen, TikZ- und Listing-Stile.
% Neue Befehle werden HIER ergänzt (nicht im Dokument) und unten protokolliert.
% Quelle: https://github.com/mq90/latex-skills-schule  (CC BY 4.0, Martin Quandt)
%
% Änderungsprotokoll
% 2026-11-09  zusammengeführt aus allen Handreichungen: \TikZ, Box »ausgabe«, Stile für
%             Programmablaufpläne (input, operation, condition, output) und UML (umlclass, inherit,
%             compose, aggregate, associate, depend); Listing-Stil »latex« wieder wie in den
%             Handreichungen ohne Zeilennummern, nummerierte Variante als »latexnummeriert«;
%             Umlaute und · € ° → ≤ ≥ ≠ – — „ “ … × ² ³ π in allen Listings erlaubt (\lstset literate); \lehrerblase, \schuelerblase; PAP-Stile mit align=center (Zeilenumbruch mit \\);
%             \ipa{...} für Lautschrift; \koordinatensystem[Schritt]{xmin}{xmax}{ymin}{ymax};
%             TikZ-Stile punkt, punktoffen, baumdiagramm, wkeit
% 2026-11-09  TikZ-Stil baumdiagrammvertikal (Baumdiagramm von oben nach unten)
% 2026-11-09  Formularfelder für ausfüllbare PDFs: \feld, \zeichen, \kreuz, \antwort, einheitliches
%             Aussehen über \DefaultOptionsofText/CheckBox/Radio/PopdownBox (nur wirksam in einer Form-Umgebung)
% 2026-11-11  Akzentfarben akzentblau, akzentrot, akzentgruen, akzentorange, akzentlila (Karten,
%             Naturwissenschaften, Schulorganisation); Box »hinweis« mit frei wählbarem Titel
% 2026-11-11  Umgebung »passend«: verkleinert den Inhalt auf \linewidth, aber nur wenn er breiter ist
% 2026-11-11  \oz{+II}{Cu}: Oxidationszahl über ein Elementsymbol (Chemie, auch innerhalb \ce{... $\oz{0}{Zn}$ ...})
% 2026-11-22  Farbige Schachbretter: \schachbrettfarben{braun|blau|gruen|grau|klassisch} (Standard: braun),
%             Farben brett...hell/brett...dunkel und brettmarke (braucht LSBC5 aus der header.tex)
%








%----------------------------------------------------------------------------------------------------%
%%%%%%%%%%%                                   Mathematische Kürzel                                       %%%%%%%%%%
%----------------------------------------------------------------------------------------------------%


%vereinfachte Schreibweise für den Displaystyle
\newcommand{\dis}[1]{\displaystyle{#1}}


%Zahlenbereiche
\newcommand{\N}{\mathbb{N}} %natürliche Zahlen
\newcommand{\Z}{\mathbb{Z}} %ganze Zahlen
\newcommand{\Q}{\mathbb{Q}} %rationale Zahlen
\newcommand{\R}{\mathbb{R}} %reelle Zahlen
\newcommand{\C}{\mathbb{C}} %komplexe Zahlen
\newcommand{\cP}{\mathcal{P}} %Potenzmenge


%Abkürzungen für Analysis
\newcommand{\eps}{\varepsilon} %Kürzel für Epsilon
\newcommand{\limn}{\underset{n\longrightarrow \infty}{\lim}} %Grenzwert für n gegen Unendlich
\newcommand{\limh}{\underset{h\longrightarrow 0}{\lim}} %Grenzwert für h gegen Null
\newcommand{\limxto}[1]{\underset{x\longrightarrow #1}{\lim}} %beliebiger funktionaler Grenzwert x gegen ....
\newcommand{\dx}{\text{\;d} x} %dx Operator für Integrale
\renewcommand{\d}[1]{\text{\;d}#1} %Operator für beliebige Integrationsvariable
\newcommand{\norm}[1]{\left\lVert#1\right\rVert} %\norm{...} für die Normdarstellung ||...||


%Abkürzungen für Analytische Geometrie
\newcommand{\vektor}[3]{\begin{tiny}\begin{pmatrix} \text{\begin{scriptsize}$#1$\end{scriptsize}}\\ \text{\begin{scriptsize}$#2$\end{scriptsize}}\\ \text{\begin{scriptsize}$#3$\end{scriptsize}}\end{pmatrix}\end{tiny}} %dreidimensionaler Vektor
\newcommand{\punktvektor}[2]{\begin{tiny}\begin{pmatrix} \text{\begin{scriptsize}$#1$\end{scriptsize}}\\ \text{\begin{scriptsize}$#2$\end{scriptsize}}\end{pmatrix}\end{tiny}} %zweidimensionaler Vektor


%Abkürzungen für Stochastik
\newcommand{\E}{\mathbb{E}} %Erwartungswert
\newcommand{\V}{\mathbb{V}} %Varianz


%Umgebung für gegeben und gesucht
\newcommand{\gegges}[2]{%
	\begin{multicols}{2}
		\begin{enumerate}
			\item[geg:] #1
			\item[ges:] #2
		\end{enumerate}
	\end{multicols}
}
%----------------------------------------------------------------------------------------------------%









%----------------------------------------------------------------------------------------------------%
%%%%%%%%%%%%                             Verschiedene  hilfreiche Kürzel                                %%%%%%%%%
%----------------------------------------------------------------------------------------------------%


%römische Darstellung einer Zahl \uproman{2021} ergibt MMXXI
\newcommand{\uproman}[1]{\uppercase\expandafter{\romannumeral#1}}


%Automatisch erstellte Lücke wie die Länge des Textes mit \luecke{<Text>}
\newcommand{\luecke}[1]{\hspace{0,1cm}\underline{\hspace{0,25cm}\hphantom{#1}\hspace{0,25cm}}\hspace{0,1cm}}


%fett-kursive Schrift für neu eingeführte Begriffe
\newcommand{\fk}[1]{\textbf{\textit{#1}}}


%Schwierigkeitsstufen für Aufgaben mit \leicht , \mittel , \schwer
\newcommand{\leicht}{\smash{\raisebox{-1pt}{{\color{black}\rule{4pt}{8pt}}\,{\color{lightgray}\rule{4pt}{8pt}}\,{\color{lightgray}\rule{4pt}{8pt}}}}\,}
\newcommand{\mittel}{\smash{\raisebox{-1pt}{{\color{black}\rule{4pt}{8pt}}\,{\color{black}\rule{4pt}{8pt}}\,{\color{lightgray}\rule{4pt}{8pt}}}}\,}
\newcommand{\schwer}{\smash{\raisebox{-1pt}{{\color{black}\rule{4pt}{8pt}}\,{\color{black}\rule{4pt}{8pt}}\,{\color{black}\rule{4pt}{8pt}}}}\,}
%----------------------------------------------------------------------------------------------------%










%----------------------------------------------------------------------------------------------------%
%%%%%%%                                            Verschiedene Farben                                                      %%%%%%%
%----------------------------------------------------------------------------------------------------%


%Formeln in herkömmlichen Farben schreiben
\newcommand{\frot}[1]{\text{\textcolor{red}{$#1$}}} % rote Formel für Lösungen
\newcommand{\fgruen}[1]{\text{\textcolor{green}{$#1$}}} % grüne Formel für Lösungen
\newcommand{\fblau}[1]{\text{\textcolor{blue}{$#1$}}} % blaue Formel für Lösungen
\newcommand{\fmagenta}[1]{\text{\textcolor{magenta}{$#1$}}} % magenta-farbige Formel für Lösungen
\newcommand{\forange}[1]{\text{\textcolor{orange}{$#1$}}} % orangene Formel für Lösungen
\newcommand{\fcyan}[1]{\text{\textcolor{cyan}{$#1$}}} % cyan-farbige Formel für Lösungen
\newcommand{\folive}[1]{\text{\textcolor{olive}{$#1$}}} % olive-farbige Formel für Lösungen


%Custom made colors
\definecolor{plum}{rgb}{0.36078, 0.20784, 0.4}
\definecolor{chameleon}{rgb}{0.30588, 0.60392, 0.023529}
\definecolor{cornflower}{rgb}{0.12549, 0.29020, 0.52941}
\definecolor{scarlet}{rgb}{0.8, 0, 0}
\definecolor{brick}{rgb}{0.64314, 0, 0}
\definecolor{sunrise}{rgb}{0.80784, 0.36078, 0}
\definecolor{lightblue}{rgb}{0.15,0.35,0.75}
\definecolor{carolina}{rgb}{153, 186, 221}

%Akzentfarben für Grafiken, Karten und Diagramme (kräftig, gut unterscheidbar; aufgehellt z.B. akzentblau!18)
\definecolor{akzentblau}{HTML}{2563EB}
\definecolor{akzentrot}{HTML}{DC2626}
\definecolor{akzentgruen}{HTML}{16A34A}
\definecolor{akzentorange}{HTML}{EA580C}
\definecolor{akzentlila}{HTML}{7C3AED}
%----------------------------------------------------------------------------------------------------%









%----------------------------------------------------------------------------------------------------%
%%%%%%%%                            Kopfzeilen für Leistungsüberprüfungen                              %%%%%%%%%
%----------------------------------------------------------------------------------------------------%


%Definition Kopfzeile für Leistungskontrolle mit \lk{Klasse}{Thema}
\newcommand{\lk}[2]{%
	\pagestyle{fancy}
	\fancyhead[L]{%
	\begin{tabular}[b]{l}
		\\Name:\\Punkte:\\Note:
	\end{tabular} 
	}
	\fancyhead[R]{%
	\begin{tabular}[b]{r}
		\\Leistungskontrolle\\Klasse #1\\ \today
	\end{tabular} 
	}
	\fancyhead[C]{%
	\textbf{\large{#2}}\\
	}
	\fancyfoot[C]{}
	\renewcommand{\headrulewidth}{1pt}
	\fancyheadoffset{1cm}
}


%Definition Kopfzeile für Klassenarbeit mit \ka{Klasse}{Thema}
\newcommand{\ka}[2]{%
	\pagestyle{fancy}
	\fancyhead[L]{%
	\begin{tabular}[b]{l}
		\\Name:\\Punkte:\\Note:
	\end{tabular} 
	}
	\fancyhead[R]{%
	\begin{tabular}[b]{r}
		\\Klassenarbeit\\Klasse #1\\ \today
	\end{tabular} 
	}
	\fancyhead[C]{%
	\textbf{\large{#2}}\\
	}
	\fancyfoot[C]{}
	\renewcommand{\headrulewidth}{1pt}
	\fancyheadoffset{1cm}
}


%Definition Kopfzeile für Klausur mit \kl{Klasse}{Thema}
\newcommand{\kl}[2]{%
	\pagestyle{fancy}
	\fancyhead[L]{%
	\begin{tabular}[b]{l}
		\\Name:\\Punkte:\\Note:
	\end{tabular} 
	}
	\fancyhead[R]{%
	\begin{tabular}[b]{r}
		\\Klausur\\Klasse #1\\ \today
	\end{tabular} 
	}
	\fancyhead[C]{%
	\textbf{\large{#2}}\\
	}
	\fancyfoot[C]{}
	\renewcommand{\headrulewidth}{1pt}
	\fancyheadoffset{1cm}
}


%Entfernt den Header für Leistungskontrollen, Klassenarbeiten oder Klausuren
\newcommand{\clearfancyheader}{%
	\fancyhf{}
	\renewcommand{\headrulewidth}{0pt}
	\renewcommand{\headheight}{0pt}
}


%Definition Kopfzeile für Tägliche Übungen mit \tu{Klasse}{Thema}
\newcommand{\tu}[2]{%
	\pagestyle{fancy}
	\fancyhead[L]{%
	\begin{tabular}[b]{l}
		\\Name:\\Punkte:\\Note:
	\end{tabular} 
	}
	\fancyhead[R]{%
	\begin{tabular}[b]{r}
		\\Tägliche Übung\\Klasse #1\\ \today
	\end{tabular} 
	}
	\fancyhead[C]{%
	\textbf{\large{#2}}\\
	}
	\fancyfoot[C]{}
	\renewcommand{\headrulewidth}{1pt}
	\fancyheadoffset{1cm}
}
%----------------------------------------------------------------------------------------------------%








%----------------------------------------------------------------------------------------------------%
%%%%%%%%%                                     Colorboxes für Hervorhebungen                              %%%%%%%%%
%----------------------------------------------------------------------------------------------------%


 %Colorboxes für Lösungen mit \begin{solution}<Text>\end{solution}
\newtcolorbox{solution}{%
        	lower separated=false,%
	colback=gray!5!white,%
	colframe=gray!75!black,%
        	fonttitle=\bfseries,%
	title=Lösung:,%
      	enhanced,%
	breakable,%
        	attach boxed title to top left={xshift=0.5cm,yshift=-2mm}%
}


%Colorboxes für das chronologische Vorgehen 
\newtcolorbox{todo}{%
        	lower separated=false,%
	colback=yellow!5!white,%
	colframe=orange!75!black,%
        	enhanced,%
	breakable%
}


%Colorboxes für Merksätze
\newtcolorbox{merke}{%
        	lower separated=false,%
	colback=red!5!white,%
	colframe=red!75!black,%
        	fonttitle=\bfseries,%
	title=Merke:,%
      	enhanced,%
	breakable,%
        	attach boxed title to top left={xshift=0.5cm,yshift=-2mm}%
}


%Hinweis-Box mit frei wählbarem Titel: \begin{hinweis}[Lizenz] ... \end{hinweis} (Standardtitel: Hinweis)
\newtcolorbox{hinweis}[1][Hinweis]{%
	enhanced,%
	breakable,%
	sharp corners,%
	colback=akzentblau!5,%
	colframe=akzentblau,%
	boxrule=0pt,%
	leftrule=3pt,%
	fonttitle=\bfseries,%
	coltitle=akzentblau!85!black,%
	attach title to upper={\ },%
	title={#1:}%
}


%Colorbox2 für Merksätze 
\newtcolorbox{merke2}{%
        	lower separated=false,%
        	colback=red!40,%
       	title=Merke,%
       	colframe=black,%
        	fonttitle=\bfseries,%
        	colbacktitle=red,%
        	coltitle=black,%
        	enhanced,%
        	breakable,%
        	attach boxed title to top left={xshift=0.5cm,yshift=-2mm},%
        	title=Merke:%
}
%----------------------------------------------------------------------------------------------------%










%----------------------------------------------------------------------------------------------------%
%%%%%%%%%%%%%                                         MC-Fragen                                    %%%%%%%%%%%%%
%----------------------------------------------------------------------------------------------------%

%Definition des Kästchens zum ankreuzen
\newcommand{\casket}{\LARGE{$\square$}}


%Umgebung für horizontale MC-Fragen
\newenvironment{mcquestionh}[1]{%
	#1%
	\begin{center}\begin{tabular}{>{\centering\arraybackslash}p{2.5cm}  >{\centering\arraybackslash}p{2.5cm} >{\centering\arraybackslash}p{2.5cm} >{\centering\arraybackslash}p{2.5cm} >{\centering\arraybackslash}p{2.5cm}}
	\casket & \casket & \casket & \casket & \casket\medskip\\
	}
	%hier wird die fehlende Tabellenzeile angegeben. Aufruf erfolgt durch:
	%\begin{mcquestionh}{<Frage>}
	%   	\answerh{<Antwort>} & \answerh{<Antwort>} & \answerh{<Antwort>} & \answerh{<Antwort>} & \answerh{<Antwort>}
	%\end{mcquestionh}
	{
	\\
	\end{tabular}\end{center}
}


%Antwort für horizontale MC-Fragen mit \answerh{<Antwort>}
\newcommand{\answerh}[1]{%
	\parbox{2.5cm}{\centering{#1}}
}


%Umgebung für vertikale MC-Fragen
\newenvironment{mcquestionv}[1]{%
    	#1%
    	\begin{enumerate}[\raisebox{-3pt}{\casket}]
    	}
    	%hier wird das fehlende \item angegeben werden. Aufruf erfolgt durch:
    	%\begin{mcquestionv}{<Frage>}
    	%   	\answerv{<Antwort>}
    	%  	\answerv{<Antwort>}
    	%   	\answerv{<Antwort>}
    	%   	\answerv{<Antwort>}
    	%   	\answerv{<Antwort>}
    	%\end{mcquestionv}
    	{
    	\end{enumerate}
}


%Antwort für vertikale MC-Frage mit \answerv{<Antwort>}
\newcommand{\answerv}[1]{%
	\item \hspace{0.1cm} #1
}
%----------------------------------------------------------------------------------------------------%









%----------------------------------------------------------------------------------------------------%
%%%%%%%%%%%                                          Tikz Settings                                          %%%%%%%%%%%%
%----------------------------------------------------------------------------------------------------%


%Koordinatensystem im Standardstil der Mathe-Arbeitsblätter (Gitter gray!30, Achsen -latex, Beschriftung \tiny)
%\koordinatensystem[Schritt]{xmin}{xmax}{ymin}{ymax};
%             TikZ-Stile punkt, punktoffen, baumdiagramm, wkeit
% 2026-11-09  TikZ-Stil baumdiagrammvertikal (Baumdiagramm von oben nach unten) – ganze Zahlen, Bereich muss 0 enthalten,
%Striche und Zahlen bei allen Vielfachen von Schritt (Standard 2). Innerhalb von tikzpicture[scale=0.42] verwenden.
\newcommand{\koordinatensystem}[5][2]{%
	\foreach \kosyx in {#2,...,#3} \draw[gray!30, thin] (\kosyx,#4) -- (\kosyx,#5);
	\foreach \kosyy in {#4,...,#5} \draw[gray!30, thin] (#2,\kosyy) -- (#3,\kosyy);
	\draw[-latex, thick] ({#2-0.3},0) -- ({#3+0.5},0) node[right] {$x$};
	\draw[-latex, thick] (0,{#4-0.3}) -- (0,{#5+0.5}) node[above] {$y$};
	\foreach \kosyx in {#2,...,#3} {%
		\pgfmathtruncatemacro{\kosyrest}{mod(abs(\kosyx),#1)}%
		\ifnum\kosyrest=0 \ifnum\kosyx=0 \else
			\draw (\kosyx,0.1) -- (\kosyx,-0.1) node[below, font=\tiny] {$\kosyx$};
		\fi\fi}
	\foreach \kosyy in {#4,...,#5} {%
		\pgfmathtruncatemacro{\kosyrest}{mod(abs(\kosyy),#1)}%
		\ifnum\kosyrest=0 \ifnum\kosyy=0 \else
			\draw (0.1,\kosyy) -- (-0.1,\kosyy) node[left, font=\tiny] {$\kosyy$};
		\fi\fi}
}


%optionale Eigenschaft für ein  Tikzpicture, damit dieses bündig am oberen Rand einer enumerate-Auflistung steht (benötigt tikzlibrary calc!!!)
\tikzset{
  topline/.style={baseline={([yshift=#1]current bounding box.north)}},
  topline/.default=-\ht\strutbox
}


%Punkte in Graphen und Figuren: \node[punkt, label=above:{$P$}] at (2,3) {}; (offen: punktoffen)
%Als Knoten, damit die Größe nicht mit scale=0.42 mitschrumpft (circle (3pt) würde mitskaliert)
\tikzset{
  punkt/.style={circle, fill, inner sep=0pt, minimum size=4pt},
  punktoffen/.style={circle, draw, fill=white, inner sep=0pt, minimum size=4pt}
}


%Baumdiagramme (Stochastik) von links nach rechts: \begin{tikzpicture}[baumdiagramm] \node[punkt]{} child{...}
%Wahrscheinlichkeiten an den Ästen: edge from parent node[wkeit, above] {$0{,}6$} (oberer Ast above, unterer below)
\tikzset{
  baumdiagramm/.style={grow'=right, sloped, level distance=2.6cm,
    level 1/.style={sibling distance=3cm}, level 2/.style={sibling distance=1.5cm},
    level 3/.style={sibling distance=0.75cm}, edge from parent/.style={draw, thin}},
  wkeit/.style={font=\scriptsize, inner sep=2pt}
}


%Baumdiagramme von oben nach unten: \begin{tikzpicture}[baumdiagrammvertikal] \node[punkt]{} child{...}
%erstes child = linker Ast; Wahrscheinlichkeit links mit node[wkeit, left], rechts mit node[wkeit, right]
\tikzset{
  baumdiagrammvertikal/.style={grow=down, level distance=1.6cm,
    level 1/.style={sibling distance=4.8cm}, level 2/.style={sibling distance=2.4cm},
    level 3/.style={sibling distance=1.2cm}, edge from parent/.style={draw, thin}}
}
%----------------------------------------------------------------------------------------------------%









%----------------------------------------------------------------------------------------------------%
%%%%%%%%%%%                                     Handreichungen                                        %%%%%%%%%%%
%----------------------------------------------------------------------------------------------------%


%Schriftzug TikZ im Fließtext (Leerzeichen danach ist eingebaut)
\newcommand{\TikZ}{Ti\textit{k}Z }


%Oxidationszahl über ein Elementsymbol: $\oz{+II}{Cu}$, in \ce{...} zwischen Dollarzeichen: \ce{$\oz{0}{Zn}$ + ...}
\newcommand{\oz}[2]{\overset{\scriptscriptstyle\mathrm{#1}}{\mathrm{\ce{#2}}}}


%Umgebung passend: Inhalt (z.B. eine große tikzpicture) höchstens so breit wie die Zeile setzen.
%Ist er schmaler, bleibt er unverändert. \begin{passend} ... \end{passend} (Inhalt nicht als Argument, # erlaubt)
\newsavebox{\passendbox}
\newenvironment{passend}{\begin{lrbox}{\passendbox}}{\end{lrbox}%
	\ifdim\wd\passendbox>\linewidth\resizebox{\linewidth}{!}{\usebox{\passendbox}}\else\usebox{\passendbox}\fi}


%Box für die Ausgabe unter einem Code-Beispiel: erst lstlisting[style=latex], darunter \begin{ausgabe} ... \end{ausgabe}
\newtcolorbox{ausgabe}{%
	colback=blue!5,%
	colframe=blue!50!black,%
	title=Ausgabe,%
	fonttitle=\bfseries%
}
%----------------------------------------------------------------------------------------------------%









%----------------------------------------------------------------------------------------------------%
%%%%%%%%%%%                                        Lautschrift                                         %%%%%%%%%%%
%----------------------------------------------------------------------------------------------------%


%Lautschrift in eckigen Klammern mit \ipa{"a:b@nt} – schaltet das babel-Kürzel " vorher aus (sonst wird "a zu ä)
\makeatletter
\newcommand{\ipa}{\begingroup\shorthandoff{"}\ipa@aux}
\newcommand{\ipa@aux}[1]{[\textipa{#1}]\endgroup}
\makeatother
%----------------------------------------------------------------------------------------------------%









%----------------------------------------------------------------------------------------------------%
%%%%%%%%%%%                                       Sprechblasen                                         %%%%%%%%%%%
%----------------------------------------------------------------------------------------------------%


%Sprechblase mit Zipfel links unten (Lehrkraft) mit \lehrerblase[TikZ-Optionen]{Text} (benötigt shapes.callouts)
\newcommand{\lehrerblase}[2][]{%
  \tikz[baseline=(B.base)]{%
    \node[rectangle callout, rounded corners=6pt, callout relative pointer={(-13mm,-7mm)},
      draw, thick, fill=blue!6, align=left, inner sep=3.5mm, #1] (B) {#2};}%
}


%Sprechblase mit Zipfel rechts unten (Schülerin/Schüler) mit \schuelerblase[TikZ-Optionen]{Text}
\newcommand{\schuelerblase}[2][]{%
  \tikz[baseline=(B.base)]{%
    \node[rectangle callout, rounded corners=6pt, callout relative pointer={(13mm,-7mm)},
      draw, thick, fill=green!6, align=left, inner sep=3.5mm, #1] (B) {#2};}%
}
%----------------------------------------------------------------------------------------------------%









%----------------------------------------------------------------------------------------------------%
%%%%%%%%%%%                            TikZ-Stile für Programmablaufpläne                            %%%%%%%%%%%
%----------------------------------------------------------------------------------------------------%


%Ein-/Ausgabe als Parallelogramm, Operation als Rechteck, Verzweigung als Raute (benötigt shapes.geometric)
%align=center, damit \\ im Knotentext eine neue Zeile ergibt (ohne align wird \\ stillschweigend verschluckt)
\tikzstyle{input}=[trapezium, trapezium left angle=60, trapezium right angle=120, fill=green!20, draw=green, align=center]
\tikzstyle{operation}=[rectangle, fill=orange!20, draw=orange, align=center]
\tikzstyle{condition}=[diamond, fill=magenta!20, draw=magenta, align=center]
\tikzstyle{output}=[trapezium, trapezium left angle=60, trapezium right angle=120, fill=red!20, draw=red, align=center]
%----------------------------------------------------------------------------------------------------%









%----------------------------------------------------------------------------------------------------%
%%%%%%%%%%%                            TikZ-Stile für UML-Klassendiagramme                           %%%%%%%%%%%
%----------------------------------------------------------------------------------------------------%


%Klasse mit drei Teilen: Name, Attribute, Methoden (benötigt shapes.multipart)
\tikzstyle{umlclass} = [rectangle split, rectangle split parts=3, rectangle split part align={center,left,left}, draw, thick, text width=4cm, minimum height=2cm]


%Beziehungen (benötigt arrows.meta)
\tikzstyle{inherit} = [-{Triangle[open]}, thick]	%Vererbung
\tikzstyle{compose} = [-{Diamond}, thick]	%Komposition
\tikzstyle{aggregate} = [-{Diamond[open]}, thick]	%Aggregation
\tikzstyle{associate} = [-, thick]	%Assoziation
\tikzstyle{depend} = [-, thick, dashed]	%Abhängigkeit
%----------------------------------------------------------------------------------------------------%








%----------------------------------------------------------------------------------------------------%
%%%%%%%%%%%                                         Listing Settings                                       %%%%%%%%%%%%
%----------------------------------------------------------------------------------------------------%


%Umlaute und häufige Sonderzeichen in ALLEN Listings (Python, Java, ...), sonst Fehler »Invalid UTF-8 byte sequence«
\lstset{literate=%
	{ä}{{\"a}}1 {ö}{{\"o}}1 {ü}{{\"u}}1 {Ä}{{\"A}}1 {Ö}{{\"O}}1 {Ü}{{\"U}}1 {ß}{{\ss}}1%
	{·}{{$\cdot$}}1 {€}{{\euro}}1 {°}{{\degree}}1 {→}{{$\rightarrow$}}2 {≤}{{$\leq$}}1 {≥}{{$\geq$}}1 {≠}{{$\neq$}}1%
	{–}{{--}}1 {—}{{---}}1 {„}{{\glqq}}1 {“}{{\grqq}}1 {…}{{\dots}}3 {×}{{$\times$}}1 {²}{{$^2$}}1 {³}{{$^3$}}1 {π}{{$\pi$}}1%
}


%Einstellung der Codeumgebung für LaTeX-Code-Beispiele (Stil der Handreichungen: ohne Zeilennummern)
%Verwendung: \begin{lstlisting}[style=latex] ... \end{lstlisting}, darunter \begin{ausgabe} ... \end{ausgabe}
\lstdefinestyle{latex}{%
  	language=TeX,%
  	frame=single,%
  	backgroundcolor=\color{yellow!20},%
  	rulecolor=\color{lightgray!40},%
  	breaklines=true,%
  	basicstyle=\scriptsize\ttfamily,%
  	keywordstyle=\bfseries\color{purple!40!black},%
  	commentstyle=\itshape\color{green!40!black},%
  	identifierstyle=\color{blue},%
	stringstyle=\color{black},%
	captionpos=b,%
	tabsize=2,%
	showspaces=false,%
	showstringspaces=false,%
  	showtabs=false,%
	deletekeywords={document},%
	morekeywords={begin}%
}


%Variante mit Zeilennummern: \begin{lstlisting}[style=latexnummeriert]
\lstdefinestyle{latexnummeriert}{%
	style=latex,%
	numbers=left,%
	numberstyle=\tiny\color{gray!90!black},%
  	numbersep=-5pt%
}
%----------------------------------------------------------------------------------------------------%






%----------------------------------------------------------------------------------------------------%
%%%%%%%%%%%                          Formularfelder (ausfüllbare PDFs)                          %%%%%%%%%%%
%----------------------------------------------------------------------------------------------------%
% Nur innerhalb von \begin{Form} ... \end{Form} (hyperref) verwendbar. Für Optionsfelder (radio)
% im Dokument zusätzlich \hypersetup{unicode=false}, sonst bekommen sie kaputte Namen.
% Jedes Feld steckt in einer eigenen Box (\raisebox/\hbox): Ohne sie übernimmt pdfTeX die Höhe
% der umgebenden Zeile, in Tabellen werden die Felder dann so hoch wie die ganze Zeile.

%Aussehen aller Felder: blauer Rand, hellblauer Hintergrund, Schrift 11 pt, wird mitgedruckt
\renewcommand{\DefaultOptionsofText}{print, bordercolor={0 0.353 0.612}, backgroundcolor={0.94 0.96 1}, charsize=11pt}
\renewcommand{\DefaultOptionsofCheckBox}{print, bordercolor={0 0.353 0.612}, backgroundcolor={0.94 0.96 1}, charsize=11pt}
\renewcommand{\DefaultOptionsofRadio}{print, notoggletooff, bordercolor={0 0.353 0.612}, backgroundcolor={0.94 0.96 1}, charsize=11pt}
\renewcommand{\DefaultOptionsofPopdownBox}{print, bordercolor={0 0.353 0.612}, backgroundcolor={0.94 0.96 1}, charsize=11pt}
\renewcommand{\DefaultHeightofText}{0.6cm}

%\feld[Breite]{Name} – einzeiliges Textfeld, sitzt auf der Grundlinie (Standardbreite 2,4 cm)
\newcommand{\feld}[2][2.4cm]{\raisebox{-0.17cm}{\TextField[name=#2, width=#1]{}}}

%\zeichen{Name} – Feld für genau ein Zeichen, zentriert (z. B. E/V/A, Buchstabe einer Reihenfolge)
\newcommand{\zeichen}[1]{\raisebox{-0.12cm}{\TextField[name=#1, width=0.8cm, height=0.5cm, maxlen=1, align=1]{}}}

%\kreuz{Name} – Ankreuzkästchen
\newcommand{\kreuz}[1]{\raisebox{-0.08cm}{\CheckBox[name=#1, width=0.42cm, height=0.42cm]{}}}

%\antwort[Höhe]{Name} – mehrzeiliges Antwortfeld über die ganze Zeilenbreite (Standardhöhe 3,2 cm ≈ 4 Zeilen)
\newcommand{\antwort}[2][3.2cm]{\par\vspace{0.15cm}\noindent\hbox{\TextField[name=#2, multiline=true, width=\linewidth, height=#1]{}}}
%----------------------------------------------------------------------------------------------------%




%----------------------------------------------------------------------------------------------------%
%%%%%%%%%%%                                   Schach (farbige Bretter)                                   %%%%%%%%%%
%----------------------------------------------------------------------------------------------------%
% Die Standardschrift von skak schraffiert die dunklen Felder; Figur und Schraffur sind dort ein
% Zeichen und lassen sich nicht getrennt einfärben. Deshalb: Schrift skaknew in der Kodierung LSBC5
% (zeichnet nur die Figuren, weiße Figuren weiß gefüllt), die Felder malt chessboard als Flächen
% darunter. Eigene Markierungen mit colorbackfields liegen über den Feldern und unter den Figuren.

%Feldfarben der Brettschemata (hell/dunkel) und Markierungsfarbe für den letzten Zug
\definecolor{brettblauhell}{HTML}{EEF3FD} %Schema blau, helle Felder
\definecolor{brettblaudunkel}{HTML}{A9B4CC} %Schema blau, dunkle Felder
\definecolor{brettbraunhell}{HTML}{F0D9B5} %Schema braun (klassisches Holzbrett), helle Felder
\definecolor{brettbraundunkel}{HTML}{B58863} %Schema braun, dunkle Felder
\definecolor{brettgruenhell}{HTML}{EEEED2} %Schema gruen (Turnierbrett), helle Felder
\definecolor{brettgruendunkel}{HTML}{769656} %Schema gruen, dunkle Felder
\definecolor{brettgrauhell}{HTML}{F2F2F2} %Schema grau (für Schwarz-Weiß-Druck), helle Felder
\definecolor{brettgraudunkel}{HTML}{B4B4B4} %Schema grau, dunkle Felder
\definecolor{brettmarke}{HTML}{F6D55C} %Start- und Zielfeld eines Zuges hervorheben: pgfstyle=color, color=brettmarke, colorbackfields={e2,e4}

% Setzt die Voreinstellungen von \setchessboard zurück – eigene \setchessboard{...} deshalb DANACH schreiben.
%\schachbrettfarben{Schema} – Farben aller folgenden Schachbretter: braun (Standard), blau, gruen, grau oder klassisch (schraffiert wie in skak)
\makeatletter
\let\brett@grundeinstellung\board@val@savedkeylist
\newcommand{\schachbrettfarben}[1]{%
	\let\board@val@savedkeylist\brett@grundeinstellung
	\ifthenelse{\equal{#1}{klassisch}}{}%
	{\setchessboard{boardfontfamily=skaknew, boardfontencoding=LSBC5,
		piececolor=black, piecemaskcolor=white, setfontcolors,
		pgfstyle=color, opacity=1, color=brett#1hell, colorbackboard, color=brett#1dunkel, colorblackbackfields,
		color=black, labelfont=\sffamily\color{black!65}}}}
\makeatother
\schachbrettfarben{braun}
%----------------------------------------------------------------------------------------------------%


% Chemisches Wörterbuch Deutsch – Englisch mit Lautschrift
% Skills: latex-grundstil, latex-lautschrift (tipa, \ipa), latex-chemie (mhchem \ce, chemfig),
%         latex-funktionsgraphen, latex-tikz-grundlagen
% Kompilieren mit pdflatex (dreimal, wegen Inhaltsverzeichnis und Seitenzahlen im Register).

%----------------------------------------------------------------------------------------------------%
% Dokumentspezifische Befehle (nur für dieses Wörterbuch, deshalb lokal und nicht in newcommands.tex)
%----------------------------------------------------------------------------------------------------%

% \eintrag{Marke}{Artikel}{Begriff}{english term}{Lautschrift EN}{Formel/Skizze}
% Schaltet das babel-Kürzel " aus, BEVOR die Argumente gelesen werden – dadurch darf die Lautschrift
% wie bei \ipa{...} mit " (Hauptbetonung) und "" (Nebenbetonung) geschrieben werden.
\makeatletter
\newcommand\eintrag{\begingroup\shorthandoff{"}\eintrag@aux}
\newcommand\eintrag@aux[6]{\endgroup
	#2 \textbf{#3}\phantomsection\label{e:#1}
	& \textbf{#4}\newline{\small\color{black!70}\ipa{#5}}
	& #6 \\ \midrule}
% \lies{Ausdruck}{english reading}{deutsche Lesart}
\newcommand\lies{\begingroup\shorthandoff{"}\lies@aux}
\newcommand\lies@aux[3]{\endgroup #1 & #2\newline{\color{black!60}#3} \\[4.5pt]}
\makeatother

% Tabellenkopf für jedes Kapitel in Teil I (wird auf jeder Folgeseite wiederholt)
\newenvironment{woerter}{%
	\arrayrulecolor{black!20}%
	\begin{longtable}{>{\raggedright\arraybackslash}m{5.3cm} >{\raggedright\arraybackslash}m{6.1cm} >{\centering\arraybackslash}m{4.6cm}}
		\rowcolor{cornflower!15}\textbf{Deutsch} & \textbf{Englisch} & \textbf{Formel \,/\, Skizze} \\ \midrule
		\endhead}{%
	\end{longtable}\arrayrulecolor{black}}

% Liste ohne Linien und Rahmen (Teil II): links der Ausdruck, rechts Englisch (schwarz) und Deutsch (grau)
\newenvironment{lesen}{%
	\begin{longtable}{@{}>{\raggedright\arraybackslash}p{4.6cm} >{\raggedright\arraybackslash}p{12.1cm}@{}}}{%
	\end{longtable}\vspace{-1.2em}}

\newcommand{\anmerkung}[1]{\par\smallskip\noindent{\small\sloppy\textbf{Hinweis:} #1\par}\smallskip}
% Verhindert, dass eine Zwischenüberschrift allein am Seitenende steht
\newcommand{\platz}{\par\ifdim\dimexpr\pagegoal-\pagetotal\relax<5.5cm \newpage\fi}
% Einheit aufrecht im Mathemodus
\newcommand{\eh}[1]{\,\mathrm{#1}}

% Kleine Skizzen: Gegenstand blau, Hervorhebung (der Begriff selbst) rot, Hilfslinien gestrichelt
\tikzset{
	fig/.style={blue, very thick},
	mass/.style={red, thick},
	hilf/.style={dashed},
	achse/.style={-{Stealth[length=4pt]}, thin},
	kraft/.style={red, very thick, -{Stealth[length=5pt]}},
}
\newcommand{\skizze}[1]{\tikz[baseline=(current bounding box.center), scale=0.45, line cap=round,
	line join=round, every node/.style={font=\scriptsize}, outer sep=0pt]{#1}}
\newcommand{\achsen}[6]{\draw[achse] (#1,0) -- (#2,0) node[below left] {$#5$}; \draw[achse] (0,#3) -- (0,#4) node[below left] {$#6$};}

% Strukturformel in kleiner Größe
\newcommand{\struktur}[1]{{\footnotesize\setchemfig{atom sep=1.7em, bond offset=1.5pt}\chemfig{#1}}}
\newcommand{\fl}[2]{\fill[blue!12] #1; \draw[blue!60, thin] #2;}

% Laborgeräte
\newcommand{\skReagenzglas}{\skizze{\fill[blue!12] (-0.25,1.1) -- (-0.25,0.25) arc[start angle=180, end angle=360, radius=0.25] -- (0.25,1.1) -- cycle;
	\draw[fig] (-0.25,2.4) -- (-0.25,0.25) arc[start angle=180, end angle=360, radius=0.25] -- (0.25,2.4);}}
\newcommand{\becher}[2]{\fill[blue!12] (#1-0.9,#2) rectangle (#1+0.9,#2+0.9); \draw[fig] (#1-1.05,#2+1.7) -- (#1-0.9,#2+1.6) -- (#1-0.9,#2) -- (#1+0.9,#2) -- (#1+0.9,#2+1.7);}
\newcommand{\skBecherglas}{\skizze{\becher{0}{0} \foreach \y in {0.5,0.8,1.1} {\draw[blue, thin] (0.9,\y) -- (0.55,\y);}}}
\newcommand{\skErlenmeyer}{\skizze{\fill[blue!12] (-0.78,0.6) -- (-1.1,0) -- (1.1,0) -- (0.78,0.6) -- cycle;
	\draw[fig] (-0.3,2.2) -- (-0.3,1.5) -- (-1.1,0) -- (1.1,0) -- (0.3,1.5) -- (0.3,2.2);}}
\newcommand{\skRundkolben}{\skizze{\begin{scope}\clip (-1,0) rectangle (1,0.8); \fill[blue!12] (0,0.9) circle (0.9);\end{scope}
	\draw[fig] (-0.25,2.5) -- (-0.25,1.765) arc[start angle=106.1, end angle=433.9, radius=0.9] -- (0.25,2.5);}}
\newcommand{\skMesszylinder}{\skizze{\fill[blue!12] (-0.3,0.15) rectangle (0.3,1.5); \draw[fig] (-0.3,2.8) -- (-0.3,0.15) -- (0.3,0.15) -- (0.3,2.8); \draw[fig] (-0.7,0.1) -- (0.7,0.1);
	\foreach \y in {0.5,0.8,...,2.6} {\draw[blue, thin] (0.3,\y) -- (0.05,\y);}}}
\newcommand{\trichter}[2]{\draw[fig] (#1-0.9,#2+1.5) -- (#1-0.1,#2+0.7) -- (#1-0.1,#2) (#1+0.9,#2+1.5) -- (#1+0.1,#2+0.7) -- (#1+0.1,#2+0.1);}
\newcommand{\skTrichter}{\skizze{\trichter{0}{0}}}
\newcommand{\brenner}[2]{\fill[orange!70] (#1,#2+1.3) .. controls (#1-0.35,#2+1.6) and (#1-0.2,#2+2.0) .. (#1,#2+2.4) .. controls (#1+0.2,#2+2.0) and (#1+0.35,#2+1.6) .. cycle;
	\fill[blue!50] (#1,#2+1.3) .. controls (#1-0.15,#2+1.45) and (#1-0.1,#2+1.65) .. (#1,#2+1.85) .. controls (#1+0.1,#2+1.65) and (#1+0.15,#2+1.45) .. cycle;
	\draw[fig] (#1-0.17,#2+0.2) rectangle (#1+0.17,#2+1.3); \draw[fig] (#1-0.6,#2) rectangle (#1+0.6,#2+0.2); \draw[fig, thick] (#1-0.6,#2+0.1) -- (#1-1.1,#2+0.1);}
\newcommand{\skBrenner}{\skizze{\brenner{0}{0}}}
\newcommand{\skDreifuss}{\skizze{\draw[fig] (-0.9,1.6) -- (0.9,1.6) (-0.8,1.6) -- (-1.1,0) (0.8,1.6) -- (1.1,0) (0,1.6) -- (0,0.15); \draw[black!60, thick, densely dashed] (-1,1.72) -- (1,1.72);}}
\newcommand{\skAbdampfschale}{\skizze{\draw[fig] (-1.1,0.8) arc[start angle=180, end angle=360, x radius=1.1, y radius=0.8] (-1.25,0.8) -- (1.25,0.8);}}
\newcommand{\skTiegel}{\skizze{\draw[fig] (-0.7,1.2) -- (-0.4,0) -- (0.4,0) -- (0.7,1.2) (-0.8,1.2) -- (0.8,1.2);}}
\newcommand{\skPipette}{\skizze{\draw[fig, thick] (-0.06,2.8) -- (-0.06,1.9) .. controls (-0.3,1.7) and (-0.3,1.3) .. (-0.06,1.1) -- (-0.06,0.3) -- (0,0) -- (0.06,0.3) -- (0.06,1.1) .. controls (0.3,1.3) and (0.3,1.7) .. (0.06,1.9) -- (0.06,2.8);}}
\newcommand{\buerette}[2]{\fill[blue!12] (#1-0.12,#2+0.9) rectangle (#1+0.12,#2+2.3); \draw[fig, thick] (#1-0.12,#2+3) -- (#1-0.12,#2+0.6) -- (#1-0.04,#2+0.45) -- (#1-0.04,#2) (#1+0.12,#2+3) -- (#1+0.12,#2+0.6) -- (#1+0.04,#2+0.45) -- (#1+0.04,#2);
	\draw[red, very thick] (#1-0.35,#2+0.52) -- (#1+0.35,#2+0.52); \foreach \k in {0,...,6} {\draw[blue, thin] (#1+0.12,#2+1+0.3*\k) -- (#1-0.02,#2+1+0.3*\k);}}
\newcommand{\skBuerette}{\skizze{\buerette{0}{0}}}
\newcommand{\skSpatel}{\skizze{\draw[fig] (0,0) -- (2.3,0.9); \draw[blue, line width=2.5pt] (2.3,0.9) -- (2.9,1.13);}}
\newcommand{\skMoerser}{\skizze{\draw[fig] (-1,1) arc[start angle=180, end angle=360, x radius=1, y radius=0.9] (-1.1,1) -- (1.1,1) (-0.4,0.1) -- (0.4,0.1);
	\draw[blue, line width=3pt, line cap=round] (0.1,0.5) -- (0.9,1.9);}}
\newcommand{\skUhrglas}{\skizze{\draw[fig] (-1.2,0.35) arc[start angle=200, end angle=340, x radius=1.277, y radius=0.9];}}
\newcommand{\stativ}{\draw[black, line width=2pt] (-1.2,0) -- (0.6,0); \draw[black, line width=1.5pt] (-0.8,0) -- (-0.8,3.6);}
\newcommand{\skStativ}{\skizze{\stativ \draw[black, thick] (-0.8,2.6) -- (0.5,2.6); \draw[fig, thick] (0.5,2.45) rectangle (0.9,2.75);}}
\newcommand{\skThermometer}{\skizze{\draw[fig, thick, rounded corners=2pt] (-0.15,0.3) rectangle (0.15,2.8); \filldraw[fill=red, draw=blue, thick] (0,0.1) circle (0.33);
	\fill[red] (-0.07,0.3) rectangle (0.07,1.7); \foreach \y in {0.8,1.2,...,2.5} {\draw[thin] (0.15,\y) -- (0.4,\y);}}}
\newcommand{\skKuehler}{\skizze{\begin{scope}[rotate=-20] \draw[fig, fill=blue!10] (0,-0.3) rectangle (3,0.3); \draw[fig, thick] (-0.6,-0.08) -- (3.6,-0.08) (-0.6,0.08) -- (3.6,0.08);
	\draw[fig, thick] (0.4,0.3) -- (0.4,0.7) (2.6,-0.3) -- (2.6,-0.7);\end{scope}}}
\newcommand{\skSchutzbrille}{\skizze{\draw[fig, rounded corners=3pt] (-1.3,0) rectangle (-0.15,0.7) (0.15,0) rectangle (1.3,0.7); \draw[fig] (-0.15,0.45) -- (0.15,0.45) (-1.3,0.5) -- (-1.7,0.7) (1.3,0.5) -- (1.7,0.7);}}
\newcommand{\skWaage}{\skizze{\draw[fig] (-1.2,0) rectangle (1.2,0.6); \draw[fig] (-0.9,0.8) -- (0.9,0.8) (0,0.6) -- (0,0.8); \node at (0,0.3) {\tiny 12{,}5\,g};}}
\newcommand{\skFiltration}{\skizze{\becher{0}{0} \trichter{0}{1.3} \draw[mass, hilf] (-0.75,2.72) -- (0,1.95) -- (0.75,2.72); \fill[black!55] (-0.3,2.3) -- (0,2.0) -- (0.3,2.3) -- cycle;}}
\newcommand{\skErhitzen}{\skizze{\brenner{0}{0} \begin{scope}[shift={(0,2.5)}, rotate=-35] \fill[blue!12] (-0.2,0.8) -- (-0.2,0.2) arc[start angle=180, end angle=360, radius=0.2] -- (0.2,0.8) -- cycle;
	\draw[fig] (-0.2,2) -- (-0.2,0.2) arc[start angle=180, end angle=360, radius=0.2] -- (0.2,2);\end{scope}}}

% Teilchen, Atome, Bindungen
\newcommand{\rahmen}{\draw[thin] (0,0) rectangle (2.1,1.5);}
\newcommand{\punkte}[3]{\foreach \p in {#3} {\fill[#1] \p circle (#2);}}
\newcommand{\skElement}{\skizze{\rahmen \punkte{blue}{2.8pt}{(0.3,0.3),(0.9,0.5),(1.6,0.3),(0.5,1.0),(1.2,1.1),(1.8,0.9)}}}
\newcommand{\molek}[2]{\fill[blue] (#1,#2) circle (2.8pt); \fill[red] (#1+0.22,#2+0.1) circle (2.8pt);}
\newcommand{\skVerbindung}{\skizze{\rahmen \molek{0.25}{0.3} \molek{0.95}{0.45} \molek{1.6}{0.25} \molek{0.45}{1.0} \molek{1.2}{1.05}}}
\newcommand{\skGemisch}{\skizze{\rahmen \punkte{blue}{2.8pt}{(0.3,0.3),(1.6,0.3),(1.2,1.1)} \punkte{red}{2.8pt}{(0.9,0.55),(0.5,1.05)} \punkte{green!55!black}{2.8pt}{(1.8,0.95),(0.95,0.2)}}}
\newcommand{\skReinstoff}{\skVerbindung}
\newcommand{\skLoesung}{\skizze{\becher{0}{0} \punkte{red}{1.6pt}{(-0.6,0.2),(-0.2,0.6),(0.3,0.25),(0.65,0.7),(-0.55,0.75),(0.1,0.45)}}}
\newcommand{\skSuspension}{\skizze{\becher{0}{0} \punkte{black!60}{2.2pt}{(-0.6,0.1),(-0.3,0.12),(0,0.1),(0.3,0.13),(0.6,0.1),(-0.45,0.28),(0.15,0.3),(0.45,0.5)}}}
\newcommand{\skSchalen}{\skizze{\fill[red] (0,0) circle (0.22); \foreach \r in {0.5,0.9,1.3} {\draw[thin] (0,0) circle (\r);}
	\punkte{blue}{2pt}{(90:0.5),(270:0.5),(0:0.9),(45:0.9),(90:0.9),(135:0.9),(180:0.9),(225:0.9),(270:0.9),(315:0.9),(30:1.3)}}}
\newcommand{\skAussenelektron}{\skizze{\fill[black!40] (0,0) circle (0.22); \foreach \r in {0.5,0.9,1.3} {\draw[thin, black!40] (0,0) circle (\r);}
	\punkte{black!40}{2pt}{(90:0.5),(270:0.5),(0:0.9),(45:0.9),(90:0.9),(135:0.9),(180:0.9),(225:0.9),(270:0.9),(315:0.9)} \fill[red] (30:1.3) circle (3pt);}}
\newcommand{\skPSE}{\skizze{\foreach \x in {0,...,7} {\foreach \y in {0,...,3} {\draw[thin, black!50] (0.3*\x,0.3*\y) rectangle ++(0.3,0.3);}}
	\draw[thin, black!50] (0,1.2) rectangle (0.3,1.5) (2.1,1.2) rectangle (2.4,1.5);}}
\newcommand{\pse}[1]{\skizze{\foreach \x in {0,...,7} {\foreach \y in {0,...,3} {\draw[thin, black!50] (0.3*\x,0.3*\y) rectangle ++(0.3,0.3);}}
	\draw[thin, black!50] (0,1.2) rectangle (0.3,1.5) (2.1,1.2) rectangle (2.4,1.5); \fill[red, opacity=0.6] #1;}}
\newcommand{\skGitter}{\skizze{\foreach \x in {0,...,3} {\foreach \y in {0,...,2} {\pgfmathparse{int(mod(\x+\y,2))}
	\ifnum\pgfmathresult=0 \fill[red] (0.5*\x,0.5*\y) circle (4pt); \else \fill[blue] (0.5*\x,0.5*\y) circle (3pt); \fi}}}}
\newcommand{\skMetall}{\skizze{\foreach \x in {0,...,3} {\foreach \y in {0,...,2} {\filldraw[fill=red!25, draw=red] (0.6*\x,0.6*\y) circle (5pt); \node[font=\tiny] at (0.6*\x,0.6*\y) {+};}}
	\punkte{blue}{1.5pt}{(0.3,0.3),(0.9,0.25),(1.5,0.35),(0.3,0.9),(0.95,0.95),(1.5,0.85),(0.6,0.62),(1.2,0.58)}}}
\newcommand{\skAtombindung}{\skizze{\draw[fig, fill=blue!10, fill opacity=0.6] (-0.45,0) circle (0.7) (0.45,0) circle (0.7); \punkte{red}{2.2pt}{(0,0.15),(0,-0.15)} \node at (-0.7,0) {H}; \node at (0.7,0) {H};}}
\newcommand{\skDipol}{\skizze{\fill[red] (0,0) circle (0.42); \fill[black!30] (-0.55,-0.42) circle (0.27) (0.55,-0.42) circle (0.27); \node at (0,0.75) {$\delta^-$}; \node at (-1.05,-0.6) {$\delta^+$}; \node at (1.05,-0.6) {$\delta^+$};}}

% Reaktionen
\newcommand{\skExotherm}{\skizze{\achsen{-0.3}{3.8}{-0.3}{2.7}{}{E}\draw[fig, thick] (0.2,1.5) -- (1,1.5) (2.6,0.5) -- (3.4,0.5); \draw[fig] (1,1.5) .. controls (1.5,1.5) and (1.5,2.4) .. (1.8,2.4) .. controls (2.1,2.4) and (2.2,0.5) .. (2.6,0.5);
	\draw[mass, -{Stealth[length=3pt]}] (3.1,1.5) -- (3.1,0.55); \draw[thin, hilf] (1,1.5) -- (3.3,1.5);}}
\newcommand{\skEndotherm}{\skizze{\achsen{-0.3}{3.8}{-0.3}{2.7}{}{E}\draw[fig, thick] (0.2,0.5) -- (1,0.5) (2.6,1.5) -- (3.4,1.5); \draw[fig] (1,0.5) .. controls (1.5,0.5) and (1.5,2.4) .. (1.8,2.4) .. controls (2.1,2.4) and (2.2,1.5) .. (2.6,1.5);
	\draw[mass, -{Stealth[length=3pt]}] (3.1,0.5) -- (3.1,1.45); \draw[thin, hilf] (1,0.5) -- (3.3,0.5);}}
\newcommand{\skAktivierung}{\skizze{\achsen{-0.3}{3.8}{-0.3}{2.7}{}{E}\draw[fig, thick] (0.2,1.5) -- (1,1.5) (2.6,0.5) -- (3.4,0.5); \draw[fig] (1,1.5) .. controls (1.5,1.5) and (1.5,2.4) .. (1.8,2.4) .. controls (2.1,2.4) and (2.2,0.5) .. (2.6,0.5);
	\draw[mass, <->] (0.6,1.5) -- node[left, inner sep=1pt] {$E_A$} (0.6,2.4); \draw[thin, hilf] (0.5,2.4) -- (1.8,2.4);}}
\newcommand{\skKatalysator}{\skizze{\achsen{-0.3}{3.8}{-0.3}{2.7}{}{E}\draw[fig, thick] (0.2,1.5) -- (1,1.5) (2.6,0.5) -- (3.4,0.5); \draw[fig] (1,1.5) .. controls (1.5,1.5) and (1.5,2.4) .. (1.8,2.4) .. controls (2.1,2.4) and (2.2,0.5) .. (2.6,0.5);
	\draw[mass, hilf] (1,1.5) .. controls (1.5,1.5) and (1.5,1.9) .. (1.8,1.9) .. controls (2.1,1.9) and (2.2,0.5) .. (2.6,0.5);}}
\newcommand{\skGeschwindigkeit}{\skizze{\achsen{-0.3}{3.8}{-0.3}{2.6}{t}{c}\draw[fig, domain=0:3.5, samples=30] plot (\x,{2.2*exp(-1.1*\x)}); \draw[mass, domain=0:3.5, samples=30] plot (\x,{2.0*(1-exp(-1.1*\x))});}}
\newcommand{\skGleichgewicht}{\skizze{\achsen{-0.3}{3.8}{-0.3}{2.6}{t}{c}\draw[fig, domain=0:3.5, samples=30] plot (\x,{1.3+0.9*exp(-1.6*\x)}); \draw[mass, domain=0:3.5, samples=30] plot (\x,{0.8*(1-exp(-1.6*\x))});}}
\newcommand{\skNiederschlag}{\skizze{\fill[blue!12] (-0.25,1.4) -- (-0.25,0.25) arc[start angle=180, end angle=360, radius=0.25] -- (0.25,1.4) -- cycle;
	\begin{scope}\clip (-0.25,0) rectangle (0.25,0.45); \fill[black!55] (0,0.25) circle (0.25); \fill[black!55] (-0.25,0.25) rectangle (0.25,0.45);\end{scope}
	\draw[fig] (-0.25,2.4) -- (-0.25,0.25) arc[start angle=180, end angle=360, radius=0.25] -- (0.25,2.4);}}

% Säuren und Basen
\newcommand{\phskala}[1]{\foreach \k/\c in {0/red, 1/red!80!orange, 2/red!50!orange, 3/orange, 4/orange!60!yellow, 5/yellow, 6/yellow!60!green, 7/green!80!black, 8/green!50!blue!70, 9/blue!50!green, 10/blue!70, 11/blue, 12/blue!80!violet, 13/violet!80!blue, 14/violet}
	{\fill[\c] (#1*\k,0) rectangle ++(#1,#1*1.4);}}
\newcommand{\skPH}{\skizze{\phskala{0.24} \node[below] at (0.12,0) {0}; \node[below] at (1.8,0) {7}; \node[below] at (3.48,0) {14};}}
\newcommand{\skTitration}{\skizze{\stativ \draw[black, thick] (-0.8,2.4) -- (0.2,2.4); \buerette{0.2}{1.2}
	\fill[blue!12] (-0.33,0.25) -- (-0.5,0.05) -- (0.9,0.05) -- (0.73,0.25) -- cycle; \draw[fig, thick] (0.05,1.1) -- (0.05,0.75) -- (-0.5,0.05) -- (0.9,0.05) -- (0.35,0.75) -- (0.35,1.1);}}
\newcommand{\skTitrationskurve}{\skizze{\achsen{-0.3}{3.8}{-0.3}{2.7}{V}{}\node[left] at (0,2.4) {pH}; \draw[fig, domain=0:3.5, samples=60] plot (\x,{1.3+0.7*rad(atan(7*(\x-1.9)))});
	\fill[red] (1.9,1.3) circle (2.5pt);}}

% Elektrochemie
\newcommand{\skElektrolyse}{\skizze{\fill[blue!12] (-1.2,0) rectangle (1.2,1.2); \draw[fig] (-1.2,1.9) -- (-1.2,0) -- (1.2,0) -- (1.2,1.9);
	\draw[black, line width=2pt] (-0.6,0.3) -- (-0.6,2.2) (0.6,0.3) -- (0.6,2.2); \draw[thick] (-0.6,2.2) -- (-0.6,2.8) -- (-0.12,2.8) (0.12,2.8) -- (0.6,2.8) -- (0.6,2.2);
	\draw[thick] (-0.12,2.55) -- (-0.12,3.05); \draw[line width=1.6pt] (0.12,2.67) -- (0.12,2.93); \node at (-0.95,2.5) {$+$}; \node at (0.95,2.5) {$-$};}}
\newcommand{\skGalvanisch}{\skizze{\fill[blue!12] (-1.9,0) rectangle (-0.5,1); \fill[blue!12] (0.5,0) rectangle (1.9,1);
	\draw[fig] (-1.9,1.5) -- (-1.9,0) -- (-0.5,0) -- (-0.5,1.5) (0.5,1.5) -- (0.5,0) -- (1.9,0) -- (1.9,1.5);
	\draw[black!60, line width=2.2pt] (-1.4,0.25) -- (-1.4,1.9); \draw[orange!80!black, line width=2.2pt] (1.4,0.25) -- (1.4,1.9);
	\draw[thick] (-1.4,1.9) -- (-1.4,2.5) -- (-0.3,2.5) (0.3,2.5) -- (1.4,2.5) -- (1.4,1.9); \draw[thick] (0,2.5) circle (0.3); \node at (0,2.5) {\tiny V};
	\draw[fig, thick] (-0.8,0.4) -- (-0.8,1.7) -- (0.8,1.7) -- (0.8,0.4);}}
\newcommand{\skRost}{\skizze{\filldraw[fill=black!25, draw=black!60, thick] (0,0) rectangle (2.6,0.5); \fill[orange!60!brown] (0.3,0.5) .. controls (0.5,0.8) and (0.9,0.75) .. (1.1,0.5) -- cycle;
	\fill[orange!60!brown] (1.5,0.5) .. controls (1.7,0.7) and (2.1,0.7) .. (2.3,0.5) -- cycle;}}
\newcommand{\skKolonne}{\skizze{\draw[fig, rounded corners=4pt] (0,0) rectangle (1,3); \foreach \y in {0.6,1.2,1.8,2.4} {\draw[blue, thin] (0,\y) -- (1,\y); \draw[thick, -{Stealth[length=3pt]}] (1,\y+0.3) -- (1.9,\y+0.3);}
	\draw[thick, -{Stealth[length=3pt]}] (-0.9,0.3) -- (0,0.3);}}
\newcommand{\skKette}{\skizze{\foreach \k in {0,...,5} {\fill[blue] (0.55*\k,{0.18*mod(\k,2)}) circle (4pt);} \draw[blue, thick] (0,0) -- (0.55,0.18) -- (1.1,0) -- (1.65,0.18) -- (2.2,0) -- (2.75,0.18);}}


\title{Chemisches Wörterbuch\\[0.3em] \large Deutsch -- Englisch mit englischer Lautschrift, Formeln und Skizzen}
\author{Martin Quandt}
\date{\today}
\pagestyle{plain}
\setlength{\tabcolsep}{4pt}


% Zweisprachige Überschriften: \teil, \kapitel, \zwischen{Deutsch}{English}{Lautschrift}
% Wie \eintrag wird " ausgeschaltet, BEVOR die Argumente gelesen werden (Betonungszeichen der Lautschrift).
% Ins Inhaltsverzeichnis und in die PDF-Lesezeichen kommt nur „Deutsch -- English“ ohne Lautschrift.
\makeatletter
\newcommand\teil{\begingroup\shorthandoff{"}\teil@aux}
\newcommand\teil@aux[3]{\endgroup\part[#1 -- #2]{\raggedright #1\\[5pt]{\normalfont\Large\color{black!55}#2\hspace{0.35em}\penalty0 \ipa{#3}}}}
\newcommand\kapitel{\begingroup\shorthandoff{"}\kapitel@aux}
\newcommand\kapitel@aux[3]{\endgroup\section[#1 -- #2]{\raggedright #1\\[0pt]{\normalfont\small\color{black!55}#2\hspace{0.35em}\penalty0 \ipa{#3}}}}
\newcommand\zwischen{\begingroup\shorthandoff{"}\zwischen@aux}
\newcommand\zwischen@aux[3]{\endgroup\subsection*{\raggedright #1 {\normalfont\small\color{black!55}\,--\, #2\hspace{0.35em}\penalty0 \ipa{#3}}}}
% Die zweite Überschriftenzeile kostet Platz: Zeilen in den Listen und Tabellen etwas enger setzen
\renewcommand\lies@aux[3]{\endgroup #1 & #2\newline{\color{black!60}#3} \\[1.5pt]}
\makeatother
\setlength{\aboverulesep}{0.15ex}\setlength{\belowrulesep}{0.3ex}\renewcommand{\arraystretch}{0.93}

\begin{document}

\maketitle

Dieses Wörterbuch enthält 260 chemische Fachbegriffe aus der Sekundarstufe I und II,
geordnet nach Sachgebieten. Zu jedem Begriff stehen der englische Fachausdruck mit seiner Aussprache und
-- wo es hilft -- eine Formel, eine Reaktionsgleichung oder eine kleine Skizze. Teil~II zeigt, wie man
Elementnamen, Formeln und Reaktionsgleichungen auf Englisch liest, wie Verbindungen benannt werden und
wie man Experimente beschreibt. Am Ende finden Sie je ein alphabetisches Register Deutsch -- Englisch
und Englisch -- Deutsch.

\zwischen{Hinweise zur Lautschrift}{Notes on the phonetic transcription}{"n@Uts 6n D@ f@"netIk tr\ae n"skrIpS\textsyllabic{n}}

Die Lautschrift folgt dem Internationalen Phonetischen Alphabet (IPA) und gibt die britische
Aussprache wieder, wie sie in Schulwörterbüchern steht (dort wird \textipa{r} statt
\textipa{\textturnr} geschrieben). Der Strich \textipa{\textprimstress} steht \textit{vor} der
betonten Silbe, \textipa{\textsecstress} vor einer Silbe mit Nebenbetonung, \textipa{:} bedeutet lang.

\medskip
{\small
\begin{tabular}{>{\centering\arraybackslash}p{0.9cm} l @{\hspace{1.2cm}} >{\centering\arraybackslash}p{0.9cm} l}
	\toprule
	\textbf{Laut} & \textbf{Beispiel} & \textbf{Laut} & \textbf{Beispiel} \\ \midrule
	\textipa{@} & wat\textbf{er} \ipa{"wO:t@} & \textipa{T} & me\textbf{th}ane \ipa{"mi:TeIn} \\
	\textipa{\ae} & \textbf{a}cid \ipa{"\ae sId} & \textipa{S} & solu\textbf{t}ion \ipa{s@"lu:S\textsyllabic{n}} \\
	\textipa{2} & s\textbf{u}lfur \ipa{"s2lf@} & \textipa{dZ} & hydro\textbf{g}en \ipa{"haIdr@dZ@n} \\
	\textipa{6} & s\textbf{o}lid \ipa{"s6lId} & \textipa{aI} & ox\textbf{i}de \ipa{"6ksaId} \\
	\textipa{3:} & b\textbf{ur}ner \ipa{"b3:n@} & \textipa{eI} & sulf\textbf{a}te \ipa{"s2lfeIt} \\
	\textipa{@U} & s\textbf{o}dium \ipa{"s@Udi@m} & \textipa{\textsyllabic{l}} & meta\textbf{l} \ipa{"met\textsyllabic{l}} \\
	\bottomrule
\end{tabular}}

\medskip
\textbf{Schreibweisen:} Die englischen Namen folgen den IUPAC-Empfehlungen (\textit{sulfur},
\textit{ethanoic acid}); verbreitete ältere Namen stehen in Klammern. \textbf{Farben in den Skizzen:}
Geräte sind blau, hervorgehobene Teile (der Begriff selbst) rot, Hilfslinien gestrichelt.

\newpage
{\footnotesize\setlength{\parskip}{0pt}\tableofcontents}
\newpage

\teil{Wörterbuch}{Dictionary}{"dIkS@nri}

%%%%  STOFFE UND IHRE EIGENSCHAFTEN  %%%%
\kapitel{Stoffe und ihre Eigenschaften}{Substances and their properties}{"s2bst@nsIz @n De@ "pr6p@tiz}

\begin{woerter}
\eintrag{1}{der}{Stoff}{substance}{"s2bst@ns}{}
\eintrag{2}{der}{Reinstoff}{pure substance}{"pjU@ "s2bst@ns}{\skReinstoff}
\eintrag{3}{das}{Element}{element}{"elIm@nt}{\skElement}
\eintrag{4}{die}{Verbindung}{compound}{"k6mpaUnd}{\skVerbindung}
\eintrag{5}{das}{Stoffgemisch}{mixture}{"mIkstS@}{\skGemisch}
\eintrag{6}{}{homogen}{homogeneous}{""h6m@"dZi:ni@s}{}
\eintrag{7}{}{heterogen}{heterogeneous}{""het@r@"dZi:ni@s}{}
\eintrag{8}{die}{Lösung}{solution}{s@"lu:S\textsyllabic{n}}{\skLoesung}
\eintrag{9}{das}{Lösungsmittel}{solvent}{"s6lv@nt}{}
\eintrag{10}{der}{gelöste Stoff}{solute}{"s6lju:t}{}
\eintrag{11}{}{löslich}{soluble}{"s6lj@b\textsyllabic{l}}{}
\eintrag{12}{}{unlöslich}{insoluble}{In"s6lj@b\textsyllabic{l}}{}
\eintrag{13}{die}{Löslichkeit}{solubility}{""s6lju"bIl@ti}{}
\eintrag{14}{die}{gesättigte Lösung}{saturated solution}{"s\ae tS@reItId s@"lu:S\textsyllabic{n}}{}
\eintrag{15}{die}{Suspension}{suspension}{s@"spenS\textsyllabic{n}}{\skSuspension}
\eintrag{16}{die}{Emulsion}{emulsion}{I"m2lS\textsyllabic{n}}{}
\eintrag{17}{die}{Legierung}{alloy}{"\ae lOI}{}
\eintrag{18}{das}{Metall}{metal}{"met\textsyllabic{l}}{}
\eintrag{19}{das}{Nichtmetall}{non-metal}{""n6n"met\textsyllabic{l}}{}
\eintrag{20}{die}{Eigenschaft}{property}{"pr6p@ti}{}
\eintrag{21}{die}{Farbe}{colour (AE: color)}{"k2l@}{}
\eintrag{22}{der}{Geruch}{smell, odour}{smel, "@Ud@}{}
\eintrag{23}{die}{Härte}{hardness}{"hA:dn@s}{}
\eintrag{24}{die}{Dichte}{density}{"dens@ti}{$\varrho = \dfrac{m}{V}$}
\eintrag{25}{die}{Schmelztemperatur}{melting point}{"meltIN pOInt}{}
\eintrag{26}{die}{Siedetemperatur}{boiling point}{"bOIlIN pOInt}{}
\eintrag{27}{die}{elektrische Leitfähigkeit}{electrical conductivity}{I"lektrIk\textsyllabic{l} ""k6nd2k"tIv@ti}{}
\eintrag{28}{}{brennbar}{flammable, combustible}{"fl\ae m@b\textsyllabic{l}, k@m"b2st@b\textsyllabic{l}}{}
\eintrag{29}{}{fest, flüssig, gasförmig}{solid, liquid, gaseous}{"s6lId, "lIkwId, "g\ae si@s}{(s), (l), (g)}
\eintrag{30}{das}{Trennverfahren}{separation technique}{""sep@"reIS\textsyllabic{n} tek""ni:k}{}
\eintrag{31}{das}{Filtrieren}{filtration}{fIl"treIS\textsyllabic{n}}{\skFiltration}
\eintrag{32}{das}{Filtrat}{filtrate}{"fIltreIt}{}
\eintrag{33}{der}{Rückstand}{residue}{"rezIdju:}{}
\eintrag{34}{das}{Eindampfen}{evaporation}{I""v\ae p@"reIS\textsyllabic{n}}{}
\eintrag{35}{die}{Destillation}{distillation}{""dIstI"leIS\textsyllabic{n}}{}
\eintrag{36}{das}{Dekantieren}{decanting}{dI"k\ae ntIN}{}
\eintrag{37}{die}{Kristallisation}{crystallisation}{""krIst@laI"zeIS\textsyllabic{n}}{}
\eintrag{38}{die}{Chromatografie}{chromatography}{""kr@Um@"t6gr@fi}{}
\end{woerter}

\platz\zwischen{Beispiele}{Examples}{Ig"zA:mp\textsyllabic{l}z}

\begin{lesen}
\lies{\skElement \; \skVerbindung}{An element consists of only one kind of atom; a compound consists of two or more elements chemically combined.}{Ein Element besteht aus nur einer Atomsorte; eine Verbindung besteht aus zwei oder mehr chemisch verbundenen Elementen.}
\lies{\skLoesung}{Salt dissolves in water: salt is the solute, water is the solvent, and salt water is the solution.}{Salz löst sich in Wasser: Salz ist der gelöste Stoff, Wasser das Lösungsmittel und Salzwasser die Lösung.}
\lies{\skFiltration}{Sand can be separated from water by filtration: the sand stays in the filter paper as the residue, the water passes through as the filtrate.}{Sand lässt sich von Wasser durch Filtrieren trennen: Der Sand bleibt als Rückstand im Filter, das Wasser läuft als Filtrat durch.}
\lies{Eigenschaften}{Copper is a reddish-brown metal. It is a good conductor of electricity and heat and it melts at 1085 degrees Celsius.}{Kupfer ist ein rotbraunes Metall. Es leitet den elektrischen Strom und Wärme gut und schmilzt bei 1085 Grad Celsius.}
\lies{Destillation}{Ethanol and water can be separated by distillation because they have different boiling points.}{Ethanol und Wasser lassen sich durch Destillation trennen, weil sie verschiedene Siedetemperaturen haben.}
\end{lesen}


%%%%  LABORGERÄTE UND ARBEITSWEISEN  %%%%
\newpage
\kapitel{Laborgeräte und Arbeitsweisen}{Laboratory apparatus and techniques}{l@"b6r@tri ""\ae p@"reIt@s @n tek"ni:ks}

\begin{woerter}
\eintrag{39}{das}{Reagenzglas}{test tube}{"test tju:b}{\skReagenzglas}
\eintrag{40}{das}{Becherglas}{beaker}{"bi:k@}{\skBecherglas}
\eintrag{41}{der}{Erlenmeyerkolben}{conical flask}{"k6nIk\textsyllabic{l} flA:sk}{\skErlenmeyer}
\eintrag{42}{der}{Rundkolben}{round-bottomed flask}{""raUnd"b6t@md flA:sk}{\skRundkolben}
\eintrag{43}{der}{Messzylinder}{measuring cylinder}{"meZ@rIN ""sIlInd@}{\skMesszylinder}
\eintrag{44}{der}{Trichter}{funnel}{"f2n\textsyllabic{l}}{\skTrichter}
\eintrag{45}{das}{Filterpapier}{filter paper}{"fIlt@ ""peIp@}{}
\eintrag{46}{der}{Gasbrenner}{Bunsen burner}{""b2ns\textsyllabic{n} "b3:n@}{\skBrenner}
\eintrag{47}{der}{Dreifuß (mit Drahtnetz)}{tripod (and gauze)}{"traIp6d, gO:z}{\skDreifuss}
\eintrag{48}{die}{Abdampfschale}{evaporating dish}{I"v\ae p@reItIN dIS}{\skAbdampfschale}
\eintrag{49}{der}{Tiegel}{crucible}{"kru:sIb\textsyllabic{l}}{\skTiegel}
\eintrag{50}{die}{Pipette}{pipette}{pI"pet}{\skPipette}
\eintrag{51}{die}{Bürette}{burette}{bju"ret}{\skBuerette}
\eintrag{52}{der}{Spatel}{spatula}{"sp\ae tS@l@}{\skSpatel}
\eintrag{53}{der}{Glasstab}{glass rod, stirring rod}{"glA:s r6d}{}
\eintrag{54}{der}{Mörser mit Pistill}{mortar and pestle}{"mO:t@r @n "pes\textsyllabic{l}}{\skMoerser}
\eintrag{55}{das}{Uhrglas}{watch glass}{"w6tS glA:s}{\skUhrglas}
\eintrag{56}{das}{Stativ}{clamp stand}{"kl\ae mp st\ae nd}{\skStativ}
\eintrag{57}{das}{Thermometer}{thermometer}{T@"m6mIt@}{\skThermometer}
\eintrag{58}{der}{Kühler}{condenser}{k@n"dens@}{\skKuehler}
\eintrag{59}{die}{Waage}{balance}{"b\ae l@ns}{\skWaage}
\eintrag{60}{der}{Stopfen}{stopper, bung}{"st6p@, b2N}{}
\eintrag{61}{die}{Tiegelzange}{tongs}{t6Nz}{}
\eintrag{62}{die}{Schutzbrille}{safety goggles}{"seIfti ""g6g\textsyllabic{l}z}{\skSchutzbrille}
\eintrag{63}{der}{Abzug}{fume cupboard (AE: fume hood)}{"fju:m ""k2b@d}{}
\eintrag{64}{die}{Chemikalie}{chemical}{"kemIk\textsyllabic{l}}{}
\eintrag{65}{}{erhitzen}{heat}{hi:t}{\skErhitzen}
\eintrag{66}{}{rühren}{stir}{st3:}{}
\eintrag{67}{}{lösen}{dissolve}{dI"z6lv}{}
\eintrag{68}{}{gießen}{pour}{pO:}{}
\eintrag{69}{}{abwiegen}{weigh out}{""weI "aUt}{}
\eintrag{70}{}{verdünnen}{dilute}{daI"lu:t}{}
\eintrag{71}{}{abkühlen}{cool (down)}{ku:l}{}
\end{woerter}

\platz\zwischen{Beispiele}{Examples}{Ig"zA:mp\textsyllabic{l}z}

\begin{lesen}
\lies{\skMesszylinder}{Measure out fifty millilitres of water using a measuring cylinder.}{Miss mit dem Messzylinder fünfzig Milliliter Wasser ab.}
\lies{\skErhitzen}{Heat the test tube gently in the Bunsen flame. Do not point the test tube at anyone.}{Erhitze das Reagenzglas vorsichtig in der Brennerflamme. Richte die Öffnung des Reagenzglases nicht auf Personen.}
\lies{\skBecherglas}{Pour the solution into a beaker and stir it with a glass rod.}{Gieße die Lösung in ein Becherglas und rühre mit einem Glasstab um.}
\lies{\skWaage}{Weigh out twelve point five grams of copper sulfate on the balance.}{Wiege auf der Waage zwölf Komma fünf Gramm Kupfersulfat ab.}
\lies{\skSchutzbrille}{Always wear safety goggles when you are working with chemicals.}{Trage beim Arbeiten mit Chemikalien immer eine Schutzbrille.}
\end{lesen}


%%%%  ATOMBAU UND PERIODENSYSTEM  %%%%
\newpage
\kapitel{Atombau und Periodensystem}{Atomic structure and the periodic table}{@"t6mIk "str2ktS@r @n D@ ""pI@ri"6dIk "teIb\textsyllabic{l}}

\begin{woerter}
\eintrag{72}{das}{Atom}{atom}{"\ae t@m}{}
\eintrag{73}{der}{Atomkern}{nucleus}{"nju:kli@s}{}
\eintrag{74}{das}{Proton}{proton}{"pr@Ut6n}{$\mathrm{p^+}$}
\eintrag{75}{das}{Neutron}{neutron}{"nju:tr6n}{$\mathrm{n}$}
\eintrag{76}{das}{Elektron}{electron}{I"lektr6n}{$\mathrm{e^-}$}
\eintrag{77}{die}{Elektronenschale}{electron shell}{I"lektr6n Sel}{\skSchalen}
\eintrag{78}{das}{Außenelektron}{outer electron, valence electron}{"aUt@r I""lektr6n, "veIl@ns}{\skAussenelektron}
\eintrag{79}{die}{Elektronenkonfiguration}{electron configuration}{I"lektr6n k@n""fIgj@"reIS\textsyllabic{n}}{Na: 2, 8, 1}
\eintrag{80}{die}{Ordnungszahl}{atomic number, proton number}{@"t6mIk "n2mb@}{$Z$}
\eintrag{81}{die}{Massenzahl}{mass number}{"m\ae s ""n2mb@}{$A$}
\eintrag{82}{die}{relative Atommasse}{relative atomic mass}{"rel@tIv @"t6mIk m\ae s}{$A_r$}
\eintrag{83}{das}{Isotop}{isotope}{"aIs@t@Up}{\ce{^{12}_{6}C}, \ce{^{14}_{6}C}}
\eintrag{84}{das}{Elementsymbol}{chemical symbol}{"kemIk\textsyllabic{l} "sImb\textsyllabic{l}}{\ce{H}, \ce{O}, \ce{Na}}
\eintrag{85}{das}{Periodensystem der Elemente}{periodic table}{""pI@ri"6dIk "teIb\textsyllabic{l}}{\skPSE}
\eintrag{86}{die}{Periode}{period}{"pI@ri@d}{\pse{(0,0.6) rectangle (2.4,0.9)}}
\eintrag{87}{die}{Hauptgruppe}{(main) group}{gru:p}{\pse{(0.9,0) rectangle (1.2,1.2)}}
\eintrag{88}{die}{Alkalimetalle}{alkali metals}{"\ae lk@laI ""met\textsyllabic{l}z}{\pse{(0,0) rectangle (0.3,1.2)}}
\eintrag{89}{die}{Erdalkalimetalle}{alkaline earth metals}{"\ae lk@laIn "3:T ""met\textsyllabic{l}z}{\pse{(0.3,0) rectangle (0.6,1.2)}}
\eintrag{90}{die}{Halogene}{halogens}{"h\ae l@dZ@nz}{\pse{(1.8,0) rectangle (2.1,1.2)}}
\eintrag{91}{die}{Edelgase}{noble gases}{"n@Ub\textsyllabic{l} ""g\ae sIz}{\pse{(2.1,0) rectangle (2.4,1.5)}}
\eintrag{92}{die}{Übergangsmetalle}{transition metals}{tr\ae n"zIS\textsyllabic{n} ""met\textsyllabic{l}z}{\ce{Fe}, \ce{Cu}, \ce{Zn}}
\eintrag{93}{das}{Ion}{ion}{"aI@n}{}
\eintrag{94}{das}{Kation}{cation}{"k\ae taI@n}{\ce{Na+}}
\eintrag{95}{das}{Anion}{anion}{"\ae naI@n}{\ce{Cl-}}
\eintrag{96}{die}{Edelgaskonfiguration}{noble gas configuration}{"n@Ub\textsyllabic{l} "g\ae s k@n""fIgj@"reIS\textsyllabic{n}}{}
\eintrag{97}{die}{Oktettregel}{octet rule}{6k"tet ru:l}{}
\eintrag{98}{die}{Elektronegativität}{electronegativity}{I""lektr@U""neg@"tIv@ti}{EN}
\eintrag{99}{die}{Ionisierungsenergie}{ionisation energy}{""aI@naI"zeIS\textsyllabic{n} ""en@dZi}{}
\eintrag{100}{das}{Molekül}{molecule}{"m6lIkju:l}{\ce{H2O}}
\eintrag{101}{die}{chemische Formel}{chemical formula}{"kemIk\textsyllabic{l} "fO:mj@l@}{}
\eintrag{102}{die}{Summenformel}{molecular formula}{m@"lekj@l@ "fO:mj@l@}{\ce{C2H6O}}
\eintrag{103}{die}{Strukturformel}{structural formula}{"str2ktS@r@l "fO:mj@l@}{\struktur{H-C(-[2]H)(-[6]H)-C(-[2]H)(-[6]H)-O-H}}
\eintrag{104}{die}{Verhältnisformel}{empirical formula}{Im"pIrIk\textsyllabic{l} "fO:mj@l@}{\ce{NaCl}}
\end{woerter}

\platz\zwischen{Beispiele}{Examples}{Ig"zA:mp\textsyllabic{l}z}

\begin{lesen}
\lies{\skSchalen}{A sodium atom has eleven protons in its nucleus and eleven electrons in three shells: two, eight, one.}{Ein Natriumatom hat elf Protonen im Kern und elf Elektronen auf drei Schalen: zwei, acht, eins.}
\lies{\ce{Na -> Na+ + e-}}{A sodium atom loses its outer electron and forms a positively charged sodium ion.}{Ein Natriumatom gibt sein Außenelektron ab und bildet ein positiv geladenes Natrium-Ion.}
\lies{\pse{(0,0) rectangle (0.3,1.2)}}{The alkali metals are in group one of the periodic table. Their atoms have one outer electron.}{Die Alkalimetalle stehen in der ersten Hauptgruppe des Periodensystems. Ihre Atome haben ein Außenelektron.}
\lies{\pse{(2.1,0) rectangle (2.4,1.5)}}{The noble gases have full outer shells and are very unreactive.}{Die Edelgase haben voll besetzte Außenschalen und sind sehr reaktionsträge.}
\lies{\ce{^{35}_{17}Cl}, \ce{^{37}_{17}Cl}}{Chlorine has two isotopes: chlorine-thirty-five and chlorine-thirty-seven. They have the same number of protons but different numbers of neutrons.}{Chlor hat zwei Isotope: Chlor-35 und Chlor-37. Sie haben gleich viele Protonen, aber verschieden viele Neutronen.}
\end{lesen}


%%%%  CHEMISCHE BINDUNG  %%%%
\newpage
\kapitel{Chemische Bindung}{Chemical bonding}{"kemIk\textsyllabic{l} "b6ndIN}

\begin{woerter}
\eintrag{105}{die}{chemische Bindung}{chemical bond}{"kemIk\textsyllabic{l} b6nd}{}
\eintrag{106}{die}{Ionenbindung}{ionic bond}{aI"6nIk b6nd}{\ce{Na+} \ce{Cl-}}
\eintrag{107}{das}{Ionengitter}{ionic lattice}{aI"6nIk "l\ae tIs}{\skGitter}
\eintrag{108}{der}{Kristall}{crystal}{"krIst\textsyllabic{l}}{}
\eintrag{109}{die}{Atombindung (Elektronenpaarbindung)}{covalent bond}{""k@U"veIl@nt b6nd}{\skAtombindung}
\eintrag{110}{das}{bindende Elektronenpaar}{bonding pair}{"b6ndIN pe@}{$\mathrm{H{-}H}$}
\eintrag{111}{das}{freie Elektronenpaar}{lone pair}{"l@Un pe@}{$\mathrm{H{-}\overline{\underline{O}}{-}H}$}
\eintrag{112}{die}{Einfachbindung}{single bond}{"sINg\textsyllabic{l} b6nd}{$\mathrm{H{-}H}$}
\eintrag{113}{die}{Doppelbindung}{double bond}{"d2b\textsyllabic{l} b6nd}{$\mathrm{O{=}O}$}
\eintrag{114}{die}{Dreifachbindung}{triple bond}{"trIp\textsyllabic{l} b6nd}{$\mathrm{N{\equiv}N}$}
\eintrag{115}{die}{Metallbindung}{metallic bond}{m@"t\ae lIk b6nd}{\skMetall}
\eintrag{116}{die}{polare Atombindung}{polar covalent bond}{"p@Ul@ ""k@U"veIl@nt b6nd}{$\overset{\delta+}{\mathrm{H}}{-}\overset{\delta-}{\mathrm{Cl}}$}
\eintrag{117}{der}{Dipol}{dipole}{"daIp@Ul}{\skDipol}
\eintrag{118}{die}{Wasserstoffbrückenbindung}{hydrogen bond}{"haIdr@dZ@n b6nd}{}
\eintrag{119}{die}{zwischenmolekularen Kräfte}{intermolecular forces}{""Int@m@"lekj@l@ "fO:sIz}{}
\eintrag{120}{die}{Van-der-Waals-Kräfte}{van der Waals forces}{""v\ae n d@ "wA:lz ""fO:sIz}{}
\eintrag{121}{die}{Wertigkeit}{valency}{"veIl@nsi}{Fe(III)}
\eintrag{122}{die}{Oxidationszahl}{oxidation number}{""6ksI"deIS\textsyllabic{n} ""n2mb@}{$\oz{+I}{H}_2\oz{-II}{O}$}
\eintrag{123}{die}{Lewis-Formel}{Lewis structure, dot-and-cross diagram}{"lu:Is ""str2ktS@}{}
\eintrag{124}{die}{delokalisierten Elektronen}{delocalised electrons}{di:"l@Uk@laIzd I"lektr6nz}{}
\end{woerter}

\platz\zwischen{Beispiele}{Examples}{Ig"zA:mp\textsyllabic{l}z}

\begin{lesen}
\lies{\skGitter}{In sodium chloride the sodium ions and the chloride ions are held together in a lattice by ionic bonds.}{Im Natriumchlorid werden die Natrium-Ionen und die Chlorid-Ionen durch Ionenbindung in einem Gitter zusammengehalten.}
\lies{\skAtombindung}{In a hydrogen molecule the two atoms share a pair of electrons: a covalent bond.}{Im Wasserstoffmolekül teilen sich die beiden Atome ein Elektronenpaar: eine Atombindung.}
\lies{$\mathrm{O{=}C{=}O}$}{A carbon dioxide molecule contains two double bonds.}{Ein Kohlenstoffdioxidmolekül enthält zwei Doppelbindungen.}
\lies{\skDipol}{The water molecule is a dipole: the oxygen atom carries a partial negative charge, the hydrogen atoms a partial positive charge.}{Das Wassermolekül ist ein Dipol: Das Sauerstoffatom trägt eine negative, die Wasserstoffatome tragen eine positive Teilladung.}
\lies{\skMetall}{Metals conduct electricity because their delocalised electrons are free to move.}{Metalle leiten den elektrischen Strom, weil ihre delokalisierten Elektronen frei beweglich sind.}
\end{lesen}


%%%%  CHEMISCHE REAKTIONEN  %%%%
\newpage
\kapitel{Chemische Reaktionen}{Chemical reactions}{"kemIk\textsyllabic{l} ri"\ae kS\textsyllabic{n}z}

\begin{woerter}
\eintrag{125}{die}{chemische Reaktion}{chemical reaction}{"kemIk\textsyllabic{l} ri"\ae kS\textsyllabic{n}}{}
\eintrag{126}{der}{Ausgangsstoff (das Edukt)}{reactant}{ri"\ae kt@nt}{{\color{red}\ce{2 H2 + O2}} \ce{-> 2 H2O}}
\eintrag{127}{das}{Reaktionsprodukt}{product}{"pr6d2kt}{\ce{2 H2 + O2 ->} {\color{red}\ce{2 H2O}}}
\eintrag{128}{die}{Wortgleichung}{word equation}{"w3:d I""kweIZ\textsyllabic{n}}{}
\eintrag{129}{die}{Reaktionsgleichung}{chemical equation}{"kemIk\textsyllabic{l} I"kweIZ\textsyllabic{n}}{\ce{2 Mg + O2 -> 2 MgO}}
\eintrag{130}{}{ausgleichen}{balance (an equation)}{"b\ae l@ns}{}
\eintrag{131}{der}{Koeffizient (Stöchiometriezahl)}{coefficient}{""k@UI"fIS\textsyllabic{n}t}{{\color{red}2}\,\ce{Mg + O2 ->} {\color{red}2}\,\ce{MgO}}
\eintrag{132}{der}{Index}{subscript}{"s2bskrIpt}{$\mathrm{H_{\color{red}2}O}$}
\eintrag{133}{die}{Synthese}{synthesis}{"sInT@sIs}{\ce{A + B -> AB}}
\eintrag{134}{die}{Zersetzung (Analyse)}{decomposition}{""di:k6mp@"zIS\textsyllabic{n}}{\ce{AB -> A + B}}
\eintrag{135}{die}{Verbrennung}{combustion}{k@m"b2stS@n}{\ce{C + O2 -> CO2}}
\eintrag{136}{die}{Fällung}{precipitation}{prI""sIpI"teIS\textsyllabic{n}}{\ce{Ag+ + Cl- -> AgCl v}}
\eintrag{137}{der}{Niederschlag}{precipitate}{prI"sIpIteIt}{\skNiederschlag}
\eintrag{138}{}{exotherm}{exothermic}{""eks@U"T3:mIk}{$\Delta H < 0$ \newline \skExotherm}
\eintrag{139}{}{endotherm}{endothermic}{""end@U"T3:mIk}{$\Delta H > 0$ \newline \skEndotherm}
\eintrag{140}{die}{Reaktionsenthalpie}{enthalpy change of reaction}{"enT@lpi tSeIndZ}{$\Delta H$ in $\mathrm{kJ/mol}$}
\eintrag{141}{die}{Aktivierungsenergie}{activation energy}{""\ae ktI"veIS\textsyllabic{n} ""en@dZi}{\skAktivierung}
\eintrag{142}{der}{Katalysator}{catalyst}{"k\ae t@lIst}{\skKatalysator}
\eintrag{143}{die}{Reaktionsgeschwindigkeit}{rate of reaction}{"reIt @v ri"\ae kS\textsyllabic{n}}{\skGeschwindigkeit}
\eintrag{144}{die}{umkehrbare Reaktion}{reversible reaction}{rI"v3:s@b\textsyllabic{l} ri"\ae kS\textsyllabic{n}}{\ce{A + B <=> C + D}}
\eintrag{145}{das}{chemische Gleichgewicht}{chemical equilibrium}{"kemIk\textsyllabic{l} ""i:kwI"lIbri@m}{\skGleichgewicht}
\eintrag{146}{die}{Gleichgewichtskonstante}{equilibrium constant}{""i:kwI"lIbri@m ""k6nst@nt}{$K_c$}
\eintrag{147}{das}{Massenwirkungsgesetz}{law of mass action}{"lO: @v "m\ae s ""\ae kS\textsyllabic{n}}{$K_c = \dfrac{c(\mathrm{C}) \cdot c(\mathrm{D})}{c(\mathrm{A}) \cdot c(\mathrm{B})}$}
\eintrag{148}{das}{Prinzip von Le Chatelier}{Le Chatelier's principle}{l@ S\ae "telieIz "prIns@p\textsyllabic{l}}{}
\eintrag{149}{das}{Gesetz von der Erhaltung der Masse}{law of conservation of mass}{""k6ns@"veIS\textsyllabic{n} @v "m\ae s}{}
\eintrag{150}{die}{Ausbeute}{yield}{ji:ld}{}
\eintrag{151}{der}{Nachweis}{test (for)}{test}{}
\end{woerter}

\platz\zwischen{Beispiele}{Examples}{Ig"zA:mp\textsyllabic{l}z}

\begin{lesen}
\lies{\ce{2 Mg + O2 -> 2 MgO}}{Magnesium reacts with oxygen to form magnesium oxide. Magnesium and oxygen are the reactants, magnesium oxide is the product.}{Magnesium reagiert mit Sauerstoff zu Magnesiumoxid. Magnesium und Sauerstoff sind die Ausgangsstoffe, Magnesiumoxid ist das Reaktionsprodukt.}
\lies{\skExotherm}{The reaction is exothermic: energy is given out to the surroundings and the temperature rises.}{Die Reaktion ist exotherm: Es wird Energie an die Umgebung abgegeben und die Temperatur steigt.}
\lies{\skKatalysator}{A catalyst speeds up a reaction by lowering the activation energy. It is not used up in the reaction.}{Ein Katalysator beschleunigt eine Reaktion, indem er die Aktivierungsenergie herabsetzt. Er wird dabei nicht verbraucht.}
\lies{\ce{N2 + 3 H2 <=> 2 NH3}}{The reaction is reversible. At equilibrium the forward and the reverse reaction take place at the same rate.}{Die Reaktion ist umkehrbar. Im Gleichgewicht laufen Hin- und Rückreaktion gleich schnell ab.}
\lies{\skNiederschlag}{When silver nitrate solution is added, a white precipitate of silver chloride forms.}{Bei Zugabe von Silbernitratlösung bildet sich ein weißer Niederschlag von Silberchlorid.}
\end{lesen}


%%%%  SÄUREN, BASEN UND SALZE  %%%%
\newpage
\kapitel{Säuren, Basen und Salze}{Acids, bases and salts}{"\ae sIdz, "beIsIz @n "sO:lts}

\begin{woerter}
\eintrag{152}{die}{Säure}{acid}{"\ae sId}{\ce{HCl}}
\eintrag{153}{die}{Base}{base}{beIs}{\ce{NH3}}
\eintrag{154}{die}{Lauge (alkalische Lösung)}{alkali, alkaline solution}{"\ae lk@laI, "\ae lk@laIn s@"lu:S\textsyllabic{n}}{\ce{NaOH(aq)}}
\eintrag{155}{}{sauer}{acidic}{@"sIdIk}{pH $< 7$}
\eintrag{156}{}{neutral}{neutral}{"nju:tr@l}{pH $= 7$}
\eintrag{157}{}{basisch, alkalisch}{basic, alkaline}{"beIsIk, "\ae lk@laIn}{pH $> 7$}
\eintrag{158}{der}{pH-Wert}{pH (value)}{""pi: "eItS}{\skPH}
\eintrag{159}{der}{Indikator}{indicator}{"IndIkeIt@}{}
\eintrag{160}{das}{Lackmus}{litmus}{"lItm@s}{}
\eintrag{161}{der}{Universalindikator}{universal indicator}{""ju:nI"v3:s\textsyllabic{l} "IndIkeIt@}{}
\eintrag{162}{der}{Protonendonator}{proton donor}{"pr@Ut6n ""d@Un@}{\ce{HCl -> H+ + Cl-}}
\eintrag{163}{der}{Protonenakzeptor}{proton acceptor}{"pr@Ut6n @k""sept@}{\ce{NH3 + H+ -> NH4+}}
\eintrag{164}{das}{Oxonium-Ion}{oxonium ion, hydronium ion}{6k"s@Uni@m ""aI@n}{\ce{H3O+}}
\eintrag{165}{das}{Hydroxid-Ion}{hydroxide ion}{haI"dr6ksaId ""aI@n}{\ce{OH-}}
\eintrag{166}{die}{Neutralisation}{neutralisation}{""nju:tr@laI"zeIS\textsyllabic{n}}{\ce{H3O+ + OH- -> 2 H2O}}
\eintrag{167}{das}{Salz}{salt}{sO:lt}{\ce{NaCl}}
\eintrag{168}{die}{starke Säure}{strong acid}{"str6N "\ae sId}{}
\eintrag{169}{die}{schwache Säure}{weak acid}{"wi:k "\ae sId}{}
\eintrag{170}{}{konzentriert}{concentrated}{"k6ns\textsyllabic{n}treItId}{}
\eintrag{171}{}{verdünnt}{dilute}{daI"lu:t}{}
\eintrag{172}{}{ätzend}{corrosive}{k@"r@UsIv}{}
\eintrag{173}{die}{Titration}{titration}{taI"treIS\textsyllabic{n}}{\skTitration}
\eintrag{174}{die}{Titrationskurve}{titration curve}{taI"treIS\textsyllabic{n} k3:v}{\skTitrationskurve}
\eintrag{175}{der}{Äquivalenzpunkt}{equivalence point}{I"kwIv@l@ns pOInt}{}
\eintrag{176}{die}{Maßlösung}{standard solution}{"st\ae nd@d s@"lu:S\textsyllabic{n}}{}
\eintrag{177}{die}{Pufferlösung}{buffer solution}{"b2f@ s@"lu:S\textsyllabic{n}}{}
\eintrag{178}{die}{Salzsäure}{hydrochloric acid}{""haIdr@"kl6rIk "\ae sId}{\ce{HCl(aq)}}
\eintrag{179}{die}{Schwefelsäure}{sulfuric acid}{s2l"fjU@rIk "\ae sId}{\ce{H2SO4}}
\eintrag{180}{die}{Salpetersäure}{nitric acid}{"naItrIk "\ae sId}{\ce{HNO3}}
\eintrag{181}{die}{Natronlauge}{sodium hydroxide solution}{"s@Udi@m haI"dr6ksaId s@"lu:S\textsyllabic{n}}{\ce{NaOH(aq)}}
\eintrag{182}{das}{Ammoniak}{ammonia}{@"m@Uni@}{\ce{NH3}}
\end{woerter}

\platz\zwischen{Beispiele}{Examples}{Ig"zA:mp\textsyllabic{l}z}

\begin{lesen}
\lies{\ce{HCl + NaOH -> NaCl + H2O}}{Hydrochloric acid reacts with sodium hydroxide solution to form sodium chloride and water: an acid and a base neutralise each other.}{Salzsäure reagiert mit Natronlauge zu Natriumchlorid und Wasser: Säure und Base neutralisieren einander.}
\lies{\skPH}{Acidic solutions have a pH below seven, alkaline solutions have a pH above seven.}{Saure Lösungen haben einen pH-Wert unter sieben, alkalische Lösungen einen pH-Wert über sieben.}
\lies{Lackmus}{Acids turn blue litmus red; alkalis turn red litmus blue.}{Säuren färben blaues Lackmus rot; Laugen färben rotes Lackmus blau.}
\lies{\skTitration}{The acid is titrated with sodium hydroxide solution until the indicator changes colour.}{Die Säure wird mit Natronlauge titriert, bis der Indikator umschlägt.}
\lies{\ce{CH3COOH}}{Ethanoic acid is a weak acid: it is only partially ionised in water.}{Essigsäure ist eine schwache Säure: Sie ist in Wasser nur teilweise dissoziiert.}
\end{lesen}


%%%%  REDOXREAKTIONEN UND ELEKTROCHEMIE  %%%%
\newpage
\kapitel{Redoxreaktionen und Elektrochemie}{Redox reactions and electrochemistry}{"ri:d6ks ri"\ae kS\textsyllabic{n}z @n I""lektr@U"kemIstri}

\begin{woerter}
\eintrag{183}{die}{Oxidation}{oxidation}{""6ksI"deIS\textsyllabic{n}}{\ce{Mg -> Mg^2+ + 2 e-}}
\eintrag{184}{die}{Reduktion}{reduction}{rI"d2kS\textsyllabic{n}}{\ce{Cu^2+ + 2 e- -> Cu}}
\eintrag{185}{die}{Redoxreaktion}{redox reaction}{"ri:d6ks ri""\ae kS\textsyllabic{n}}{\ce{Mg + Cu^2+ -> Mg^2+ + Cu}}
\eintrag{186}{das}{Oxidationsmittel}{oxidising agent}{"6ksIdaIzIN ""eIdZ@nt}{}
\eintrag{187}{das}{Reduktionsmittel}{reducing agent}{rI"dju:sIN ""eIdZ@nt}{}
\eintrag{188}{das}{Oxid}{oxide}{"6ksaId}{\ce{MgO}}
\eintrag{189}{der}{Elektronenübergang}{electron transfer}{I"lektr6n ""tr\ae nsf3:}{}
\eintrag{190}{die}{Korrosion}{corrosion}{k@"r@UZ\textsyllabic{n}}{}
\eintrag{191}{der}{Rost}{rust}{r2st}{\skRost}
\eintrag{192}{das}{Edelmetall}{noble metal, precious metal}{"n@Ub\textsyllabic{l} "met\textsyllabic{l}, "preS@s}{\ce{Au}, \ce{Ag}, \ce{Pt}}
\eintrag{193}{das}{unedle Metall}{base metal}{"beIs "met\textsyllabic{l}}{\ce{Zn}, \ce{Fe}, \ce{Mg}}
\eintrag{194}{die}{Redoxreihe der Metalle}{reactivity series}{""ri:\ae k"tIv@ti ""sI@ri:z}{}
\eintrag{195}{die}{Elektrolyse}{electrolysis}{I""lek"tr6l@sIs}{\skElektrolyse}
\eintrag{196}{der}{Elektrolyt}{electrolyte}{I"lektr@laIt}{}
\eintrag{197}{die}{Elektrode}{electrode}{I"lektr@Ud}{}
\eintrag{198}{die}{Anode}{anode}{"\ae n@Ud}{}
\eintrag{199}{die}{Kathode}{cathode}{"k\ae T@Ud}{}
\eintrag{200}{die}{galvanische Zelle}{galvanic cell, electrochemical cell}{g\ae l"v\ae nIk sel}{\skGalvanisch}
\eintrag{201}{die}{Halbzelle}{half-cell}{"hA:f sel}{\ce{Zn}/\ce{Zn^2+}}
\eintrag{202}{das}{Standardpotential}{standard electrode potential}{"st\ae nd@d I"lektr@Ud p@""tenS\textsyllabic{l}}{$E^0$ in V}
\eintrag{203}{die}{Brennstoffzelle}{fuel cell}{"fju:@l sel}{}
\eintrag{204}{der}{Akkumulator}{rechargeable battery}{""ri:"tSA:dZ@b\textsyllabic{l} "b\ae tri}{}
\eintrag{205}{das}{Galvanisieren}{electroplating}{I"lektr@U""pleItIN}{}
\eintrag{206}{das}{Erz}{ore}{O:}{}
\eintrag{207}{der}{Hochofen}{blast furnace}{"blA:st ""f3:nIs}{}
\end{woerter}

\platz\zwischen{Beispiele}{Examples}{Ig"zA:mp\textsyllabic{l}z}

\begin{lesen}
\lies{\ce{2 Cu + O2 -> 2 CuO}}{Copper is oxidised to copper(II) oxide: the copper atoms lose electrons.}{Kupfer wird zu Kupfer(II)-oxid oxidiert: Die Kupferatome geben Elektronen ab.}
\lies{\ce{CuO + H2 -> Cu + H2O}}{Copper(II) oxide is reduced to copper; hydrogen is the reducing agent.}{Kupfer(II)-oxid wird zu Kupfer reduziert; Wasserstoff ist das Reduktionsmittel.}
\lies{Merkhilfe}{OIL RIG: Oxidation Is Loss, Reduction Is Gain (of electrons).}{Oxidation ist Elektronenabgabe, Reduktion ist Elektronenaufnahme.}
\lies{\skElektrolyse}{During electrolysis the positive ions move to the cathode and the negative ions move to the anode.}{Bei der Elektrolyse wandern die positiven Ionen zur Kathode und die negativen Ionen zur Anode.}
\lies{\skRost}{Iron rusts in the presence of water and oxygen.}{Eisen rostet, wenn Wasser und Sauerstoff vorhanden sind.}
\end{lesen}


%%%%  ORGANISCHE CHEMIE  %%%%
\newpage
\kapitel{Organische Chemie}{Organic chemistry}{O:"g\ae nIk "kemIstri}

\begin{woerter}
\eintrag{208}{die}{organische Chemie}{organic chemistry}{O:"g\ae nIk "kemIstri}{}
\eintrag{209}{der}{Kohlenwasserstoff}{hydrocarbon}{""haIdr@"kA:b@n}{}
\eintrag{210}{das}{Alkan}{alkane}{"\ae lkeIn}{\ce{H3C-CH3}}
\eintrag{211}{das}{Alken}{alkene}{"\ae lki:n}{\ce{H2C=CH2}}
\eintrag{212}{das}{Alkin}{alkyne}{"\ae lkaIn}{\ce{HC\bond{3}CH}}
\eintrag{213}{die}{homologe Reihe}{homologous series}{h@"m6l@g@s "sI@ri:z}{$\mathrm{C}_n\mathrm{H}_{2n+2}$}
\eintrag{214}{}{gesättigt}{saturated}{"s\ae tS@reItId}{}
\eintrag{215}{}{ungesättigt}{unsaturated}{2n"s\ae tS@reItId}{}
\eintrag{216}{das}{Isomer}{isomer}{"aIs@m@}{}
\eintrag{217}{der}{Aromat}{aromatic compound}{""\ae r@"m\ae tIk "k6mpaUnd}{\struktur{*6(-=-=-=)}}
\eintrag{218}{die}{funktionelle Gruppe}{functional group}{"f2NkS@n\textsyllabic{l} gru:p}{}
\eintrag{219}{der}{Alkohol}{alcohol}{"\ae lk@h6l}{\ce{R-OH}}
\eintrag{220}{die}{Hydroxygruppe}{hydroxyl group}{haI"dr6ksIl gru:p}{\ce{-OH}}
\eintrag{221}{das}{Aldehyd}{aldehyde}{"\ae ldIhaId}{\ce{R-CHO}}
\eintrag{222}{das}{Keton}{ketone}{"ki:t@Un}{$\mathrm{R{-}CO{-}R'}$}
\eintrag{223}{die}{Carbonsäure}{carboxylic acid}{""kA:b6k"sIlIk "\ae sId}{\ce{R-COOH}}
\eintrag{224}{der}{Ester}{ester}{"est@}{$\mathrm{R{-}COO{-}R'}$}
\eintrag{225}{die}{Veresterung}{esterification}{e""sterIfI"keIS\textsyllabic{n}}{}
\eintrag{226}{das}{Amin}{amine}{@"mi:n}{\ce{R-NH2}}
\eintrag{227}{die}{Aminosäure}{amino acid}{@""mi:n@U "\ae sId}{}
\eintrag{228}{die}{Additionsreaktion}{addition reaction}{@"dIS\textsyllabic{n} ri""\ae kS\textsyllabic{n}}{}
\eintrag{229}{die}{Substitutionsreaktion}{substitution reaction}{""s2bstI"tju:S\textsyllabic{n} ri""\ae kS\textsyllabic{n}}{}
\eintrag{230}{die}{Eliminierung}{elimination}{I""lImI"neIS\textsyllabic{n}}{}
\eintrag{231}{das}{Erdöl}{crude oil}{""kru:d "OIl}{}
\eintrag{232}{das}{Erdgas}{natural gas}{"n\ae tSr@l g\ae s}{}
\eintrag{233}{die}{fraktionierte Destillation}{fractional distillation}{"fr\ae kS@n\textsyllabic{l} ""dIstI"leIS\textsyllabic{n}}{\skKolonne}
\eintrag{234}{das}{Cracken}{cracking}{"kr\ae kIN}{}
\eintrag{235}{das}{Benzin}{petrol (AE: gasoline)}{"petr@l}{}
\eintrag{236}{das}{Monomer}{monomer}{"m6n@m@}{}
\eintrag{237}{das}{Polymer}{polymer}{"p6lIm@}{\skKette}
\eintrag{238}{die}{Polymerisation}{polymerisation}{""p6lIm@raI"zeIS\textsyllabic{n}}{}
\eintrag{239}{der}{Kunststoff}{plastic}{"pl\ae stIk}{}
\eintrag{240}{das}{Fett}{fat}{f\ae t}{}
\eintrag{241}{das}{Kohlenhydrat}{carbohydrate}{""kA:b@U"haIdreIt}{}
\eintrag{242}{der}{Traubenzucker (die Glucose)}{glucose}{"glu:k@Us}{\ce{C6H12O6}}
\eintrag{243}{die}{Stärke}{starch}{stA:tS}{}
\eintrag{244}{das}{Eiweiß (Protein)}{protein}{"pr@Uti:n}{}
\eintrag{245}{die}{Seife}{soap}{s@Up}{}
\eintrag{246}{das}{Tensid}{surfactant, detergent}{s3:"f\ae kt@nt, dI"t3:dZ@nt}{}
\eintrag{247}{die}{alkoholische Gärung}{fermentation}{""f3:men"teIS\textsyllabic{n}}{}
\end{woerter}

\platz\zwischen{Beispiele}{Examples}{Ig"zA:mp\textsyllabic{l}z}

\begin{lesen}
\lies{\ce{CH4 + 2 O2 -> CO2 + 2 H2O}}{Methane burns in oxygen to form carbon dioxide and water.}{Methan verbrennt mit Sauerstoff zu Kohlenstoffdioxid und Wasser.}
\lies{\ce{H2C=CH2}}{Ethene is an unsaturated hydrocarbon: its molecule contains a carbon–carbon double bond. It decolourises bromine water.}{Ethen ist ein ungesättigter Kohlenwasserstoff: Sein Molekül enthält eine C=C-Doppelbindung. Es entfärbt Bromwasser.}
\end{lesen}


%%%%  STOFFMENGE UND CHEMISCHES RECHNEN  %%%%
\newpage
\kapitel{Stoffmenge und chemisches Rechnen}{Amount of substance and chemical calculations}{@"maUnt @v "s2bst@ns @n "kemIk\textsyllabic{l} ""k\ae lkju"leIS\textsyllabic{n}z}

\begin{woerter}
\eintrag{248}{die}{Stoffmenge}{amount of substance}{@"maUnt @v "s2bst@ns}{$n$ in $\mathrm{mol}$}
\eintrag{249}{das}{Mol}{mole}{m@Ul}{$1\eh{mol} \mathrel{\widehat{=}} 6{,}022 \cdot 10^{23}$ Teilchen}
\eintrag{250}{die}{Avogadro-Konstante}{Avogadro constant}{""\ae v@"gA:dr@U ""k6nst@nt}{\small$N_A = 6{,}022 \cdot 10^{23}\eh{mol^{-1}}$}
\eintrag{251}{die}{Teilchenanzahl}{number of particles}{"n2mb@r @v "pA:tIk\textsyllabic{l}z}{$N = n \cdot N_A$}
\eintrag{252}{die}{molare Masse}{molar mass}{"m@Ul@ m\ae s}{$M$ in $\mathrm{g/mol}$ \newline $M = \dfrac{m}{n}$}
\eintrag{253}{das}{molare Volumen}{molar volume}{"m@Ul@ "v6lju:m}{$V_m = 22{,}4\eh{l/mol}$}
\eintrag{254}{die}{relative Molekülmasse}{relative molecular mass}{"rel@tIv m@"lekj@l@ m\ae s}{$M_r$}
\eintrag{255}{die}{Stoffmengenkonzentration}{concentration}{""k6ns\textsyllabic{n}"treIS\textsyllabic{n}}{$c$ in $\mathrm{mol/l}$ \newline $c = \dfrac{n}{V}$}
\eintrag{256}{die}{Massenkonzentration}{mass concentration}{"m\ae s ""k6ns\textsyllabic{n}"treIS\textsyllabic{n}}{$\beta$ in $\mathrm{g/l}$}
\eintrag{257}{der}{Massenanteil}{mass fraction, percentage by mass}{"m\ae s ""fr\ae kS\textsyllabic{n}}{$w = \dfrac{m(\text{Stoff})}{m(\text{Gemisch})}$}
\eintrag{258}{die}{Stöchiometrie}{stoichiometry}{""stOIki"6m@tri}{}
\eintrag{259}{das}{Stoffmengenverhältnis}{mole ratio}{"m@Ul ""reISi@U}{$n(\mathrm{Mg}) : n(\mathrm{O_2}) = 2 : 1$}
\eintrag{260}{der}{Überschuss}{excess}{Ik"ses}{}
\end{woerter}

\platz\zwischen{Beispiele}{Examples}{Ig"zA:mp\textsyllabic{l}z}

\begin{lesen}
\lies{$M(\mathrm{H_2O}) = 18\eh{g/mol}$}{The molar mass of water is eighteen grams per mole.}{Die molare Masse von Wasser beträgt achtzehn Gramm pro Mol.}
\lies{$n = \dfrac{m}{M} = \dfrac{36\eh{g}}{18\eh{g/mol}} = 2\eh{mol}$}{Thirty-six grams of water contain two moles of water molecules.}{Sechsunddreißig Gramm Wasser enthalten zwei Mol Wassermoleküle.}
\lies{$c = \dfrac{n}{V} = \dfrac{0{,}5\eh{mol}}{2\eh{l}} = 0{,}25\eh{mol/l}$}{The concentration of the solution is nought point two five moles per litre.}{Die Konzentration der Lösung beträgt null Komma zwei fünf Mol pro Liter.}
\lies{$V = n \cdot V_m = 2\eh{mol} \cdot 22{,}4\eh{l/mol}$ \newline $V = 44{,}8\eh{l}$}{Under standard conditions two moles of a gas occupy forty-four point eight litres.}{Unter Normbedingungen nehmen zwei Mol eines Gases vierundvierzig Komma acht Liter ein.}
\end{lesen}



%%%%  TEIL II: CHEMIE LESEN UND SPRECHEN  %%%%
\newpage
\teil{Chemie lesen und sprechen}{Reading and speaking chemistry}{"ri:dIN @n "spi:kIN "kemIstri}

In diesem Teil steht, wie man Elementnamen, Formeln und Reaktionsgleichungen auf Englisch (und zum
Vergleich auf Deutsch) liest, wie Verbindungen benannt werden und wie man Experimente beschreibt.
Links steht der Ausdruck, rechts die englische Lesart und darunter grau die deutsche.

%--------------------------------------------------------------------------------------------------
\kapitel{Die Elemente: Namen und Aussprache}{The elements: names and pronunciation}{Di "elIm@nts, "neImz @n pr@""n2nsi"eIS\textsyllabic{n}}

Viele Elemente heißen im Englischen anders als im Deutschen oder werden deutlich anders
ausgesprochen. Das Elementsymbol ist in beiden Sprachen gleich.

\begin{center}
{\small\renewcommand{\arraystretch}{1.2}
\begin{tabular}{@{}l l l @{\hspace{1cm}} l l l@{}}
	\toprule
	& \textbf{Englisch} & \textbf{Deutsch} & & \textbf{Englisch} & \textbf{Deutsch} \\ \midrule
	H & hydrogen \ipa{"haIdr@dZ@n} & Wasserstoff & K & potassium \ipa{p@"t\ae si@m} & Kalium \\
	He & helium \ipa{"hi:li@m} & Helium & Ca & calcium \ipa{"k\ae lsi@m} & Calcium \\
	Li & lithium \ipa{"lITi@m} & Lithium & Cr & chromium \ipa{"kr@Umi@m} & Chrom \\
	C & carbon \ipa{"kA:b@n} & Kohlenstoff & Mn & manganese \ipa{"m\ae Ng@ni:z} & Mangan \\
	N & nitrogen \ipa{"naItr@dZ@n} & Stickstoff & Fe & iron \ipa{"aI@n} & Eisen \\
	O & oxygen \ipa{"6ksIdZ@n} & Sauerstoff & Cu & copper \ipa{"k6p@} & Kupfer \\
	F & fluorine \ipa{"flO:ri:n} & Fluor & Zn & zinc \ipa{zINk} & Zink \\
	Ne & neon \ipa{"ni:6n} & Neon & Br & bromine \ipa{"br@Umi:n} & Brom \\
	Na & sodium \ipa{"s@Udi@m} & Natrium & Ag & silver \ipa{"sIlv@} & Silber \\
	Mg & magnesium \ipa{m\ae g"ni:zi@m} & Magnesium & Sn & tin \ipa{tIn} & Zinn \\
	Al & aluminium \ipa{""\ae lj@"mIni@m} & Aluminium & I & iodine \ipa{"aI@di:n} & Iod \\
	Si & silicon \ipa{"sIlIk@n} & Silicium & Ba & barium \ipa{"be@ri@m} & Barium \\
	P & phosphorus \ipa{"f6sf@r@s} & Phosphor & W & tungsten \ipa{"t2Nst@n} & Wolfram \\
	S & sulfur \ipa{"s2lf@} & Schwefel & Pt & platinum \ipa{"pl\ae tIn@m} & Platin \\
	Cl & chlorine \ipa{"klO:ri:n} & Chlor & Au & gold \ipa{g@Uld} & Gold \\
	Ar & argon \ipa{"A:g6n} & Argon & Hg & mercury \ipa{"m3:kj@ri} & Quecksilber \\
	 & & & Pb & lead \ipa{led} & Blei \\
	\bottomrule
\end{tabular}}
\end{center}

\anmerkung{Die Halogene enden im Englischen auf \textit{-ine} (\textit{chlorine}), ihre Ionen auf
\textit{-ide} (\textit{chloride}). Die IUPAC empfiehlt die Schreibung \textit{sulfur}; in älteren
britischen Texten steht \textit{sulphur}. Amerikanisch: \textit{aluminum} \ipa{@"lu:mIn@m}. Im Deutschen
gelten heute Iod (früher Jod), Calcium (Kalzium), Oxid (Oxyd) und Ethan (Äthan).}

\platz\zwischen{Beispiele}{Examples}{Ig"zA:mp\textsyllabic{l}z}

\begin{lesen}
\lies{Na}{Sodium is a soft, silvery metal which reacts vigorously with water.}{Natrium ist ein weiches, silbrig glänzendes Metall, das heftig mit Wasser reagiert.}
\lies{Fe}{Iron \ipa{"aI@n} is the most widely used metal. (The r is silent.)}{Eisen ist das am häufigsten verwendete Metall.}
\lies{Pb}{Lead \ipa{led} is a heavy, soft metal -- not to be confused with the verb \textit{to lead} \ipa{li:d}.}{Blei ist ein schweres, weiches Metall.}
\lies{\ce{N2}, \ce{O2}}{Air is a mixture of about seventy-eight per cent nitrogen and twenty-one per cent oxygen.}{Luft ist ein Gemisch aus etwa achtundsiebzig Prozent Stickstoff und einundzwanzig Prozent Sauerstoff.}
\lies{Hg, Br}{Mercury and bromine are the only elements that are liquid at room temperature.}{Quecksilber und Brom sind die einzigen Elemente, die bei Raumtemperatur flüssig sind.}
\end{lesen}

%--------------------------------------------------------------------------------------------------
\newpage
\kapitel{Formeln und Reaktionsgleichungen lesen}{Reading formulae and equations}{"ri:dIN "fO:mjuli: @n I"kweIZ\textsyllabic{n}z}

\begin{center}
\begin{tikzpicture}[font=\small]
	\node[font=\LARGE] (g) at (0,0) {$\color{red}2\,\color{black}\mathrm{H}_{\color{blue}2}\mathrm{O}\,\color{green!50!black}\mathrm{(l)}$};
	\node[align=center] (a) at (-3.6,1.2) {coefficient\\[-2pt]\color{black!60}Koeffizient};
	\node[align=center] (b) at (0,1.6) {chemical symbol\\[-2pt]\color{black!60}Elementsymbol};
	\node[align=center] (c) at (-2.6,-1.3) {subscript\\[-2pt]\color{black!60}Index};
	\node[align=center] (d) at (3.6,1.2) {state symbol\\[-2pt]\color{black!60}Zustandssymbol};
	\draw[-{Stealth}, thin] (a) -- (-1.05,0.2); \draw[-{Stealth}, thin] (b) -- (-0.45,0.35); \draw[-{Stealth}, thin] (c) -- (-0.15,-0.3); \draw[-{Stealth}, thin] (d) -- (1.0,0.25);
\end{tikzpicture}
\end{center}

\begin{lesen}
\lies{\ce{H2O}}{H two O -- water}{H zwei O -- Wasser}
\lies{\ce{CO2}}{C O two -- carbon dioxide}{C O zwei -- Kohlenstoffdioxid}
\lies{\ce{H2SO4}}{H two S O four -- sulfuric acid}{H zwei S O vier -- Schwefelsäure}
\lies{\ce{Ca(OH)2}}{C A, O H twice \;/\; C A, open bracket O H close bracket two -- calcium hydroxide}{C A, O H zweimal \;/\; C A, in Klammern O H, zwei -- Calciumhydroxid}
\lies{\ce{Na+}, \ce{Cl-}}{N A plus -- a sodium ion; C L minus -- a chloride ion}{N A plus -- ein Natrium-Ion; C L minus -- ein Chlorid-Ion}
\lies{\ce{SO4^2-}}{S O four two minus -- the sulfate ion}{S O vier zwei minus -- das Sulfat-Ion}
\lies{(s), (l), (g), (aq)}{state symbols: solid, liquid, gas, aqueous \ipa{"eIkwi@s} (dissolved in water)}{Zustandssymbole: fest, flüssig, gasförmig, in Wasser gelöst}
\lies{\ce{->}}{gives \;/\; yields \;/\; reacts to form}{reagiert zu}
\lies{\ce{<=>}}{is in equilibrium with (a reversible reaction)}{steht im Gleichgewicht mit (umkehrbare Reaktion)}
\lies{\ce{^}, \ce{v}}{a gas is given off; a precipitate is formed}{ein Gas entweicht; ein Niederschlag fällt aus}
\lies{\ce{2 H2 + O2 -> 2 H2O}}{two H two plus O two gives two H two O. \newline In words: Hydrogen reacts with oxygen to form water.}{zwei H zwei plus O zwei reagiert zu zwei H zwei O. \newline In Worten: Wasserstoff reagiert mit Sauerstoff zu Wasser.}
\lies{$\Delta H = -286\eh{kJ/mol}$}{delta H equals minus two hundred and eighty-six kilojoules per mole: the reaction is exothermic}{Delta H gleich minus zweihundertsechsundachtzig Kilojoule pro Mol: Die Reaktion ist exotherm}
\end{lesen}

\platz\zwischen{Beispiele}{Examples}{Ig"zA:mp\textsyllabic{l}z}

\begin{lesen}
\lies{\ce{Zn(s) + 2 HCl(aq) -> ZnCl2(aq) + H2(g)}}{Zinc reacts with hydrochloric acid to form zinc chloride solution and hydrogen gas.}{Zink reagiert mit Salzsäure zu Zinkchloridlösung und gasförmigem Wasserstoff.}
\lies{\ce{CaCO3 -> CaO + CO2 ^}}{On heating, calcium carbonate decomposes into calcium oxide and carbon dioxide.}{Beim Erhitzen zerfällt Calciumcarbonat in Calciumoxid und Kohlenstoffdioxid.}
\lies{\ce{N2 + 3 H2 <=> 2 NH3}}{Nitrogen and hydrogen react to form ammonia; the reaction is reversible.}{Stickstoff und Wasserstoff reagieren zu Ammoniak; die Reaktion ist umkehrbar.}
\lies{\ce{Ba^2+ + SO4^2- -> BaSO4 v}}{Barium ions and sulfate ions form a white precipitate of barium sulfate.}{Barium-Ionen und Sulfat-Ionen bilden einen weißen Niederschlag von Bariumsulfat.}
\lies{Ausgleichen}{Balance the equation: there must be the same number of atoms of each element on both sides.}{Gleiche die Reaktionsgleichung aus: Auf beiden Seiten müssen von jedem Element gleich viele Atome stehen.}
\end{lesen}

%--------------------------------------------------------------------------------------------------
\newpage
\kapitel{Verbindungen aus zwei Elementen benennen}{Naming binary compounds}{"neImIN "baIn@ri "k6mpaUndz}

Der Name eines Salzes aus zwei Elementen besteht aus dem Namen des Metalls (Kation), bei Bedarf
seiner Wertigkeit als römischer Zahl und dem Namen des Nichtmetalls (Anion) mit der Endung
\textit{-ide} (deutsch: \textit{-id}).

\begin{center}
\begin{tikzpicture}[font=\small]
	\node[font=\LARGE, text=blue] (fe) at (0,0) {Fe};
	\node[font=\LARGE, text=red, anchor=west, inner sep=0pt] (cl) at (0.42,0) {Cl$_3$};
	\node[font=\Large, anchor=west] (n1) at (3.6,0.5) {\textcolor{blue}{iron}\textcolor{green!45!black}{(III)} \textcolor{red}{chloride}};
	\node[font=\Large, anchor=west, text=black!60] (n2) at (3.6,-0.5) {\textcolor{blue!50}{Eisen}\textcolor{green!45!black!60}{(III)}-\textcolor{red!55}{chlorid}};
	\draw[-{Stealth}, thick] (1.6,0) -- (3.3,0);
	\node[align=center, text=blue] (l1) at (2.6,1.9) {metal (cation)\\[-2pt]\color{black!60}Metall (Kation)};
	\node[align=center, text=green!45!black] (l2) at (5.0,1.9) {valency\\[-2pt]\color{black!60}Wertigkeit};
	\node[align=center, text=red] (l3) at (7.6,1.9) {non-metal + \textit{-ide}\\[-2pt]\color{black!60}Nichtmetall + \textit{-id}};
	\draw[thin, blue] (l1.south) -- (4.0,0.8); \draw[thin, green!45!black] (l2.south) -- (4.95,0.8); \draw[thin, red] (l3.south) -- (6.1,0.8);
\end{tikzpicture}
\end{center}

\begin{lesen}
\lies{\ce{NaCl}}{sodium chloride \ipa{"klO:raId}}{Natriumchlorid}
\lies{\ce{MgO}}{magnesium oxide \ipa{"6ksaId}}{Magnesiumoxid}
\lies{\ce{CaF2}}{calcium fluoride \ipa{"flO:raId}}{Calciumfluorid}
\lies{\ce{KBr}, \ce{KI}}{potassium bromide \ipa{"br@UmaId}, potassium iodide \ipa{"aI@daId}}{Kaliumbromid, Kaliumiodid}
\lies{\ce{ZnS}}{zinc sulfide \ipa{"s2lfaId}}{Zinksulfid}
\lies{\ce{FeCl2}, \ce{FeCl3}}{iron(II) chloride, iron(III) chloride -- read: iron two chloride, iron three chloride}{Eisen(II)-chlorid, Eisen(III)-chlorid -- lies: Eisen-zwei-chlorid, Eisen-drei-chlorid}
\lies{\ce{Cu2O}, \ce{CuO}}{copper(I) oxide, copper(II) oxide}{Kupfer(I)-oxid, Kupfer(II)-oxid}
\lies{\ce{PbO2}}{lead(IV) oxide \ipa{""led "fO:r "6ksaId}}{Blei(IV)-oxid}
\end{lesen}

\platz\zwischen{Zahlwörter bei Verbindungen aus Nichtmetallen}{Number prefixes for non-metal compounds}{"n2mb@ "pri:fIksIz f@ ""n6n"met\textsyllabic{l} "k6mpaUndz}

\begin{center}
{\small
\begin{tabular}{@{}l *{10}{c}@{}}
	\toprule
	 & 1 & 2 & 3 & 4 & 5 & 6 & 7 & 8 & 9 & 10 \\ \midrule
	englisch & mono- & di- & tri- & tetra- & penta- & hexa- & hepta- & octa- & nona- & deca- \\
	 & \ipa{"m6n@U} & \ipa{daI} & \ipa{traI} & \ipa{"tetr@} & \ipa{"pent@} & \ipa{"heks@} & \ipa{"hept@} & \ipa{"6kt@} & \ipa{"n6n@} & \ipa{"dek@} \\
	deutsch & mono- & di- & tri- & tetra- & penta- & hexa- & hepta- & octa- & nona- & deca- \\
	\bottomrule
\end{tabular}}
\end{center}

\begin{lesen}
\lies{\ce{CO}}{carbon monoxide \ipa{m@"n6ksaId}}{Kohlenstoffmonoxid}
\lies{\ce{CO2}}{carbon dioxide \ipa{daI"6ksaId}}{Kohlenstoffdioxid}
\lies{\ce{SO2}, \ce{SO3}}{sulfur dioxide, sulfur trioxide \ipa{traI"6ksaId}}{Schwefeldioxid, Schwefeltrioxid}
\lies{\ce{N2O}}{dinitrogen \ipa{daI"naItr@dZ@n} monoxide \;/\; nitrous oxide (laughing gas)}{Distickstoffmonoxid (Lachgas)}
\lies{\ce{NO}, \ce{NO2}}{nitrogen monoxide, nitrogen dioxide}{Stickstoffmonoxid, Stickstoffdioxid}
\lies{\ce{N2O4}}{dinitrogen tetroxide \ipa{te"tr6ksaId}}{Distickstofftetraoxid}
\lies{\ce{N2O5}}{dinitrogen pentoxide \ipa{pen"t6ksaId}}{Distickstoffpentaoxid}
\lies{\ce{HCl}, \ce{H2S}}{hydrogen chloride, hydrogen sulfide}{Chlorwasserstoff, Schwefelwasserstoff}
\lies{\ce{NH3}, \ce{CH4}}{ammonia \ipa{@"m@Uni@}, methane \ipa{"mi:TeIn}}{Ammoniak, Methan}
\end{lesen}

\anmerkung{In älteren englischen Texten stehen statt der römischen Zahlen lateinische Namen: die
Endung \textit{-ous} für die niedrigere, \textit{-ic} für die höhere Wertigkeit -- \textit{ferrous}
\ipa{"fer@s} \textit{chloride} = iron(II) chloride, \textit{ferric} \ipa{"ferIk} \textit{chloride} =
iron(III) chloride, \textit{cuprous} \ipa{"kju:pr@s} und \textit{cupric} \ipa{"kju:prIk} für Kupfer(I) und Kupfer(II).}

\platz\zwischen{Beispiele}{Examples}{Ig"zA:mp\textsyllabic{l}z}

\begin{lesen}
\lies{\ce{Al2O3}}{aluminium oxide -- no Roman numeral is needed, because aluminium always has a valency of three}{Aluminiumoxid -- die römische Zahl entfällt, weil Aluminium immer dreiwertig ist}
\lies{\ce{Fe2O3}}{iron(III) oxide: the main constituent of rust}{Eisen(III)-oxid: der Hauptbestandteil von Rost}
\lies{\ce{P4O10}}{tetraphosphorus decoxide}{Tetraphosphordecaoxid}
\lies{\ce{CCl4}}{carbon tetrachloride \;/\; tetrachloromethane}{Tetrachlormethan}
\lies{\ce{CO}}{Carbon monoxide is a colourless, odourless and very poisonous gas.}{Kohlenstoffmonoxid ist ein farbloses, geruchloses und sehr giftiges Gas.}
\end{lesen}

%--------------------------------------------------------------------------------------------------
\newpage
\kapitel{Säuren, Basen und ihre Salze benennen}{Naming acids, bases and their salts}{"neImIN "\ae sIdz, "beIsIz @n De@ "sO:lts}

\begin{center}
\begin{tikzpicture}[font=\small]
	\foreach \y/\a/\s/\da/\ds in {1.6/hydro\ldots\textbf{ic} acid/\ldots\textbf{ide}/\ldots wasserstoffsäure/\ldots\textbf{id},
			0.8/\ldots\textbf{ic} acid/\ldots\textbf{ate}/\ldots säure/\ldots\textbf{at}, 0/\ldots\textbf{ous} acid/\ldots\textbf{ite}/\ldots ige Säure/\ldots\textbf{it}} {
		\node[anchor=east] at (0,\y) {\a}; \draw[-{Stealth}, thick, red] (0.3,\y) -- (1.5,\y); \node[anchor=west] at (1.8,\y) {\s};
		\node[anchor=east, black!60] at (8.6,\y) {\da}; \draw[-{Stealth}, thick, red!50] (8.9,\y) -- (10.1,\y); \node[anchor=west, black!60] at (10.4,\y) {\ds};}
	\node[font=\small\bfseries] at (0.9,2.35) {acid $\rightarrow$ salt}; \node[font=\small\bfseries, black!60] at (9.5,2.35) {Säure $\rightarrow$ Salz};
\end{tikzpicture}
\end{center}

\begin{center}
{\small\renewcommand{\arraystretch}{1.3}
\begin{tabular}{@{}>{\raggedright}p{1.9cm} >{\raggedright}p{4.5cm} >{\raggedright}p{3.7cm} >{\raggedright}p{3.2cm} >{\raggedright\arraybackslash}p{2.1cm}@{}}
	\toprule
	\textbf{Formel} & \textbf{Säure (englisch)} & \textbf{Säure (deutsch)} & \textbf{Salz (englisch)} & \textbf{Salz (deutsch)} \\ \midrule
	\ce{HF} & hydrofluoric acid \ipa{""haIdr@"flO:rIk} & Flusssäure & fluoride & Fluorid \\
	\ce{HCl} & hydrochloric acid \ipa{""haIdr@"kl6rIk} & Salzsäure & chloride & Chlorid \\
	\ce{HBr} & hydrobromic acid \ipa{""haIdr@"br@UmIk} & Bromwasserstoffsäure & bromide & Bromid \\
	\ce{H2S} & hydrogen sulfide (solution) & Schwefelwasserstoff & sulfide & Sulfid \\
	\ce{H2SO4} & sulfuric acid \ipa{s2l"fjU@rIk} & Schwefelsäure & sulfate \ipa{"s2lfeIt} & Sulfat \\
	\ce{H2SO3} & sulfurous acid \ipa{"s2lf@r@s} & schweflige Säure & sulfite \ipa{"s2lfaIt} & Sulfit \\
	\ce{HNO3} & nitric acid \ipa{"naItrIk} & Salpetersäure & nitrate \ipa{"naItreIt} & Nitrat \\
	\ce{HNO2} & nitrous acid \ipa{"naItr@s} & salpetrige Säure & nitrite \ipa{"naItraIt} & Nitrit \\
	\ce{H3PO4} & phosphoric acid \ipa{f6s"f6rIk} & Phosphorsäure & phosphate \ipa{"f6sfeIt} & Phosphat \\
	\ce{H2CO3} & carbonic acid \ipa{kA:"b6nIk} & Kohlensäure & carbonate \ipa{"kA:b@neIt} & Carbonat \\
	\ce{HClO} & hypochlorous acid \ipa{""haIp@"klO:r@s} & hypochlorige Säure & hypochlorite \ipa{""haIp@"klO:raIt} & Hypochlorit \\
	\ce{HClO3} & chloric acid \ipa{"klO:rIk} & Chlorsäure & chlorate \ipa{"klO:reIt} & Chlorat \\
	\ce{HClO4} & perchloric acid \ipa{p@"klO:rIk} & Perchlorsäure & perchlorate \ipa{p@"klO:reIt} & Perchlorat \\
	\ce{CH3COOH} & ethanoic acid \ipa{""eT@"n@UIk} \newline (acetic acid \ipa{@"si:tIk}) & Essigsäure (Ethansäure) & ethanoate \ipa{""eT@"n@UeIt} (acetate \ipa{"\ae sIteIt}) & Acetat (Ethanoat) \\
	\bottomrule
\end{tabular}}
\end{center}

\platz\zwischen{Basen und bekannte Salze}{Bases and common salts}{"beIsIz @n "k6m@n "sO:lts}

\begin{lesen}
\lies{\ce{NaOH}}{sodium hydroxide \ipa{haI"dr6ksaId} (caustic soda \ipa{""kO:stIk "s@Ud@}); in water: sodium hydroxide solution}{Natriumhydroxid (Ätznatron); in Wasser gelöst: Natronlauge}
\lies{\ce{KOH}}{potassium hydroxide (caustic potash \ipa{"p6t\ae S}); in water: potassium hydroxide solution}{Kaliumhydroxid (Ätzkali); in Wasser gelöst: Kalilauge}
\lies{\ce{Ca(OH)2}}{calcium hydroxide (slaked lime \ipa{""sleIkt "laIm}); its solution is called limewater}{Calciumhydroxid (gelöschter Kalk); die Lösung heißt Kalkwasser}
\lies{\ce{NH3(aq)}}{ammonia solution (aqueous ammonia)}{Ammoniaklösung (Ammoniakwasser)}
\lies{\ce{NaCl}}{sodium chloride (common salt, table salt)}{Natriumchlorid (Kochsalz)}
\lies{\ce{CaCO3}}{calcium carbonate (limestone, chalk, marble)}{Calciumcarbonat (Kalkstein, Kreide, Marmor)}
\lies{\ce{NaHCO3}}{sodium hydrogencarbonate (sodium bicarbonate, baking soda)}{Natriumhydrogencarbonat (Natron)}
\lies{\ce{KNO3}}{potassium nitrate (saltpetre \ipa{""sO:lt"pi:t@})}{Kaliumnitrat (Kalisalpeter)}
\lies{\ce{CuSO4.5H2O}}{copper(II) sulfate pentahydrate \ipa{""pent@"haIdreIt} (hydrated copper sulfate)}{Kupfer(II)-sulfat-Pentahydrat (blaues Kupfersulfat)}
\end{lesen}

\anmerkung{Für \ce{HCl} gibt es zwei Namen: Das Gas heißt \textit{hydrogen chloride}, seine Lösung in
Wasser \textit{hydrochloric acid}. Die Endung \textit{-ide} wird \ipa{aId} gesprochen, \textit{-ate}
\ipa{eIt} und \textit{-ite} \ipa{aIt} -- \textit{sulfide, sulfate, sulfite} klingen also deutlich verschieden.}

\platz\zwischen{Beispiele}{Examples}{Ig"zA:mp\textsyllabic{l}z}

\begin{lesen}
\lies{\ce{H2SO4 + 2 NaOH -> Na2SO4 + 2 H2O}}{Sulfuric acid reacts with sodium hydroxide to form sodium sulfate and water.}{Schwefelsäure reagiert mit Natriumhydroxid zu Natriumsulfat und Wasser.}
\end{lesen}

%--------------------------------------------------------------------------------------------------
\newpage
\kapitel{Der pH-Wert und Indikatoren}{pH and indicators}{""pi: "eItS @n "IndIkeIt@z}

\begin{center}
\begin{tikzpicture}[font=\small]
	\phskala{0.9}
	\foreach \k in {0,...,14} {\node[below] at (0.9*\k+0.45,0) {\k};}
	\draw[{Stealth}-, thick] (0,1.6) -- node[above] {acidic -- \color{black!60}sauer} (5.9,1.6);
	\draw[-{Stealth}, thick] (7.6,1.6) -- node[above] {alkaline -- \color{black!60}alkalisch (basisch)} (13.5,1.6);
	\node[above] at (6.75,1.5) {neutral};
	\foreach \k/\e/\d/\y in {1/stomach acid/Magensäure/-0.75, 3/vinegar/Essig/-1.75, 5/rain/Regen/-0.75, 7/pure water/reines Wasser/-1.75, 9/soap solution/Seifenlösung/-0.75, 11/ammonia solution/Ammoniaklösung/-1.75, 13/drain cleaner/Rohrreiniger/-0.75}
		{\draw[thin] (0.9*\k+0.45,-0.45) -- (0.9*\k+0.45,\y); \node[align=center, font=\footnotesize, anchor=north] at (0.9*\k+0.45,\y) {\e\\[-2pt]\color{black!60}\d};}
\end{tikzpicture}
\end{center}

\begin{lesen}
\lies{pH $= 3$}{p H three: the solution is acidic}{p H drei: Die Lösung ist sauer}
\lies{$\mathrm{pH} = -\lg c(\ce{H3O+})$}{pH equals minus log of the concentration of oxonium ions}{pH gleich minus Logarithmus der Konzentration der Oxonium-Ionen}
\lies{Lackmus}{litmus \ipa{"lItm@s}: red in acidic solutions, blue in alkaline solutions}{Lackmus: in saurer Lösung rot, in alkalischer Lösung blau}
\lies{Phenolphthalein}{phenolphthalein \ipa{""fi:n6l"TeIli:n}: colourless in acidic and neutral solutions, pink in alkaline solutions}{Phenolphthalein: in saurer und neutraler Lösung farblos, in alkalischer Lösung rosa bis pink}
\lies{Bromthymolblau}{bromothymol \ipa{""br@Um@U"TaIm6l} blue: yellow in acidic, green in neutral and blue in alkaline solutions}{Bromthymolblau: in saurer Lösung gelb, in neutraler grün und in alkalischer blau}
\end{lesen}

\platz\zwischen{Beispiele}{Examples}{Ig"zA:mp\textsyllabic{l}z}

\begin{lesen}
\lies{$c(\ce{H3O+}) = 0{,}01\eh{mol/l}$ \newline pH $= 2$}{The concentration of oxonium ions is nought point nought one moles per litre, so the pH is two.}{Die Konzentration der Oxonium-Ionen beträgt null Komma null eins Mol pro Liter, der pH-Wert ist also zwei.}
\lies{pH $3 \rightarrow$ pH $2$}{A decrease of one pH unit means that the solution is ten times more acidic.}{Sinkt der pH-Wert um eine Einheit, ist die Lösung zehnmal so sauer.}
\end{lesen}

%--------------------------------------------------------------------------------------------------
\newpage
\kapitel{Organische Verbindungen benennen}{Naming organic compounds}{"neImIN O:"g\ae nIk "k6mpaUndz}

Der Stamm des Namens gibt die Zahl der Kohlenstoffatome an, die Endung die Stoffklasse.

\begin{center}
{\small\renewcommand{\arraystretch}{1.15}
\begin{tabular}{@{}c l l l @{\hspace{1cm}} c l l l@{}}
	\toprule
	\textbf{C} & \textbf{Formel} & \textbf{Englisch} & \textbf{Deutsch} & \textbf{C} & \textbf{Formel} & \textbf{Englisch} & \textbf{Deutsch} \\ \midrule
	1 & \ce{CH4} & methane \ipa{"mi:TeIn} & Methan & 6 & \ce{C6H14} & hexane \ipa{"hekseIn} & Hexan \\
	2 & \ce{C2H6} & ethane \ipa{"i:TeIn} & Ethan & 7 & \ce{C7H16} & heptane \ipa{"hepteIn} & Heptan \\
	3 & \ce{C3H8} & propane \ipa{"pr@UpeIn} & Propan & 8 & \ce{C8H18} & octane \ipa{"6kteIn} & Octan \\
	4 & \ce{C4H10} & butane \ipa{"bju:teIn} & Butan & 9 & \ce{C9H20} & nonane \ipa{"n@UneIn} & Nonan \\
	5 & \ce{C5H12} & pentane \ipa{"penteIn} & Pentan & 10 & \ce{C10H22} & decane \ipa{"dekeIn} & Decan \\
	\bottomrule
\end{tabular}}
\end{center}

\begin{center}
{\footnotesize\renewcommand{\arraystretch}{1.25}
\begin{tabular}{@{}l l l l l@{}}
	\toprule
	\textbf{Stoffklasse (englisch)} & \textbf{deutsch} & \textbf{Endung} & \textbf{Beispiel (englisch)} & \textbf{deutsch} \\ \midrule
	alkanes & Alkane & -ane \,/\, -an & ethane & Ethan \\
	alkenes & Alkene & -ene \,/\, -en & ethene \ipa{"eTi:n} & Ethen \\
	alkynes & Alkine & -yne \,/\, -in & ethyne \ipa{"i:TaIn} & Ethin \\
	alcohols & Alkohole (Alkanole) & -ol \,/\, -ol & ethanol \ipa{"eT@n6l} & Ethanol \\
	aldehydes & Aldehyde (Alkanale) & -al \,/\, -al & ethanal \ipa{"eT@n\ae l} & Ethanal \\
	ketones & Ketone (Alkanone) & -one \,/\, -on & propanone & Propanon \\
	carboxylic acids & Carbonsäuren & -oic acid \,/\, -säure & ethanoic acid & Ethansäure \\
	esters & Ester & -yl -oate \,/\, -ester & ethyl ethanoate & Ethansäureethylester \\
	\bottomrule
\end{tabular}}
\end{center}

\begin{lesen}
\lies{\struktur{C(-[:120]H)(-[:240]H)=C(-[:60]H)(-[:-60]H)}}{ethene (common name: ethylene \ipa{"eT@li:n})}{Ethen (Trivialname: Ethylen)}
\lies{\struktur{H-C~C-H}}{ethyne (common name: acetylene \ipa{@"set@li:n})}{Ethin (Trivialname: Acetylen)}
\lies{\struktur{H_3C-CH(-[6]CH_3)-CH_3}}{2-methylpropane -- read: two methyl \ipa{"meTIl} propane (common name: isobutane)}{2-Methylpropan -- lies: zwei Methylpropan (Trivialname: Isobutan)}
\lies{\struktur{H_3C-CH_2-OH}}{ethanol (ethyl alcohol)}{Ethanol (Ethylalkohol, Trinkalkohol)}
\lies{\struktur{HO-CH_2-CH(-[6]OH)-CH_2-OH}}{propane-1,2,3-triol -- read: propane one two three triol (common name: glycerol \ipa{"glIs@r6l})}{Propan-1,2,3-triol (Trivialname: Glycerin)}
\lies{\struktur{H-C(=[2]O)-H}}{methanal (common name: formaldehyde \ipa{fO:"m\ae ldIhaId})}{Methanal (Trivialname: Formaldehyd)}
\lies{\struktur{H_3C-C(=[2]O)-CH_3}}{propanone (common name: acetone \ipa{"\ae sIt@Un})}{Propanon (Trivialname: Aceton)}
\lies{\struktur{H-C(=[2]O)-OH}}{methanoic acid \ipa{""meT@"n@UIk} (common name: formic \ipa{"fO:mIk} acid)}{Methansäure (Trivialname: Ameisensäure)}
\lies{\struktur{H_3C-C(=[2]O)-OH}}{ethanoic acid (common name: acetic acid)}{Ethansäure (Trivialname: Essigsäure)}
\lies{\struktur{H_3C-C(=[2]O)-O-CH_2-CH_3}}{ethyl ethanoate (common name: ethyl acetate)}{Ethansäureethylester (Trivialname: Essigsäureethylester)}
\lies{\struktur{*6(-=-=-=)}}{benzene \ipa{"benzi:n}}{Benzol (Benzen)}
\lies{\struktur{*6(-=-=(-CH_3)-=)}}{methylbenzene (common name: toluene \ipa{"t6lju:i:n})}{Methylbenzol (Trivialname: Toluol)}
\lies{\struktur{*6(-=-=(-OH)-=)}}{phenol \ipa{"fi:n6l}}{Phenol (Hydroxybenzol)}
\end{lesen}

\anmerkung{Falsche Freunde: Benzol heißt auf Englisch \textit{benzene}; das deutsche Benzin ist
\textit{petrol} (britisch) bzw.\ \textit{gasoline} (amerikanisch). Das englische \textit{benzine}
bezeichnet Waschbenzin. Stellungsziffern liest man einzeln: \textit{but-1-ene} -- „but one ene“,
\textit{propan-2-ol} -- „propan two ol“, \textit{2,2-dimethylpropane} -- „two two dimethylpropane“.}

\platz\zwischen{Beispiele}{Examples}{Ig"zA:mp\textsyllabic{l}z}

\begin{lesen}
\lies{\ce{C3H8}}{Propane is the third member of the homologous series of the alkanes.}{Propan ist das dritte Glied der homologen Reihe der Alkane.}
\lies{\ce{CH3-CH=CH2}}{propene: an alkene with three carbon atoms}{Propen: ein Alken mit drei Kohlenstoffatomen}
\lies{\ce{CH3-CH(OH)-CH3}}{propan-2-ol: the hydroxyl group is attached to the second carbon atom}{Propan-2-ol: Die Hydroxygruppe sitzt am zweiten Kohlenstoffatom}
\lies{\ce{CH2Cl-CH2Cl}}{1,2-dichloroethane -- read: one two dichloro ethane}{1,2-Dichlorethan -- lies: eins zwei Dichlorethan}
\lies{\ce{C2H5OH + CH3COOH <=> CH3COOC2H5 + H2O}}{Ethanol reacts with ethanoic acid to form ethyl ethanoate and water (esterification).}{Ethanol reagiert mit Ethansäure zu Ethansäureethylester und Wasser (Veresterung).}
\end{lesen}

%--------------------------------------------------------------------------------------------------
\newpage
\kapitel{Rechnen mit Stoffmengen}{Calculating with moles}{"k\ae lkjuleItIN wID "m@Ulz}

\begin{lesen}
\lies{$n = \dfrac{m}{M}$}{$n$ equals $m$ over capital $M$: the amount of substance equals the mass divided by the molar mass}{$n$ gleich $m$ durch groß $M$: Die Stoffmenge ist die Masse durch die molare Masse}
\lies{$N = n \cdot N_A$}{the number of particles equals the amount of substance times the Avogadro constant}{die Teilchenanzahl ist die Stoffmenge mal die Avogadro-Konstante}
\lies{$c = \dfrac{n}{V}$}{$c$ equals $n$ over $V$: the concentration equals the amount of substance divided by the volume of the solution}{$c$ gleich $n$ durch $V$: Die Konzentration ist die Stoffmenge durch das Volumen der Lösung}
\lies{$V = n \cdot V_m$}{the volume of a gas equals the amount of substance times the molar volume}{das Volumen eines Gases ist die Stoffmenge mal das molare Volumen}
\lies{$c(\ce{HCl}) = 0{,}1\eh{mol/l}$}{the concentration of the hydrochloric acid is nought point one moles per litre (nought point one molar)}{die Konzentration der Salzsäure beträgt null Komma eins Mol pro Liter}
\lies{$M(\ce{NaCl}) = 58{,}5\eh{g/mol}$}{the molar mass of sodium chloride is fifty-eight point five grams per mole}{die molare Masse von Natriumchlorid beträgt achtundfünfzig Komma fünf Gramm pro Mol}
\end{lesen}

\anmerkung{Im Englischen schreibt man die Konzentration meist in $\mathrm{mol/dm^3}$ oder kurz M
(\textit{molar}): $0{,}1\eh{mol/l} = 0.1\eh{mol\,dm^{-3}} = 0.1\eh{M}$. Eckige Klammern stehen für
Konzentrationen: $[\ce{H3O+}]$.}

\platz\zwischen{Beispiele}{Examples}{Ig"zA:mp\textsyllabic{l}z}

\begin{lesen}
\lies{$M(\ce{CO2}) = (12 + 2 \cdot 16)\eh{g/mol}$ \newline $M(\ce{CO2}) = 44\eh{g/mol}$}{To find the molar mass, add up the atomic masses: twelve plus two times sixteen is forty-four grams per mole.}{Für die molare Masse addiert man die Atommassen: zwölf plus zwei mal sechzehn ergibt vierundvierzig Gramm pro Mol.}
\lies{\ce{2 Mg + O2 -> 2 MgO} \newline $n(\ce{Mg}) = 1\eh{mol}$}{What mass of magnesium oxide is formed when twenty-four point three grams of magnesium burn? First calculate the amount of magnesium: one mole.}{Welche Masse Magnesiumoxid entsteht, wenn vierundzwanzig Komma drei Gramm Magnesium verbrennen? Berechne zuerst die Stoffmenge des Magnesiums: ein Mol.}
\lies{$n(\ce{MgO}) = n(\ce{Mg}) = 1\eh{mol}$ \newline $m(\ce{MgO}) = 1\eh{mol} \cdot 40{,}3\eh{g/mol}$ \newline $m(\ce{MgO}) = 40{,}3\eh{g}$}{From the equation, two moles of magnesium give two moles of magnesium oxide. So forty point three grams of magnesium oxide are formed.}{Nach der Reaktionsgleichung entstehen aus zwei Mol Magnesium zwei Mol Magnesiumoxid. Es entstehen also vierzig Komma drei Gramm Magnesiumoxid.}
\lies{$m = c \cdot V \cdot M$ \newline $m = 0{,}1\eh{\tfrac{mol}{l}} \cdot 1\eh{l} \cdot 40\eh{\tfrac{g}{mol}} = 4\eh{g}$}{To make one litre of sodium hydroxide solution with a concentration of nought point one moles per litre, dissolve four grams of sodium hydroxide in water and make the solution up to one litre.}{Für einen Liter Natronlauge der Konzentration null Komma eins Mol pro Liter löst man vier Gramm Natriumhydroxid in Wasser und füllt auf einen Liter auf.}
\lies{$c(\ce{HCl}) = \dfrac{0{,}1\eh{mol/l} \cdot 25\eh{ml}}{20\eh{ml}}$ \newline $c(\ce{HCl}) = 0{,}125\eh{mol/l}$}{Twenty millilitres of hydrochloric acid are neutralised by twenty-five millilitres of sodium hydroxide solution of concentration nought point one moles per litre. The concentration of the acid is nought point one two five moles per litre.}{Zwanzig Milliliter Salzsäure werden von fünfundzwanzig Millilitern Natronlauge der Konzentration null Komma eins Mol pro Liter neutralisiert. Die Konzentration der Säure beträgt null Komma eins zwei fünf Mol pro Liter.}
\end{lesen}

%--------------------------------------------------------------------------------------------------
\newpage
\kapitel{Experimente beschreiben und sicher arbeiten}{Describing experiments and working safely}{dI"skraIbIN Ik"sperIm@nts @n "w3:kIN "seIfli}

\begin{center}
\begin{tikzpicture}[font=\small, line cap=round, line join=round]
	\begin{scope}[scale=1.15]
		\brenner{0}{0}
		\draw[black, very thick] (-1.3,2.6) -- (1.3,2.6) (-1.15,2.6) -- (-1.5,0) (1.15,2.6) -- (1.5,0);
		\draw[black!60, thick, densely dashed] (-1.4,2.72) -- (1.4,2.72);
		\fill[blue!12] (-0.9,2.78) rectangle (0.9,3.8); \draw[blue, very thick] (-1.05,4.7) -- (-0.9,4.6) -- (-0.9,2.78) -- (0.9,2.78) -- (0.9,4.7);
		\draw[blue, thick, rounded corners=1.5pt] (0.3,3.1) rectangle (0.48,5.3); \fill[red] (0.39,3.1) circle (0.13); \fill[red] (0.36,3.1) rectangle (0.42,4.2);
	\end{scope}
	\foreach \y/\e/\d/\px/\py in {5.4/thermometer/Thermometer/0.6/5.6, 4.5/beaker/Becherglas/1.1/4.6, 3.6/gauze/Drahtnetz/1.5/3.15, 2.7/tripod/Dreifuß/1.6/1.9, 1.8/Bunsen burner/Gasbrenner/0.3/1.0, 0.9/gas supply/Gaszufuhr/-1.1/0.15}
		{\node[anchor=west] (n) at (4,\y) {\e{} -- \color{black!60}\d}; \draw[-{Stealth}, thin] (n.west) -- (\px,\py);}
\end{tikzpicture}
\end{center}

Die Teile eines Versuchsprotokolls (\textit{lab report}):

\begin{lesen}
\lies{\textbf{Aufgabe}}{\textbf{Aim:} To investigate how temperature affects the solubility of potassium nitrate in water.}{Untersuche, wie die Temperatur die Löslichkeit von Kaliumnitrat in Wasser beeinflusst.}
\lies{\textbf{Geräte und Chemikalien}}{\textbf{Apparatus and chemicals:} beaker, tripod, gauze, Bunsen burner, thermometer, balance; potassium nitrate, distilled water}{Becherglas, Dreifuß, Drahtnetz, Gasbrenner, Thermometer, Waage; Kaliumnitrat, destilliertes Wasser}
\lies{\textbf{Durchführung}}{\textbf{Method:} Pour fifty millilitres of water into the beaker. Add potassium nitrate a spatula at a time and stir until no more dissolves. Heat the mixture and repeat.}{Gib fünfzig Milliliter Wasser in das Becherglas. Füge spatelweise Kaliumnitrat hinzu und rühre, bis sich nichts mehr löst. Erwärme das Gemisch und wiederhole den Vorgang.}
\lies{\textbf{Beobachtung}}{\textbf{Observations:} More potassium nitrate dissolves in hot water than in cold water. On cooling, crystals form.}{In heißem Wasser löst sich mehr Kaliumnitrat als in kaltem. Beim Abkühlen bilden sich Kristalle.}
\lies{\textbf{Auswertung}}{\textbf{Conclusion:} The solubility of potassium nitrate increases with temperature.}{Die Löslichkeit von Kaliumnitrat nimmt mit steigender Temperatur zu.}
\end{lesen}

\platz\zwischen{Beobachtungen und Nachweise}{Observations and tests}{""6bz@"veIS\textsyllabic{n}z @n "tests}

\begin{lesen}
\lies{Gasentwicklung}{Bubbles of gas are given off (effervescence \ipa{""ef@"ves=ns}).}{Es steigen Gasblasen auf (die Lösung schäumt).}
\lies{Farbänderung}{The solution turns from blue to colourless.}{Die Lösung entfärbt sich von Blau nach Farblos.}
\lies{Niederschlag}{A white precipitate forms.}{Es bildet sich ein weißer Niederschlag.}
\lies{Sauerstoff}{Test for oxygen: a glowing splint relights.}{Nachweis von Sauerstoff (Glimmspanprobe): Ein glimmender Holzspan flammt auf.}
\lies{Wasserstoff}{Test for hydrogen: a lighted splint gives a squeaky pop.}{Nachweis von Wasserstoff (Knallgasprobe): Ein brennender Span erzeugt ein pfeifendes Geräusch.}
\lies{Kohlenstoffdioxid}{Test for carbon dioxide: limewater turns milky.}{Nachweis von Kohlenstoffdioxid: Kalkwasser trübt sich.}
\lies{Flammenfärbung}{Flame test: sodium compounds colour the flame yellow, potassium compounds lilac.}{Flammenfärbung: Natriumverbindungen färben die Flamme gelb, Kaliumverbindungen violett.}
\end{lesen}

\platz\zwischen{Gefahren und Sicherheit}{Hazards and safety}{"h\ae z@dz @n "seIfti}

\begin{lesen}
\lies{entzündbar}{flammable \ipa{"fl\ae m@b=l}}{}
\lies{ätzend}{corrosive \ipa{k@"r@UsIv}}{}
\lies{giftig}{toxic \ipa{"t6ksIk}}{}
\lies{gesundheitsschädlich, reizend}{harmful, irritant \ipa{"IrIt@nt}}{}
\lies{brandfördernd}{oxidising \ipa{"6ksIdaIzIN}}{}
\lies{explosiv}{explosive \ipa{Ik"spl@UsIv}}{}
\lies{umweltgefährlich}{hazardous to the environment \ipa{"h\ae z@d@s}}{}
\end{lesen}

\platz\zwischen{Beispiele}{Examples}{Ig"zA:mp\textsyllabic{l}z}

\begin{lesen}
\lies{Schutz}{Wear safety goggles and tie back long hair.}{Trage eine Schutzbrille und binde lange Haare zusammen.}
\lies{Verdünnen}{Always add acid to water, never water to acid.}{Gib immer die Säure ins Wasser, nie das Wasser in die Säure.}
\lies{Geruchsprobe}{To smell a substance, waft the vapour towards your nose with your hand.}{Fächle dir zur Geruchsprobe die Dämpfe mit der Hand zu.}
\lies{Abzug}{Carry out the experiment in a fume cupboard.}{Führe den Versuch unter dem Abzug durch.}
\lies{Entsorgung}{Do not pour the residues down the sink; put them in the waste container.}{Gieße die Reste nicht in den Ausguss; gib sie in den Sammelbehälter.}
\end{lesen}


\newpage
%%%%  REGISTER  %%%%
\teil{Register}{Index}{"Indeks}
\section{Register Deutsch -- Englisch}
\begin{multicols}{2}\small\raggedright

\medskip{\large\textbf{\color{cornflower}A}}\par\nopagebreak
\textbf{Abdampfschale} (die) -- evaporating dish\dotfill\hyperref[e:48]{\pageref*{e:48}}\par
\textbf{abkühlen} -- cool (down)\dotfill\hyperref[e:71]{\pageref*{e:71}}\par
\textbf{abwiegen} -- weigh out\dotfill\hyperref[e:69]{\pageref*{e:69}}\par
\textbf{Abzug} (der) -- fume cupboard (AE: fume hood)\dotfill\hyperref[e:63]{\pageref*{e:63}}\par
\textbf{Additionsreaktion} (die) -- addition reaction\dotfill\hyperref[e:228]{\pageref*{e:228}}\par
\textbf{Akkumulator} (der) -- rechargeable battery\dotfill\hyperref[e:204]{\pageref*{e:204}}\par
\textbf{Aktivierungsenergie} (die) -- activation energy\dotfill\hyperref[e:141]{\pageref*{e:141}}\par
\textbf{Aldehyd} (das) -- aldehyde\dotfill\hyperref[e:221]{\pageref*{e:221}}\par
\textbf{Alkalimetalle} (die) -- alkali metals\dotfill\hyperref[e:88]{\pageref*{e:88}}\par
\textbf{Alkan} (das) -- alkane\dotfill\hyperref[e:210]{\pageref*{e:210}}\par
\textbf{Alken} (das) -- alkene\dotfill\hyperref[e:211]{\pageref*{e:211}}\par
\textbf{Alkin} (das) -- alkyne\dotfill\hyperref[e:212]{\pageref*{e:212}}\par
\textbf{Alkohol} (der) -- alcohol\dotfill\hyperref[e:219]{\pageref*{e:219}}\par
\textbf{alkoholische Gärung} (die) -- fermentation\dotfill\hyperref[e:247]{\pageref*{e:247}}\par
\textbf{Amin} (das) -- amine\dotfill\hyperref[e:226]{\pageref*{e:226}}\par
\textbf{Aminosäure} (die) -- amino acid\dotfill\hyperref[e:227]{\pageref*{e:227}}\par
\textbf{Ammoniak} (das) -- ammonia\dotfill\hyperref[e:182]{\pageref*{e:182}}\par
\textbf{Anion} (das) -- anion\dotfill\hyperref[e:95]{\pageref*{e:95}}\par
\textbf{Anode} (die) -- anode\dotfill\hyperref[e:198]{\pageref*{e:198}}\par
\textbf{Äquivalenzpunkt} (der) -- equivalence point\dotfill\hyperref[e:175]{\pageref*{e:175}}\par
\textbf{Aromat} (der) -- aromatic compound\dotfill\hyperref[e:217]{\pageref*{e:217}}\par
\textbf{Atom} (das) -- atom\dotfill\hyperref[e:72]{\pageref*{e:72}}\par
\textbf{Atombindung (Elektronenpaarbindung)} (die) -- covalent bond\dotfill\hyperref[e:109]{\pageref*{e:109}}\par
\textbf{Atomkern} (der) -- nucleus\dotfill\hyperref[e:73]{\pageref*{e:73}}\par
\textbf{ätzend} -- corrosive\dotfill\hyperref[e:172]{\pageref*{e:172}}\par
\textbf{Ausbeute} (die) -- yield\dotfill\hyperref[e:150]{\pageref*{e:150}}\par
\textbf{Ausgangsstoff (das Edukt)} (der) -- reactant\dotfill\hyperref[e:126]{\pageref*{e:126}}\par
\textbf{ausgleichen} -- balance (an equation)\dotfill\hyperref[e:130]{\pageref*{e:130}}\par
\textbf{Außenelektron} (das) -- outer electron, valence electron\dotfill\hyperref[e:78]{\pageref*{e:78}}\par
\textbf{Avogadro-Konstante} (die) -- Avogadro constant\dotfill\hyperref[e:250]{\pageref*{e:250}}\par

\medskip{\large\textbf{\color{cornflower}B}}\par\nopagebreak
\textbf{Base} (die) -- base\dotfill\hyperref[e:153]{\pageref*{e:153}}\par
\textbf{basisch, alkalisch} -- basic, alkaline\dotfill\hyperref[e:157]{\pageref*{e:157}}\par
\textbf{Becherglas} (das) -- beaker\dotfill\hyperref[e:40]{\pageref*{e:40}}\par
\textbf{Benzin} (das) -- petrol (AE: gasoline)\dotfill\hyperref[e:235]{\pageref*{e:235}}\par
\textbf{bindende Elektronenpaar} (das) -- bonding pair\dotfill\hyperref[e:110]{\pageref*{e:110}}\par
\textbf{brennbar} -- flammable, combustible\dotfill\hyperref[e:28]{\pageref*{e:28}}\par
\textbf{Brennstoffzelle} (die) -- fuel cell\dotfill\hyperref[e:203]{\pageref*{e:203}}\par
\textbf{Bürette} (die) -- burette\dotfill\hyperref[e:51]{\pageref*{e:51}}\par

\medskip{\large\textbf{\color{cornflower}C}}\par\nopagebreak
\textbf{Carbonsäure} (die) -- carboxylic acid\dotfill\hyperref[e:223]{\pageref*{e:223}}\par
\textbf{Chemikalie} (die) -- chemical\dotfill\hyperref[e:64]{\pageref*{e:64}}\par
\textbf{chemische Bindung} (die) -- chemical bond\dotfill\hyperref[e:105]{\pageref*{e:105}}\par
\textbf{chemische Formel} (die) -- chemical formula\dotfill\hyperref[e:101]{\pageref*{e:101}}\par
\textbf{chemische Gleichgewicht} (das) -- chemical equilibrium\dotfill\hyperref[e:145]{\pageref*{e:145}}\par
\textbf{chemische Reaktion} (die) -- chemical reaction\dotfill\hyperref[e:125]{\pageref*{e:125}}\par
\textbf{Chromatografie} (die) -- chromatography\dotfill\hyperref[e:38]{\pageref*{e:38}}\par
\textbf{Cracken} (das) -- cracking\dotfill\hyperref[e:234]{\pageref*{e:234}}\par

\medskip{\large\textbf{\color{cornflower}D}}\par\nopagebreak
\textbf{Dekantieren} (das) -- decanting\dotfill\hyperref[e:36]{\pageref*{e:36}}\par
\textbf{delokalisierten Elektronen} (die) -- delocalised electrons\dotfill\hyperref[e:124]{\pageref*{e:124}}\par
\textbf{Destillation} (die) -- distillation\dotfill\hyperref[e:35]{\pageref*{e:35}}\par
\textbf{Dichte} (die) -- density\dotfill\hyperref[e:24]{\pageref*{e:24}}\par
\textbf{Dipol} (der) -- dipole\dotfill\hyperref[e:117]{\pageref*{e:117}}\par
\textbf{Doppelbindung} (die) -- double bond\dotfill\hyperref[e:113]{\pageref*{e:113}}\par
\textbf{Dreifachbindung} (die) -- triple bond\dotfill\hyperref[e:114]{\pageref*{e:114}}\par
\textbf{Dreifuß (mit Drahtnetz)} (der) -- tripod (and gauze)\dotfill\hyperref[e:47]{\pageref*{e:47}}\par

\medskip{\large\textbf{\color{cornflower}E}}\par\nopagebreak
\textbf{Edelgase} (die) -- noble gases\dotfill\hyperref[e:91]{\pageref*{e:91}}\par
\textbf{Edelgaskonfiguration} (die) -- noble gas configuration\dotfill\hyperref[e:96]{\pageref*{e:96}}\par
\textbf{Edelmetall} (das) -- noble metal, precious metal\dotfill\hyperref[e:192]{\pageref*{e:192}}\par
\textbf{Eigenschaft} (die) -- property\dotfill\hyperref[e:20]{\pageref*{e:20}}\par
\textbf{Eindampfen} (das) -- evaporation\dotfill\hyperref[e:34]{\pageref*{e:34}}\par
\textbf{Einfachbindung} (die) -- single bond\dotfill\hyperref[e:112]{\pageref*{e:112}}\par
\textbf{Eiweiß (Protein)} (das) -- protein\dotfill\hyperref[e:244]{\pageref*{e:244}}\par
\textbf{elektrische Leitfähigkeit} (die) -- electrical conductivity\dotfill\hyperref[e:27]{\pageref*{e:27}}\par
\textbf{Elektrode} (die) -- electrode\dotfill\hyperref[e:197]{\pageref*{e:197}}\par
\textbf{Elektrolyse} (die) -- electrolysis\dotfill\hyperref[e:195]{\pageref*{e:195}}\par
\textbf{Elektrolyt} (der) -- electrolyte\dotfill\hyperref[e:196]{\pageref*{e:196}}\par
\textbf{Elektron} (das) -- electron\dotfill\hyperref[e:76]{\pageref*{e:76}}\par
\textbf{Elektronegativität} (die) -- electronegativity\dotfill\hyperref[e:98]{\pageref*{e:98}}\par
\textbf{Elektronenkonfiguration} (die) -- electron configuration\dotfill\hyperref[e:79]{\pageref*{e:79}}\par
\textbf{Elektronenschale} (die) -- electron shell\dotfill\hyperref[e:77]{\pageref*{e:77}}\par
\textbf{Elektronenübergang} (der) -- electron transfer\dotfill\hyperref[e:189]{\pageref*{e:189}}\par
\textbf{Element} (das) -- element\dotfill\hyperref[e:3]{\pageref*{e:3}}\par
\textbf{Elementsymbol} (das) -- chemical symbol\dotfill\hyperref[e:84]{\pageref*{e:84}}\par
\textbf{Eliminierung} (die) -- elimination\dotfill\hyperref[e:230]{\pageref*{e:230}}\par
\textbf{Emulsion} (die) -- emulsion\dotfill\hyperref[e:16]{\pageref*{e:16}}\par
\textbf{endotherm} -- endothermic\dotfill\hyperref[e:139]{\pageref*{e:139}}\par
\textbf{Erdalkalimetalle} (die) -- alkaline earth metals\dotfill\hyperref[e:89]{\pageref*{e:89}}\par
\textbf{Erdgas} (das) -- natural gas\dotfill\hyperref[e:232]{\pageref*{e:232}}\par
\textbf{Erdöl} (das) -- crude oil\dotfill\hyperref[e:231]{\pageref*{e:231}}\par
\textbf{erhitzen} -- heat\dotfill\hyperref[e:65]{\pageref*{e:65}}\par
\textbf{Erlenmeyerkolben} (der) -- conical flask\dotfill\hyperref[e:41]{\pageref*{e:41}}\par
\textbf{Erz} (das) -- ore\dotfill\hyperref[e:206]{\pageref*{e:206}}\par
\textbf{Ester} (der) -- ester\dotfill\hyperref[e:224]{\pageref*{e:224}}\par
\textbf{exotherm} -- exothermic\dotfill\hyperref[e:138]{\pageref*{e:138}}\par

\medskip{\large\textbf{\color{cornflower}F}}\par\nopagebreak
\textbf{Fällung} (die) -- precipitation\dotfill\hyperref[e:136]{\pageref*{e:136}}\par
\textbf{Farbe} (die) -- colour (AE: color)\dotfill\hyperref[e:21]{\pageref*{e:21}}\par
\textbf{fest, flüssig, gasförmig} -- solid, liquid, gaseous\dotfill\hyperref[e:29]{\pageref*{e:29}}\par
\textbf{Fett} (das) -- fat\dotfill\hyperref[e:240]{\pageref*{e:240}}\par
\textbf{Filterpapier} (das) -- filter paper\dotfill\hyperref[e:45]{\pageref*{e:45}}\par
\textbf{Filtrat} (das) -- filtrate\dotfill\hyperref[e:32]{\pageref*{e:32}}\par
\textbf{Filtrieren} (das) -- filtration\dotfill\hyperref[e:31]{\pageref*{e:31}}\par
\textbf{fraktionierte Destillation} (die) -- fractional distillation\dotfill\hyperref[e:233]{\pageref*{e:233}}\par
\textbf{freie Elektronenpaar} (das) -- lone pair\dotfill\hyperref[e:111]{\pageref*{e:111}}\par
\textbf{funktionelle Gruppe} (die) -- functional group\dotfill\hyperref[e:218]{\pageref*{e:218}}\par

\medskip{\large\textbf{\color{cornflower}G}}\par\nopagebreak
\textbf{galvanische Zelle} (die) -- galvanic cell, electrochemical cell\dotfill\hyperref[e:200]{\pageref*{e:200}}\par
\textbf{Galvanisieren} (das) -- electroplating\dotfill\hyperref[e:205]{\pageref*{e:205}}\par
\textbf{Gasbrenner} (der) -- Bunsen burner\dotfill\hyperref[e:46]{\pageref*{e:46}}\par
\textbf{gelöste Stoff} (der) -- solute\dotfill\hyperref[e:10]{\pageref*{e:10}}\par
\textbf{Geruch} (der) -- smell, odour\dotfill\hyperref[e:22]{\pageref*{e:22}}\par
\textbf{gesättigt} -- saturated\dotfill\hyperref[e:214]{\pageref*{e:214}}\par
\textbf{gesättigte Lösung} (die) -- saturated solution\dotfill\hyperref[e:14]{\pageref*{e:14}}\par
\textbf{Gesetz von der Erhaltung der Masse} (das) -- law of conservation of mass\dotfill\hyperref[e:149]{\pageref*{e:149}}\par
\textbf{gießen} -- pour\dotfill\hyperref[e:68]{\pageref*{e:68}}\par
\textbf{Glasstab} (der) -- glass rod, stirring rod\dotfill\hyperref[e:53]{\pageref*{e:53}}\par
\textbf{Gleichgewichtskonstante} (die) -- equilibrium constant\dotfill\hyperref[e:146]{\pageref*{e:146}}\par

\medskip{\large\textbf{\color{cornflower}H}}\par\nopagebreak
\textbf{Halbzelle} (die) -- half-cell\dotfill\hyperref[e:201]{\pageref*{e:201}}\par
\textbf{Halogene} (die) -- halogens\dotfill\hyperref[e:90]{\pageref*{e:90}}\par
\textbf{Härte} (die) -- hardness\dotfill\hyperref[e:23]{\pageref*{e:23}}\par
\textbf{Hauptgruppe} (die) -- (main) group\dotfill\hyperref[e:87]{\pageref*{e:87}}\par
\textbf{heterogen} -- heterogeneous\dotfill\hyperref[e:7]{\pageref*{e:7}}\par
\textbf{Hochofen} (der) -- blast furnace\dotfill\hyperref[e:207]{\pageref*{e:207}}\par
\textbf{homogen} -- homogeneous\dotfill\hyperref[e:6]{\pageref*{e:6}}\par
\textbf{homologe Reihe} (die) -- homologous series\dotfill\hyperref[e:213]{\pageref*{e:213}}\par
\textbf{Hydroxid-Ion} (das) -- hydroxide ion\dotfill\hyperref[e:165]{\pageref*{e:165}}\par
\textbf{Hydroxygruppe} (die) -- hydroxyl group\dotfill\hyperref[e:220]{\pageref*{e:220}}\par

\medskip{\large\textbf{\color{cornflower}I}}\par\nopagebreak
\textbf{Index} (der) -- subscript\dotfill\hyperref[e:132]{\pageref*{e:132}}\par
\textbf{Indikator} (der) -- indicator\dotfill\hyperref[e:159]{\pageref*{e:159}}\par
\textbf{Ion} (das) -- ion\dotfill\hyperref[e:93]{\pageref*{e:93}}\par
\textbf{Ionenbindung} (die) -- ionic bond\dotfill\hyperref[e:106]{\pageref*{e:106}}\par
\textbf{Ionengitter} (das) -- ionic lattice\dotfill\hyperref[e:107]{\pageref*{e:107}}\par
\textbf{Ionisierungsenergie} (die) -- ionisation energy\dotfill\hyperref[e:99]{\pageref*{e:99}}\par
\textbf{Isomer} (das) -- isomer\dotfill\hyperref[e:216]{\pageref*{e:216}}\par
\textbf{Isotop} (das) -- isotope\dotfill\hyperref[e:83]{\pageref*{e:83}}\par

\medskip{\large\textbf{\color{cornflower}K}}\par\nopagebreak
\textbf{Katalysator} (der) -- catalyst\dotfill\hyperref[e:142]{\pageref*{e:142}}\par
\textbf{Kathode} (die) -- cathode\dotfill\hyperref[e:199]{\pageref*{e:199}}\par
\textbf{Kation} (das) -- cation\dotfill\hyperref[e:94]{\pageref*{e:94}}\par
\textbf{Keton} (das) -- ketone\dotfill\hyperref[e:222]{\pageref*{e:222}}\par
\textbf{Koeffizient (Stöchiometriezahl)} (der) -- coefficient\dotfill\hyperref[e:131]{\pageref*{e:131}}\par
\textbf{Kohlenhydrat} (das) -- carbohydrate\dotfill\hyperref[e:241]{\pageref*{e:241}}\par
\textbf{Kohlenwasserstoff} (der) -- hydrocarbon\dotfill\hyperref[e:209]{\pageref*{e:209}}\par
\textbf{konzentriert} -- concentrated\dotfill\hyperref[e:170]{\pageref*{e:170}}\par
\textbf{Korrosion} (die) -- corrosion\dotfill\hyperref[e:190]{\pageref*{e:190}}\par
\textbf{Kristall} (der) -- crystal\dotfill\hyperref[e:108]{\pageref*{e:108}}\par
\textbf{Kristallisation} (die) -- crystallisation\dotfill\hyperref[e:37]{\pageref*{e:37}}\par
\textbf{Kühler} (der) -- condenser\dotfill\hyperref[e:58]{\pageref*{e:58}}\par
\textbf{Kunststoff} (der) -- plastic\dotfill\hyperref[e:239]{\pageref*{e:239}}\par

\medskip{\large\textbf{\color{cornflower}L}}\par\nopagebreak
\textbf{Lackmus} (das) -- litmus\dotfill\hyperref[e:160]{\pageref*{e:160}}\par
\textbf{Lauge (alkalische Lösung)} (die) -- alkali, alkaline solution\dotfill\hyperref[e:154]{\pageref*{e:154}}\par
\textbf{Legierung} (die) -- alloy\dotfill\hyperref[e:17]{\pageref*{e:17}}\par
\textbf{Lewis-Formel} (die) -- Lewis structure, dot-and-cross diagram\dotfill\hyperref[e:123]{\pageref*{e:123}}\par
\textbf{lösen} -- dissolve\dotfill\hyperref[e:67]{\pageref*{e:67}}\par
\textbf{löslich} -- soluble\dotfill\hyperref[e:11]{\pageref*{e:11}}\par
\textbf{Löslichkeit} (die) -- solubility\dotfill\hyperref[e:13]{\pageref*{e:13}}\par
\textbf{Lösung} (die) -- solution\dotfill\hyperref[e:8]{\pageref*{e:8}}\par
\textbf{Lösungsmittel} (das) -- solvent\dotfill\hyperref[e:9]{\pageref*{e:9}}\par

\medskip{\large\textbf{\color{cornflower}M}}\par\nopagebreak
\textbf{Massenanteil} (der) -- mass fraction, percentage by mass\dotfill\hyperref[e:257]{\pageref*{e:257}}\par
\textbf{Massenkonzentration} (die) -- mass concentration\dotfill\hyperref[e:256]{\pageref*{e:256}}\par
\textbf{Massenwirkungsgesetz} (das) -- law of mass action\dotfill\hyperref[e:147]{\pageref*{e:147}}\par
\textbf{Massenzahl} (die) -- mass number\dotfill\hyperref[e:81]{\pageref*{e:81}}\par
\textbf{Maßlösung} (die) -- standard solution\dotfill\hyperref[e:176]{\pageref*{e:176}}\par
\textbf{Messzylinder} (der) -- measuring cylinder\dotfill\hyperref[e:43]{\pageref*{e:43}}\par
\textbf{Metall} (das) -- metal\dotfill\hyperref[e:18]{\pageref*{e:18}}\par
\textbf{Metallbindung} (die) -- metallic bond\dotfill\hyperref[e:115]{\pageref*{e:115}}\par
\textbf{Mol} (das) -- mole\dotfill\hyperref[e:249]{\pageref*{e:249}}\par
\textbf{molare Masse} (die) -- molar mass\dotfill\hyperref[e:252]{\pageref*{e:252}}\par
\textbf{molare Volumen} (das) -- molar volume\dotfill\hyperref[e:253]{\pageref*{e:253}}\par
\textbf{Molekül} (das) -- molecule\dotfill\hyperref[e:100]{\pageref*{e:100}}\par
\textbf{Monomer} (das) -- monomer\dotfill\hyperref[e:236]{\pageref*{e:236}}\par
\textbf{Mörser mit Pistill} (der) -- mortar and pestle\dotfill\hyperref[e:54]{\pageref*{e:54}}\par

\medskip{\large\textbf{\color{cornflower}N}}\par\nopagebreak
\textbf{Nachweis} (der) -- test (for)\dotfill\hyperref[e:151]{\pageref*{e:151}}\par
\textbf{Natronlauge} (die) -- sodium hydroxide solution\dotfill\hyperref[e:181]{\pageref*{e:181}}\par
\textbf{neutral} -- neutral\dotfill\hyperref[e:156]{\pageref*{e:156}}\par
\textbf{Neutralisation} (die) -- neutralisation\dotfill\hyperref[e:166]{\pageref*{e:166}}\par
\textbf{Neutron} (das) -- neutron\dotfill\hyperref[e:75]{\pageref*{e:75}}\par
\textbf{Nichtmetall} (das) -- non-metal\dotfill\hyperref[e:19]{\pageref*{e:19}}\par
\textbf{Niederschlag} (der) -- precipitate\dotfill\hyperref[e:137]{\pageref*{e:137}}\par

\medskip{\large\textbf{\color{cornflower}O}}\par\nopagebreak
\textbf{Oktettregel} (die) -- octet rule\dotfill\hyperref[e:97]{\pageref*{e:97}}\par
\textbf{Ordnungszahl} (die) -- atomic number, proton number\dotfill\hyperref[e:80]{\pageref*{e:80}}\par
\textbf{organische Chemie} (die) -- organic chemistry\dotfill\hyperref[e:208]{\pageref*{e:208}}\par
\textbf{Oxid} (das) -- oxide\dotfill\hyperref[e:188]{\pageref*{e:188}}\par
\textbf{Oxidation} (die) -- oxidation\dotfill\hyperref[e:183]{\pageref*{e:183}}\par
\textbf{Oxidationsmittel} (das) -- oxidising agent\dotfill\hyperref[e:186]{\pageref*{e:186}}\par
\textbf{Oxidationszahl} (die) -- oxidation number\dotfill\hyperref[e:122]{\pageref*{e:122}}\par
\textbf{Oxonium-Ion} (das) -- oxonium ion, hydronium ion\dotfill\hyperref[e:164]{\pageref*{e:164}}\par

\medskip{\large\textbf{\color{cornflower}P}}\par\nopagebreak
\textbf{Periode} (die) -- period\dotfill\hyperref[e:86]{\pageref*{e:86}}\par
\textbf{Periodensystem der Elemente} (das) -- periodic table\dotfill\hyperref[e:85]{\pageref*{e:85}}\par
\textbf{pH-Wert} (der) -- pH (value)\dotfill\hyperref[e:158]{\pageref*{e:158}}\par
\textbf{Pipette} (die) -- pipette\dotfill\hyperref[e:50]{\pageref*{e:50}}\par
\textbf{polare Atombindung} (die) -- polar covalent bond\dotfill\hyperref[e:116]{\pageref*{e:116}}\par
\textbf{Polymer} (das) -- polymer\dotfill\hyperref[e:237]{\pageref*{e:237}}\par
\textbf{Polymerisation} (die) -- polymerisation\dotfill\hyperref[e:238]{\pageref*{e:238}}\par
\textbf{Prinzip von Le Chatelier} (das) -- Le Chatelier's principle\dotfill\hyperref[e:148]{\pageref*{e:148}}\par
\textbf{Proton} (das) -- proton\dotfill\hyperref[e:74]{\pageref*{e:74}}\par
\textbf{Protonenakzeptor} (der) -- proton acceptor\dotfill\hyperref[e:163]{\pageref*{e:163}}\par
\textbf{Protonendonator} (der) -- proton donor\dotfill\hyperref[e:162]{\pageref*{e:162}}\par
\textbf{Pufferlösung} (die) -- buffer solution\dotfill\hyperref[e:177]{\pageref*{e:177}}\par

\medskip{\large\textbf{\color{cornflower}R}}\par\nopagebreak
\textbf{Reagenzglas} (das) -- test tube\dotfill\hyperref[e:39]{\pageref*{e:39}}\par
\textbf{Reaktionsenthalpie} (die) -- enthalpy change of reaction\dotfill\hyperref[e:140]{\pageref*{e:140}}\par
\textbf{Reaktionsgeschwindigkeit} (die) -- rate of reaction\dotfill\hyperref[e:143]{\pageref*{e:143}}\par
\textbf{Reaktionsgleichung} (die) -- chemical equation\dotfill\hyperref[e:129]{\pageref*{e:129}}\par
\textbf{Reaktionsprodukt} (das) -- product\dotfill\hyperref[e:127]{\pageref*{e:127}}\par
\textbf{Redoxreaktion} (die) -- redox reaction\dotfill\hyperref[e:185]{\pageref*{e:185}}\par
\textbf{Redoxreihe der Metalle} (die) -- reactivity series\dotfill\hyperref[e:194]{\pageref*{e:194}}\par
\textbf{Reduktion} (die) -- reduction\dotfill\hyperref[e:184]{\pageref*{e:184}}\par
\textbf{Reduktionsmittel} (das) -- reducing agent\dotfill\hyperref[e:187]{\pageref*{e:187}}\par
\textbf{Reinstoff} (der) -- pure substance\dotfill\hyperref[e:2]{\pageref*{e:2}}\par
\textbf{relative Atommasse} (die) -- relative atomic mass\dotfill\hyperref[e:82]{\pageref*{e:82}}\par
\textbf{relative Molekülmasse} (die) -- relative molecular mass\dotfill\hyperref[e:254]{\pageref*{e:254}}\par
\textbf{Rost} (der) -- rust\dotfill\hyperref[e:191]{\pageref*{e:191}}\par
\textbf{Rückstand} (der) -- residue\dotfill\hyperref[e:33]{\pageref*{e:33}}\par
\textbf{rühren} -- stir\dotfill\hyperref[e:66]{\pageref*{e:66}}\par
\textbf{Rundkolben} (der) -- round-bottomed flask\dotfill\hyperref[e:42]{\pageref*{e:42}}\par

\medskip{\large\textbf{\color{cornflower}S}}\par\nopagebreak
\textbf{Salpetersäure} (die) -- nitric acid\dotfill\hyperref[e:180]{\pageref*{e:180}}\par
\textbf{Salz} (das) -- salt\dotfill\hyperref[e:167]{\pageref*{e:167}}\par
\textbf{Salzsäure} (die) -- hydrochloric acid\dotfill\hyperref[e:178]{\pageref*{e:178}}\par
\textbf{sauer} -- acidic\dotfill\hyperref[e:155]{\pageref*{e:155}}\par
\textbf{Säure} (die) -- acid\dotfill\hyperref[e:152]{\pageref*{e:152}}\par
\textbf{Schmelztemperatur} (die) -- melting point\dotfill\hyperref[e:25]{\pageref*{e:25}}\par
\textbf{Schutzbrille} (die) -- safety goggles\dotfill\hyperref[e:62]{\pageref*{e:62}}\par
\textbf{schwache Säure} (die) -- weak acid\dotfill\hyperref[e:169]{\pageref*{e:169}}\par
\textbf{Schwefelsäure} (die) -- sulfuric acid\dotfill\hyperref[e:179]{\pageref*{e:179}}\par
\textbf{Seife} (die) -- soap\dotfill\hyperref[e:245]{\pageref*{e:245}}\par
\textbf{Siedetemperatur} (die) -- boiling point\dotfill\hyperref[e:26]{\pageref*{e:26}}\par
\textbf{Spatel} (der) -- spatula\dotfill\hyperref[e:52]{\pageref*{e:52}}\par
\textbf{Standardpotential} (das) -- standard electrode potential\dotfill\hyperref[e:202]{\pageref*{e:202}}\par
\textbf{Stärke} (die) -- starch\dotfill\hyperref[e:243]{\pageref*{e:243}}\par
\textbf{starke Säure} (die) -- strong acid\dotfill\hyperref[e:168]{\pageref*{e:168}}\par
\textbf{Stativ} (das) -- clamp stand\dotfill\hyperref[e:56]{\pageref*{e:56}}\par
\textbf{Stöchiometrie} (die) -- stoichiometry\dotfill\hyperref[e:258]{\pageref*{e:258}}\par
\textbf{Stoff} (der) -- substance\dotfill\hyperref[e:1]{\pageref*{e:1}}\par
\textbf{Stoffgemisch} (das) -- mixture\dotfill\hyperref[e:5]{\pageref*{e:5}}\par
\textbf{Stoffmenge} (die) -- amount of substance\dotfill\hyperref[e:248]{\pageref*{e:248}}\par
\textbf{Stoffmengenkonzentration} (die) -- concentration\dotfill\hyperref[e:255]{\pageref*{e:255}}\par
\textbf{Stoffmengenverhältnis} (das) -- mole ratio\dotfill\hyperref[e:259]{\pageref*{e:259}}\par
\textbf{Stopfen} (der) -- stopper, bung\dotfill\hyperref[e:60]{\pageref*{e:60}}\par
\textbf{Strukturformel} (die) -- structural formula\dotfill\hyperref[e:103]{\pageref*{e:103}}\par
\textbf{Substitutionsreaktion} (die) -- substitution reaction\dotfill\hyperref[e:229]{\pageref*{e:229}}\par
\textbf{Summenformel} (die) -- molecular formula\dotfill\hyperref[e:102]{\pageref*{e:102}}\par
\textbf{Suspension} (die) -- suspension\dotfill\hyperref[e:15]{\pageref*{e:15}}\par
\textbf{Synthese} (die) -- synthesis\dotfill\hyperref[e:133]{\pageref*{e:133}}\par

\medskip{\large\textbf{\color{cornflower}T}}\par\nopagebreak
\textbf{Teilchenanzahl} (die) -- number of particles\dotfill\hyperref[e:251]{\pageref*{e:251}}\par
\textbf{Tensid} (das) -- surfactant, detergent\dotfill\hyperref[e:246]{\pageref*{e:246}}\par
\textbf{Thermometer} (das) -- thermometer\dotfill\hyperref[e:57]{\pageref*{e:57}}\par
\textbf{Tiegel} (der) -- crucible\dotfill\hyperref[e:49]{\pageref*{e:49}}\par
\textbf{Tiegelzange} (die) -- tongs\dotfill\hyperref[e:61]{\pageref*{e:61}}\par
\textbf{Titration} (die) -- titration\dotfill\hyperref[e:173]{\pageref*{e:173}}\par
\textbf{Titrationskurve} (die) -- titration curve\dotfill\hyperref[e:174]{\pageref*{e:174}}\par
\textbf{Traubenzucker (die Glucose)} (der) -- glucose\dotfill\hyperref[e:242]{\pageref*{e:242}}\par
\textbf{Trennverfahren} (das) -- separation technique\dotfill\hyperref[e:30]{\pageref*{e:30}}\par
\textbf{Trichter} (der) -- funnel\dotfill\hyperref[e:44]{\pageref*{e:44}}\par

\medskip{\large\textbf{\color{cornflower}U}}\par\nopagebreak
\textbf{Übergangsmetalle} (die) -- transition metals\dotfill\hyperref[e:92]{\pageref*{e:92}}\par
\textbf{Überschuss} (der) -- excess\dotfill\hyperref[e:260]{\pageref*{e:260}}\par
\textbf{Uhrglas} (das) -- watch glass\dotfill\hyperref[e:55]{\pageref*{e:55}}\par
\textbf{umkehrbare Reaktion} (die) -- reversible reaction\dotfill\hyperref[e:144]{\pageref*{e:144}}\par
\textbf{unedle Metall} (das) -- base metal\dotfill\hyperref[e:193]{\pageref*{e:193}}\par
\textbf{ungesättigt} -- unsaturated\dotfill\hyperref[e:215]{\pageref*{e:215}}\par
\textbf{Universalindikator} (der) -- universal indicator\dotfill\hyperref[e:161]{\pageref*{e:161}}\par
\textbf{unlöslich} -- insoluble\dotfill\hyperref[e:12]{\pageref*{e:12}}\par

\medskip{\large\textbf{\color{cornflower}V}}\par\nopagebreak
\textbf{Van-der-Waals-Kräfte} (die) -- van der Waals forces\dotfill\hyperref[e:120]{\pageref*{e:120}}\par
\textbf{Verbindung} (die) -- compound\dotfill\hyperref[e:4]{\pageref*{e:4}}\par
\textbf{Verbrennung} (die) -- combustion\dotfill\hyperref[e:135]{\pageref*{e:135}}\par
\textbf{verdünnen} -- dilute\dotfill\hyperref[e:70]{\pageref*{e:70}}\par
\textbf{verdünnt} -- dilute\dotfill\hyperref[e:171]{\pageref*{e:171}}\par
\textbf{Veresterung} (die) -- esterification\dotfill\hyperref[e:225]{\pageref*{e:225}}\par
\textbf{Verhältnisformel} (die) -- empirical formula\dotfill\hyperref[e:104]{\pageref*{e:104}}\par

\medskip{\large\textbf{\color{cornflower}W}}\par\nopagebreak
\textbf{Waage} (die) -- balance\dotfill\hyperref[e:59]{\pageref*{e:59}}\par
\textbf{Wasserstoffbrückenbindung} (die) -- hydrogen bond\dotfill\hyperref[e:118]{\pageref*{e:118}}\par
\textbf{Wertigkeit} (die) -- valency\dotfill\hyperref[e:121]{\pageref*{e:121}}\par
\textbf{Wortgleichung} (die) -- word equation\dotfill\hyperref[e:128]{\pageref*{e:128}}\par

\medskip{\large\textbf{\color{cornflower}Z}}\par\nopagebreak
\textbf{Zersetzung (Analyse)} (die) -- decomposition\dotfill\hyperref[e:134]{\pageref*{e:134}}\par
\textbf{zwischenmolekularen Kräfte} (die) -- intermolecular forces\dotfill\hyperref[e:119]{\pageref*{e:119}}\par
\end{multicols}

\newpage
\section{Index English -- German}
\begin{multicols}{2}\small\raggedright

\medskip{\large\textbf{\color{cornflower}(}}\par\nopagebreak
\textbf{(main) group} -- die Hauptgruppe\dotfill\hyperref[e:87]{\pageref*{e:87}}\par

\medskip{\large\textbf{\color{cornflower}A}}\par\nopagebreak
\textbf{acid} -- die Säure\dotfill\hyperref[e:152]{\pageref*{e:152}}\par
\textbf{acidic} -- sauer\dotfill\hyperref[e:155]{\pageref*{e:155}}\par
\textbf{activation energy} -- die Aktivierungsenergie\dotfill\hyperref[e:141]{\pageref*{e:141}}\par
\textbf{addition reaction} -- die Additionsreaktion\dotfill\hyperref[e:228]{\pageref*{e:228}}\par
\textbf{alcohol} -- der Alkohol\dotfill\hyperref[e:219]{\pageref*{e:219}}\par
\textbf{aldehyde} -- das Aldehyd\dotfill\hyperref[e:221]{\pageref*{e:221}}\par
\textbf{alkali metals} -- die Alkalimetalle\dotfill\hyperref[e:88]{\pageref*{e:88}}\par
\textbf{alkali, alkaline solution} -- die Lauge (alkalische Lösung)\dotfill\hyperref[e:154]{\pageref*{e:154}}\par
\textbf{alkaline earth metals} -- die Erdalkalimetalle\dotfill\hyperref[e:89]{\pageref*{e:89}}\par
\textbf{alkane} -- das Alkan\dotfill\hyperref[e:210]{\pageref*{e:210}}\par
\textbf{alkene} -- das Alken\dotfill\hyperref[e:211]{\pageref*{e:211}}\par
\textbf{alkyne} -- das Alkin\dotfill\hyperref[e:212]{\pageref*{e:212}}\par
\textbf{alloy} -- die Legierung\dotfill\hyperref[e:17]{\pageref*{e:17}}\par
\textbf{amine} -- das Amin\dotfill\hyperref[e:226]{\pageref*{e:226}}\par
\textbf{amino acid} -- die Aminosäure\dotfill\hyperref[e:227]{\pageref*{e:227}}\par
\textbf{ammonia} -- das Ammoniak\dotfill\hyperref[e:182]{\pageref*{e:182}}\par
\textbf{amount of substance} -- die Stoffmenge\dotfill\hyperref[e:248]{\pageref*{e:248}}\par
\textbf{anion} -- das Anion\dotfill\hyperref[e:95]{\pageref*{e:95}}\par
\textbf{anode} -- die Anode\dotfill\hyperref[e:198]{\pageref*{e:198}}\par
\textbf{aromatic compound} -- der Aromat\dotfill\hyperref[e:217]{\pageref*{e:217}}\par
\textbf{atom} -- das Atom\dotfill\hyperref[e:72]{\pageref*{e:72}}\par
\textbf{atomic number, proton number} -- die Ordnungszahl\dotfill\hyperref[e:80]{\pageref*{e:80}}\par
\textbf{Avogadro constant} -- die Avogadro-Konstante\dotfill\hyperref[e:250]{\pageref*{e:250}}\par

\medskip{\large\textbf{\color{cornflower}B}}\par\nopagebreak
\textbf{balance} -- die Waage\dotfill\hyperref[e:59]{\pageref*{e:59}}\par
\textbf{balance (an equation)} -- ausgleichen\dotfill\hyperref[e:130]{\pageref*{e:130}}\par
\textbf{base} -- die Base\dotfill\hyperref[e:153]{\pageref*{e:153}}\par
\textbf{base metal} -- das unedle Metall\dotfill\hyperref[e:193]{\pageref*{e:193}}\par
\textbf{basic, alkaline} -- basisch, alkalisch\dotfill\hyperref[e:157]{\pageref*{e:157}}\par
\textbf{beaker} -- das Becherglas\dotfill\hyperref[e:40]{\pageref*{e:40}}\par
\textbf{blast furnace} -- der Hochofen\dotfill\hyperref[e:207]{\pageref*{e:207}}\par
\textbf{boiling point} -- die Siedetemperatur\dotfill\hyperref[e:26]{\pageref*{e:26}}\par
\textbf{bonding pair} -- das bindende Elektronenpaar\dotfill\hyperref[e:110]{\pageref*{e:110}}\par
\textbf{buffer solution} -- die Pufferlösung\dotfill\hyperref[e:177]{\pageref*{e:177}}\par
\textbf{Bunsen burner} -- der Gasbrenner\dotfill\hyperref[e:46]{\pageref*{e:46}}\par
\textbf{burette} -- die Bürette\dotfill\hyperref[e:51]{\pageref*{e:51}}\par

\medskip{\large\textbf{\color{cornflower}C}}\par\nopagebreak
\textbf{carbohydrate} -- das Kohlenhydrat\dotfill\hyperref[e:241]{\pageref*{e:241}}\par
\textbf{carboxylic acid} -- die Carbonsäure\dotfill\hyperref[e:223]{\pageref*{e:223}}\par
\textbf{catalyst} -- der Katalysator\dotfill\hyperref[e:142]{\pageref*{e:142}}\par
\textbf{cathode} -- die Kathode\dotfill\hyperref[e:199]{\pageref*{e:199}}\par
\textbf{cation} -- das Kation\dotfill\hyperref[e:94]{\pageref*{e:94}}\par
\textbf{chemical} -- die Chemikalie\dotfill\hyperref[e:64]{\pageref*{e:64}}\par
\textbf{chemical bond} -- die chemische Bindung\dotfill\hyperref[e:105]{\pageref*{e:105}}\par
\textbf{chemical equation} -- die Reaktionsgleichung\dotfill\hyperref[e:129]{\pageref*{e:129}}\par
\textbf{chemical equilibrium} -- das chemische Gleichgewicht\dotfill\hyperref[e:145]{\pageref*{e:145}}\par
\textbf{chemical formula} -- die chemische Formel\dotfill\hyperref[e:101]{\pageref*{e:101}}\par
\textbf{chemical reaction} -- die chemische Reaktion\dotfill\hyperref[e:125]{\pageref*{e:125}}\par
\textbf{chemical symbol} -- das Elementsymbol\dotfill\hyperref[e:84]{\pageref*{e:84}}\par
\textbf{chromatography} -- die Chromatografie\dotfill\hyperref[e:38]{\pageref*{e:38}}\par
\textbf{clamp stand} -- das Stativ\dotfill\hyperref[e:56]{\pageref*{e:56}}\par
\textbf{coefficient} -- der Koeffizient (Stöchiometriezahl)\dotfill\hyperref[e:131]{\pageref*{e:131}}\par
\textbf{colour (AE: color)} -- die Farbe\dotfill\hyperref[e:21]{\pageref*{e:21}}\par
\textbf{combustion} -- die Verbrennung\dotfill\hyperref[e:135]{\pageref*{e:135}}\par
\textbf{compound} -- die Verbindung\dotfill\hyperref[e:4]{\pageref*{e:4}}\par
\textbf{concentrated} -- konzentriert\dotfill\hyperref[e:170]{\pageref*{e:170}}\par
\textbf{concentration} -- die Stoffmengenkonzentration\dotfill\hyperref[e:255]{\pageref*{e:255}}\par
\textbf{condenser} -- der Kühler\dotfill\hyperref[e:58]{\pageref*{e:58}}\par
\textbf{conical flask} -- der Erlenmeyerkolben\dotfill\hyperref[e:41]{\pageref*{e:41}}\par
\textbf{cool (down)} -- abkühlen\dotfill\hyperref[e:71]{\pageref*{e:71}}\par
\textbf{corrosion} -- die Korrosion\dotfill\hyperref[e:190]{\pageref*{e:190}}\par
\textbf{corrosive} -- ätzend\dotfill\hyperref[e:172]{\pageref*{e:172}}\par
\textbf{covalent bond} -- die Atombindung (Elektronenpaarbindung)\dotfill\hyperref[e:109]{\pageref*{e:109}}\par
\textbf{cracking} -- das Cracken\dotfill\hyperref[e:234]{\pageref*{e:234}}\par
\textbf{crucible} -- der Tiegel\dotfill\hyperref[e:49]{\pageref*{e:49}}\par
\textbf{crude oil} -- das Erdöl\dotfill\hyperref[e:231]{\pageref*{e:231}}\par
\textbf{crystal} -- der Kristall\dotfill\hyperref[e:108]{\pageref*{e:108}}\par
\textbf{crystallisation} -- die Kristallisation\dotfill\hyperref[e:37]{\pageref*{e:37}}\par

\medskip{\large\textbf{\color{cornflower}D}}\par\nopagebreak
\textbf{decanting} -- das Dekantieren\dotfill\hyperref[e:36]{\pageref*{e:36}}\par
\textbf{decomposition} -- die Zersetzung (Analyse)\dotfill\hyperref[e:134]{\pageref*{e:134}}\par
\textbf{delocalised electrons} -- die delokalisierten Elektronen\dotfill\hyperref[e:124]{\pageref*{e:124}}\par
\textbf{density} -- die Dichte\dotfill\hyperref[e:24]{\pageref*{e:24}}\par
\textbf{dilute} -- verdünnen\dotfill\hyperref[e:70]{\pageref*{e:70}}\par
\textbf{dilute} -- verdünnt\dotfill\hyperref[e:171]{\pageref*{e:171}}\par
\textbf{dipole} -- der Dipol\dotfill\hyperref[e:117]{\pageref*{e:117}}\par
\textbf{dissolve} -- lösen\dotfill\hyperref[e:67]{\pageref*{e:67}}\par
\textbf{distillation} -- die Destillation\dotfill\hyperref[e:35]{\pageref*{e:35}}\par
\textbf{double bond} -- die Doppelbindung\dotfill\hyperref[e:113]{\pageref*{e:113}}\par

\medskip{\large\textbf{\color{cornflower}E}}\par\nopagebreak
\textbf{electrical conductivity} -- die elektrische Leitfähigkeit\dotfill\hyperref[e:27]{\pageref*{e:27}}\par
\textbf{electrode} -- die Elektrode\dotfill\hyperref[e:197]{\pageref*{e:197}}\par
\textbf{electrolysis} -- die Elektrolyse\dotfill\hyperref[e:195]{\pageref*{e:195}}\par
\textbf{electrolyte} -- der Elektrolyt\dotfill\hyperref[e:196]{\pageref*{e:196}}\par
\textbf{electron} -- das Elektron\dotfill\hyperref[e:76]{\pageref*{e:76}}\par
\textbf{electron configuration} -- die Elektronenkonfiguration\dotfill\hyperref[e:79]{\pageref*{e:79}}\par
\textbf{electron shell} -- die Elektronenschale\dotfill\hyperref[e:77]{\pageref*{e:77}}\par
\textbf{electron transfer} -- der Elektronenübergang\dotfill\hyperref[e:189]{\pageref*{e:189}}\par
\textbf{electronegativity} -- die Elektronegativität\dotfill\hyperref[e:98]{\pageref*{e:98}}\par
\textbf{electroplating} -- das Galvanisieren\dotfill\hyperref[e:205]{\pageref*{e:205}}\par
\textbf{element} -- das Element\dotfill\hyperref[e:3]{\pageref*{e:3}}\par
\textbf{elimination} -- die Eliminierung\dotfill\hyperref[e:230]{\pageref*{e:230}}\par
\textbf{empirical formula} -- die Verhältnisformel\dotfill\hyperref[e:104]{\pageref*{e:104}}\par
\textbf{emulsion} -- die Emulsion\dotfill\hyperref[e:16]{\pageref*{e:16}}\par
\textbf{endothermic} -- endotherm\dotfill\hyperref[e:139]{\pageref*{e:139}}\par
\textbf{enthalpy change of reaction} -- die Reaktionsenthalpie\dotfill\hyperref[e:140]{\pageref*{e:140}}\par
\textbf{equilibrium constant} -- die Gleichgewichtskonstante\dotfill\hyperref[e:146]{\pageref*{e:146}}\par
\textbf{equivalence point} -- der Äquivalenzpunkt\dotfill\hyperref[e:175]{\pageref*{e:175}}\par
\textbf{ester} -- der Ester\dotfill\hyperref[e:224]{\pageref*{e:224}}\par
\textbf{esterification} -- die Veresterung\dotfill\hyperref[e:225]{\pageref*{e:225}}\par
\textbf{evaporating dish} -- die Abdampfschale\dotfill\hyperref[e:48]{\pageref*{e:48}}\par
\textbf{evaporation} -- das Eindampfen\dotfill\hyperref[e:34]{\pageref*{e:34}}\par
\textbf{excess} -- der Überschuss\dotfill\hyperref[e:260]{\pageref*{e:260}}\par
\textbf{exothermic} -- exotherm\dotfill\hyperref[e:138]{\pageref*{e:138}}\par

\medskip{\large\textbf{\color{cornflower}F}}\par\nopagebreak
\textbf{fat} -- das Fett\dotfill\hyperref[e:240]{\pageref*{e:240}}\par
\textbf{fermentation} -- die alkoholische Gärung\dotfill\hyperref[e:247]{\pageref*{e:247}}\par
\textbf{filter paper} -- das Filterpapier\dotfill\hyperref[e:45]{\pageref*{e:45}}\par
\textbf{filtrate} -- das Filtrat\dotfill\hyperref[e:32]{\pageref*{e:32}}\par
\textbf{filtration} -- das Filtrieren\dotfill\hyperref[e:31]{\pageref*{e:31}}\par
\textbf{flammable, combustible} -- brennbar\dotfill\hyperref[e:28]{\pageref*{e:28}}\par
\textbf{fractional distillation} -- die fraktionierte Destillation\dotfill\hyperref[e:233]{\pageref*{e:233}}\par
\textbf{fuel cell} -- die Brennstoffzelle\dotfill\hyperref[e:203]{\pageref*{e:203}}\par
\textbf{fume cupboard (AE: fume hood)} -- der Abzug\dotfill\hyperref[e:63]{\pageref*{e:63}}\par
\textbf{functional group} -- die funktionelle Gruppe\dotfill\hyperref[e:218]{\pageref*{e:218}}\par
\textbf{funnel} -- der Trichter\dotfill\hyperref[e:44]{\pageref*{e:44}}\par

\medskip{\large\textbf{\color{cornflower}G}}\par\nopagebreak
\textbf{galvanic cell, electrochemical cell} -- die galvanische Zelle\dotfill\hyperref[e:200]{\pageref*{e:200}}\par
\textbf{glass rod, stirring rod} -- der Glasstab\dotfill\hyperref[e:53]{\pageref*{e:53}}\par
\textbf{glucose} -- der Traubenzucker (die Glucose)\dotfill\hyperref[e:242]{\pageref*{e:242}}\par

\medskip{\large\textbf{\color{cornflower}H}}\par\nopagebreak
\textbf{half-cell} -- die Halbzelle\dotfill\hyperref[e:201]{\pageref*{e:201}}\par
\textbf{halogens} -- die Halogene\dotfill\hyperref[e:90]{\pageref*{e:90}}\par
\textbf{hardness} -- die Härte\dotfill\hyperref[e:23]{\pageref*{e:23}}\par
\textbf{heat} -- erhitzen\dotfill\hyperref[e:65]{\pageref*{e:65}}\par
\textbf{heterogeneous} -- heterogen\dotfill\hyperref[e:7]{\pageref*{e:7}}\par
\textbf{homogeneous} -- homogen\dotfill\hyperref[e:6]{\pageref*{e:6}}\par
\textbf{homologous series} -- die homologe Reihe\dotfill\hyperref[e:213]{\pageref*{e:213}}\par
\textbf{hydrocarbon} -- der Kohlenwasserstoff\dotfill\hyperref[e:209]{\pageref*{e:209}}\par
\textbf{hydrochloric acid} -- die Salzsäure\dotfill\hyperref[e:178]{\pageref*{e:178}}\par
\textbf{hydrogen bond} -- die Wasserstoffbrückenbindung\dotfill\hyperref[e:118]{\pageref*{e:118}}\par
\textbf{hydroxide ion} -- das Hydroxid-Ion\dotfill\hyperref[e:165]{\pageref*{e:165}}\par
\textbf{hydroxyl group} -- die Hydroxygruppe\dotfill\hyperref[e:220]{\pageref*{e:220}}\par

\medskip{\large\textbf{\color{cornflower}I}}\par\nopagebreak
\textbf{indicator} -- der Indikator\dotfill\hyperref[e:159]{\pageref*{e:159}}\par
\textbf{insoluble} -- unlöslich\dotfill\hyperref[e:12]{\pageref*{e:12}}\par
\textbf{intermolecular forces} -- die zwischenmolekularen Kräfte\dotfill\hyperref[e:119]{\pageref*{e:119}}\par
\textbf{ion} -- das Ion\dotfill\hyperref[e:93]{\pageref*{e:93}}\par
\textbf{ionic bond} -- die Ionenbindung\dotfill\hyperref[e:106]{\pageref*{e:106}}\par
\textbf{ionic lattice} -- das Ionengitter\dotfill\hyperref[e:107]{\pageref*{e:107}}\par
\textbf{ionisation energy} -- die Ionisierungsenergie\dotfill\hyperref[e:99]{\pageref*{e:99}}\par
\textbf{isomer} -- das Isomer\dotfill\hyperref[e:216]{\pageref*{e:216}}\par
\textbf{isotope} -- das Isotop\dotfill\hyperref[e:83]{\pageref*{e:83}}\par

\medskip{\large\textbf{\color{cornflower}K}}\par\nopagebreak
\textbf{ketone} -- das Keton\dotfill\hyperref[e:222]{\pageref*{e:222}}\par

\medskip{\large\textbf{\color{cornflower}L}}\par\nopagebreak
\textbf{law of conservation of mass} -- das Gesetz von der Erhaltung der Masse\dotfill\hyperref[e:149]{\pageref*{e:149}}\par
\textbf{law of mass action} -- das Massenwirkungsgesetz\dotfill\hyperref[e:147]{\pageref*{e:147}}\par
\textbf{Le Chatelier's principle} -- das Prinzip von Le Chatelier\dotfill\hyperref[e:148]{\pageref*{e:148}}\par
\textbf{Lewis structure, dot-and-cross diagram} -- die Lewis-Formel\dotfill\hyperref[e:123]{\pageref*{e:123}}\par
\textbf{litmus} -- das Lackmus\dotfill\hyperref[e:160]{\pageref*{e:160}}\par
\textbf{lone pair} -- das freie Elektronenpaar\dotfill\hyperref[e:111]{\pageref*{e:111}}\par

\medskip{\large\textbf{\color{cornflower}M}}\par\nopagebreak
\textbf{mass concentration} -- die Massenkonzentration\dotfill\hyperref[e:256]{\pageref*{e:256}}\par
\textbf{mass fraction, percentage by mass} -- der Massenanteil\dotfill\hyperref[e:257]{\pageref*{e:257}}\par
\textbf{mass number} -- die Massenzahl\dotfill\hyperref[e:81]{\pageref*{e:81}}\par
\textbf{measuring cylinder} -- der Messzylinder\dotfill\hyperref[e:43]{\pageref*{e:43}}\par
\textbf{melting point} -- die Schmelztemperatur\dotfill\hyperref[e:25]{\pageref*{e:25}}\par
\textbf{metal} -- das Metall\dotfill\hyperref[e:18]{\pageref*{e:18}}\par
\textbf{metallic bond} -- die Metallbindung\dotfill\hyperref[e:115]{\pageref*{e:115}}\par
\textbf{mixture} -- das Stoffgemisch\dotfill\hyperref[e:5]{\pageref*{e:5}}\par
\textbf{molar mass} -- die molare Masse\dotfill\hyperref[e:252]{\pageref*{e:252}}\par
\textbf{molar volume} -- das molare Volumen\dotfill\hyperref[e:253]{\pageref*{e:253}}\par
\textbf{mole} -- das Mol\dotfill\hyperref[e:249]{\pageref*{e:249}}\par
\textbf{mole ratio} -- das Stoffmengenverhältnis\dotfill\hyperref[e:259]{\pageref*{e:259}}\par
\textbf{molecular formula} -- die Summenformel\dotfill\hyperref[e:102]{\pageref*{e:102}}\par
\textbf{molecule} -- das Molekül\dotfill\hyperref[e:100]{\pageref*{e:100}}\par
\textbf{monomer} -- das Monomer\dotfill\hyperref[e:236]{\pageref*{e:236}}\par
\textbf{mortar and pestle} -- der Mörser mit Pistill\dotfill\hyperref[e:54]{\pageref*{e:54}}\par

\medskip{\large\textbf{\color{cornflower}N}}\par\nopagebreak
\textbf{natural gas} -- das Erdgas\dotfill\hyperref[e:232]{\pageref*{e:232}}\par
\textbf{neutral} -- neutral\dotfill\hyperref[e:156]{\pageref*{e:156}}\par
\textbf{neutralisation} -- die Neutralisation\dotfill\hyperref[e:166]{\pageref*{e:166}}\par
\textbf{neutron} -- das Neutron\dotfill\hyperref[e:75]{\pageref*{e:75}}\par
\textbf{nitric acid} -- die Salpetersäure\dotfill\hyperref[e:180]{\pageref*{e:180}}\par
\textbf{noble gas configuration} -- die Edelgaskonfiguration\dotfill\hyperref[e:96]{\pageref*{e:96}}\par
\textbf{noble gases} -- die Edelgase\dotfill\hyperref[e:91]{\pageref*{e:91}}\par
\textbf{noble metal, precious metal} -- das Edelmetall\dotfill\hyperref[e:192]{\pageref*{e:192}}\par
\textbf{non-metal} -- das Nichtmetall\dotfill\hyperref[e:19]{\pageref*{e:19}}\par
\textbf{nucleus} -- der Atomkern\dotfill\hyperref[e:73]{\pageref*{e:73}}\par
\textbf{number of particles} -- die Teilchenanzahl\dotfill\hyperref[e:251]{\pageref*{e:251}}\par

\medskip{\large\textbf{\color{cornflower}O}}\par\nopagebreak
\textbf{octet rule} -- die Oktettregel\dotfill\hyperref[e:97]{\pageref*{e:97}}\par
\textbf{ore} -- das Erz\dotfill\hyperref[e:206]{\pageref*{e:206}}\par
\textbf{organic chemistry} -- die organische Chemie\dotfill\hyperref[e:208]{\pageref*{e:208}}\par
\textbf{outer electron, valence electron} -- das Außenelektron\dotfill\hyperref[e:78]{\pageref*{e:78}}\par
\textbf{oxidation} -- die Oxidation\dotfill\hyperref[e:183]{\pageref*{e:183}}\par
\textbf{oxidation number} -- die Oxidationszahl\dotfill\hyperref[e:122]{\pageref*{e:122}}\par
\textbf{oxide} -- das Oxid\dotfill\hyperref[e:188]{\pageref*{e:188}}\par
\textbf{oxidising agent} -- das Oxidationsmittel\dotfill\hyperref[e:186]{\pageref*{e:186}}\par
\textbf{oxonium ion, hydronium ion} -- das Oxonium-Ion\dotfill\hyperref[e:164]{\pageref*{e:164}}\par

\medskip{\large\textbf{\color{cornflower}P}}\par\nopagebreak
\textbf{period} -- die Periode\dotfill\hyperref[e:86]{\pageref*{e:86}}\par
\textbf{periodic table} -- das Periodensystem der Elemente\dotfill\hyperref[e:85]{\pageref*{e:85}}\par
\textbf{petrol (AE: gasoline)} -- das Benzin\dotfill\hyperref[e:235]{\pageref*{e:235}}\par
\textbf{pH (value)} -- der pH-Wert\dotfill\hyperref[e:158]{\pageref*{e:158}}\par
\textbf{pipette} -- die Pipette\dotfill\hyperref[e:50]{\pageref*{e:50}}\par
\textbf{plastic} -- der Kunststoff\dotfill\hyperref[e:239]{\pageref*{e:239}}\par
\textbf{polar covalent bond} -- die polare Atombindung\dotfill\hyperref[e:116]{\pageref*{e:116}}\par
\textbf{polymer} -- das Polymer\dotfill\hyperref[e:237]{\pageref*{e:237}}\par
\textbf{polymerisation} -- die Polymerisation\dotfill\hyperref[e:238]{\pageref*{e:238}}\par
\textbf{pour} -- gießen\dotfill\hyperref[e:68]{\pageref*{e:68}}\par
\textbf{precipitate} -- der Niederschlag\dotfill\hyperref[e:137]{\pageref*{e:137}}\par
\textbf{precipitation} -- die Fällung\dotfill\hyperref[e:136]{\pageref*{e:136}}\par
\textbf{product} -- das Reaktionsprodukt\dotfill\hyperref[e:127]{\pageref*{e:127}}\par
\textbf{property} -- die Eigenschaft\dotfill\hyperref[e:20]{\pageref*{e:20}}\par
\textbf{protein} -- das Eiweiß (Protein)\dotfill\hyperref[e:244]{\pageref*{e:244}}\par
\textbf{proton} -- das Proton\dotfill\hyperref[e:74]{\pageref*{e:74}}\par
\textbf{proton acceptor} -- der Protonenakzeptor\dotfill\hyperref[e:163]{\pageref*{e:163}}\par
\textbf{proton donor} -- der Protonendonator\dotfill\hyperref[e:162]{\pageref*{e:162}}\par
\textbf{pure substance} -- der Reinstoff\dotfill\hyperref[e:2]{\pageref*{e:2}}\par

\medskip{\large\textbf{\color{cornflower}R}}\par\nopagebreak
\textbf{rate of reaction} -- die Reaktionsgeschwindigkeit\dotfill\hyperref[e:143]{\pageref*{e:143}}\par
\textbf{reactant} -- der Ausgangsstoff (das Edukt)\dotfill\hyperref[e:126]{\pageref*{e:126}}\par
\textbf{reactivity series} -- die Redoxreihe der Metalle\dotfill\hyperref[e:194]{\pageref*{e:194}}\par
\textbf{rechargeable battery} -- der Akkumulator\dotfill\hyperref[e:204]{\pageref*{e:204}}\par
\textbf{redox reaction} -- die Redoxreaktion\dotfill\hyperref[e:185]{\pageref*{e:185}}\par
\textbf{reducing agent} -- das Reduktionsmittel\dotfill\hyperref[e:187]{\pageref*{e:187}}\par
\textbf{reduction} -- die Reduktion\dotfill\hyperref[e:184]{\pageref*{e:184}}\par
\textbf{relative atomic mass} -- die relative Atommasse\dotfill\hyperref[e:82]{\pageref*{e:82}}\par
\textbf{relative molecular mass} -- die relative Molekülmasse\dotfill\hyperref[e:254]{\pageref*{e:254}}\par
\textbf{residue} -- der Rückstand\dotfill\hyperref[e:33]{\pageref*{e:33}}\par
\textbf{reversible reaction} -- die umkehrbare Reaktion\dotfill\hyperref[e:144]{\pageref*{e:144}}\par
\textbf{round-bottomed flask} -- der Rundkolben\dotfill\hyperref[e:42]{\pageref*{e:42}}\par
\textbf{rust} -- der Rost\dotfill\hyperref[e:191]{\pageref*{e:191}}\par

\medskip{\large\textbf{\color{cornflower}S}}\par\nopagebreak
\textbf{safety goggles} -- die Schutzbrille\dotfill\hyperref[e:62]{\pageref*{e:62}}\par
\textbf{salt} -- das Salz\dotfill\hyperref[e:167]{\pageref*{e:167}}\par
\textbf{saturated} -- gesättigt\dotfill\hyperref[e:214]{\pageref*{e:214}}\par
\textbf{saturated solution} -- die gesättigte Lösung\dotfill\hyperref[e:14]{\pageref*{e:14}}\par
\textbf{separation technique} -- das Trennverfahren\dotfill\hyperref[e:30]{\pageref*{e:30}}\par
\textbf{single bond} -- die Einfachbindung\dotfill\hyperref[e:112]{\pageref*{e:112}}\par
\textbf{smell, odour} -- der Geruch\dotfill\hyperref[e:22]{\pageref*{e:22}}\par
\textbf{soap} -- die Seife\dotfill\hyperref[e:245]{\pageref*{e:245}}\par
\textbf{sodium hydroxide solution} -- die Natronlauge\dotfill\hyperref[e:181]{\pageref*{e:181}}\par
\textbf{solid, liquid, gaseous} -- fest, flüssig, gasförmig\dotfill\hyperref[e:29]{\pageref*{e:29}}\par
\textbf{solubility} -- die Löslichkeit\dotfill\hyperref[e:13]{\pageref*{e:13}}\par
\textbf{soluble} -- löslich\dotfill\hyperref[e:11]{\pageref*{e:11}}\par
\textbf{solute} -- der gelöste Stoff\dotfill\hyperref[e:10]{\pageref*{e:10}}\par
\textbf{solution} -- die Lösung\dotfill\hyperref[e:8]{\pageref*{e:8}}\par
\textbf{solvent} -- das Lösungsmittel\dotfill\hyperref[e:9]{\pageref*{e:9}}\par
\textbf{spatula} -- der Spatel\dotfill\hyperref[e:52]{\pageref*{e:52}}\par
\textbf{standard electrode potential} -- das Standardpotential\dotfill\hyperref[e:202]{\pageref*{e:202}}\par
\textbf{standard solution} -- die Maßlösung\dotfill\hyperref[e:176]{\pageref*{e:176}}\par
\textbf{starch} -- die Stärke\dotfill\hyperref[e:243]{\pageref*{e:243}}\par
\textbf{stir} -- rühren\dotfill\hyperref[e:66]{\pageref*{e:66}}\par
\textbf{stoichiometry} -- die Stöchiometrie\dotfill\hyperref[e:258]{\pageref*{e:258}}\par
\textbf{stopper, bung} -- der Stopfen\dotfill\hyperref[e:60]{\pageref*{e:60}}\par
\textbf{strong acid} -- die starke Säure\dotfill\hyperref[e:168]{\pageref*{e:168}}\par
\textbf{structural formula} -- die Strukturformel\dotfill\hyperref[e:103]{\pageref*{e:103}}\par
\textbf{subscript} -- der Index\dotfill\hyperref[e:132]{\pageref*{e:132}}\par
\textbf{substance} -- der Stoff\dotfill\hyperref[e:1]{\pageref*{e:1}}\par
\textbf{substitution reaction} -- die Substitutionsreaktion\dotfill\hyperref[e:229]{\pageref*{e:229}}\par
\textbf{sulfuric acid} -- die Schwefelsäure\dotfill\hyperref[e:179]{\pageref*{e:179}}\par
\textbf{surfactant, detergent} -- das Tensid\dotfill\hyperref[e:246]{\pageref*{e:246}}\par
\textbf{suspension} -- die Suspension\dotfill\hyperref[e:15]{\pageref*{e:15}}\par
\textbf{synthesis} -- die Synthese\dotfill\hyperref[e:133]{\pageref*{e:133}}\par

\medskip{\large\textbf{\color{cornflower}T}}\par\nopagebreak
\textbf{test (for)} -- der Nachweis\dotfill\hyperref[e:151]{\pageref*{e:151}}\par
\textbf{test tube} -- das Reagenzglas\dotfill\hyperref[e:39]{\pageref*{e:39}}\par
\textbf{thermometer} -- das Thermometer\dotfill\hyperref[e:57]{\pageref*{e:57}}\par
\textbf{titration} -- die Titration\dotfill\hyperref[e:173]{\pageref*{e:173}}\par
\textbf{titration curve} -- die Titrationskurve\dotfill\hyperref[e:174]{\pageref*{e:174}}\par
\textbf{tongs} -- die Tiegelzange\dotfill\hyperref[e:61]{\pageref*{e:61}}\par
\textbf{transition metals} -- die Übergangsmetalle\dotfill\hyperref[e:92]{\pageref*{e:92}}\par
\textbf{triple bond} -- die Dreifachbindung\dotfill\hyperref[e:114]{\pageref*{e:114}}\par
\textbf{tripod (and gauze)} -- der Dreifuß (mit Drahtnetz)\dotfill\hyperref[e:47]{\pageref*{e:47}}\par

\medskip{\large\textbf{\color{cornflower}U}}\par\nopagebreak
\textbf{universal indicator} -- der Universalindikator\dotfill\hyperref[e:161]{\pageref*{e:161}}\par
\textbf{unsaturated} -- ungesättigt\dotfill\hyperref[e:215]{\pageref*{e:215}}\par

\medskip{\large\textbf{\color{cornflower}V}}\par\nopagebreak
\textbf{valency} -- die Wertigkeit\dotfill\hyperref[e:121]{\pageref*{e:121}}\par
\textbf{van der Waals forces} -- die Van-der-Waals-Kräfte\dotfill\hyperref[e:120]{\pageref*{e:120}}\par

\medskip{\large\textbf{\color{cornflower}W}}\par\nopagebreak
\textbf{watch glass} -- das Uhrglas\dotfill\hyperref[e:55]{\pageref*{e:55}}\par
\textbf{weak acid} -- die schwache Säure\dotfill\hyperref[e:169]{\pageref*{e:169}}\par
\textbf{weigh out} -- abwiegen\dotfill\hyperref[e:69]{\pageref*{e:69}}\par
\textbf{word equation} -- die Wortgleichung\dotfill\hyperref[e:128]{\pageref*{e:128}}\par

\medskip{\large\textbf{\color{cornflower}Y}}\par\nopagebreak
\textbf{yield} -- die Ausbeute\dotfill\hyperref[e:150]{\pageref*{e:150}}\par
\end{multicols}

\end{document}
